% Equilibrium Forces, Moments and Couples — Pure Mathematics with Mechanics Upper Sixth — Cours signé OnBuch+
% Official MINESEC syllabus. Compiler avec : tectonic PMM14-equilibrium-forces-moments-and-couples.tex
\newcommand{\DOCMATIERE}{Pure Mathematics with Mechanics}
\newcommand{\DOCNIVEAU}{Upper Sixth}
\newcommand{\DOCMODULE}{Upper Sixth — Pure mathematics with mechanics}
\newcommand{\DOCLECON}{14}
\newcommand{\DOCTITRE}{Equilibrium Forces, Moments and Couples}
% OnBuch+ — préambule « cours signés » ENRICHI (compilé avec tectonic / XeLaTeX)
% Système visuel « Pop » d'OnBuch+ : fond crème (jamais blanc), bordures encre
% épaisses, ombres « dures » décalées, titres Archivo Black en capitales.
% Version MATHÉMATIQUES : graphes (pgfplots/tikz), boîtes pédagogiques
% (définition, propriété, méthode, attention, le savais-tu, exercice résolu,
% exercice, TP, bilan), page de garde et signature OnBuch+ ancrée en bas.
%
% Macros attendues dans main.tex AVANT \input{preamble} :
%   \DOCMATIERE  \DOCNIVEAU  \DOCMODULE  \DOCLECON (numéro)  \DOCTITRE
\documentclass[11pt]{article}

\usepackage{fontspec}
\usepackage[english]{babel}
\usepackage[a4paper,margin=2cm,top=3.4cm,bottom=2.8cm,headheight=1.4cm,footskip=1.3cm]{geometry}
\usepackage{amsmath,amssymb,amsfonts}
\usepackage{mathtools} % \xmapsto, \coloneqq... souvent utilisés par les modèles
\usepackage{cancel} % simplifications barrées (bilans de forces)
% \abs{z} (module, valeur absolue) et \norm{u} (norme), utilisés par les modèles
\providecommand{\abs}{}\let\abs\relax\DeclarePairedDelimiter{\abs}{\lvert}{\rvert}
\providecommand{\norm}{}\let\norm\relax\DeclarePairedDelimiter{\norm}{\lVert}{\rVert}
\usepackage[table]{xcolor} % [table] : fournit aussi \rowcolors (alternance de lignes)
\usepackage{pagecolor}
\usepackage{tikz}
\usetikzlibrary{arrows.meta,positioning,calc,shapes.geometric,shapes.misc,shapes.arrows,decorations.pathmorphing,decorations.pathreplacing,decorations.markings,patterns,shadows,mindmap,trees,fit,backgrounds,matrix,intersections,angles,quotes}
\usepackage{pgfplots}
\usepackage[version=4]{mhchem} % équations nucléaires \ce{^{235}_{92}U}
\usepackage[european,siunitx]{circuitikz} % schémas électriques
\pgfplotsset{compat=1.18}
\usepgfplotslibrary{fillbetween,groupplots} % aires entre courbes (calcul intégral)
\usepackage{enumitem}
\usepackage{fancyhdr}
% [most] charge toutes les bibliothèques tcolorbox (listings, minted, raster,
% documentation...) et gonfle inutilement la mémoire TeX sur un document long
% (~300 boîtes) : on ne charge que ce qui est réellement utilisé.
\usepackage[skins,breakable]{tcolorbox}
\usepackage{graphicx}
\usepackage{colortbl}
\usepackage{booktabs}
\usepackage{tabularx}
\usepackage{diagbox} % case d'en-tête diagonale des tableaux à double entrée
% Colonnes extensibles centrée / gauche / droite (C, L, R), souvent utilisées par les modèles
\newcolumntype{C}{>{\centering\arraybackslash}X}
\newcolumntype{L}{>{\raggedright\arraybackslash}X}
\newcolumntype{R}{>{\raggedleft\arraybackslash}X}
\usepackage{array}
\usepackage{multicol}
\usepackage{siunitx}
% parse-numbers=false : accepte aussi \SI{2,5 \times 10^{-3}}{...}
\sisetup{locale=UK,output-decimal-marker={.},group-minimum-digits=4,parse-numbers=false}
\DeclareSIUnit\ohm{\ensuremath{\symup{\Omega}}} % U+2126 absent de la police
\DeclareSIUnit\division{div} % réglages d'oscilloscope (ms/div, V/div)
\DeclareSIUnit\tr{tr} % tours (vitesse de rotation, ex. \SI{1500}{\tr\per\minute})
\usepackage{qrcode}
% \degree : commande fréquemment utilisée par les modèles pour un angle ou une
% température en dehors de \SI{}{} (non fournie par les paquets chargés).
\providecommand{\degree}{^{\circ}}
% \qed (fin de démonstration), utilisé par les modèles sans amsthm
\providecommand{\qed}{\hfill$\square$}
% \barre : double barre des valeurs interdites dans un tableau de variations (tkz-tab)
\providecommand{\barre}{\ensuremath{\|}}
% environnement proof (amsthm non chargé) utilisé par les modèles
\ifdefined\proof\else\newenvironment{proof}[1][Proof]{\par\noindent\textit{#1.}\ }{\qed\par}\fi
\usepackage{lastpage}
\usepackage{hyperref}
\hypersetup{hidelinks}

% ============================================================
% Palette « Pop » OnBuch+ — mêmes hex que la classe P (plus_ui.dart)
% ============================================================
\definecolor{poporange}{HTML}{F59321}
\definecolor{popdark}{HTML}{E07A0C}
\definecolor{popcream}{HTML}{FFF6E8}
\definecolor{popink}{HTML}{1C1714}
\definecolor{poppurple}{HTML}{7A5AE0}
\definecolor{popgreen}{HTML}{2E9E6A}
\definecolor{poppink}{HTML}{E0457B}
\definecolor{popblue}{HTML}{2F7DE1}
\definecolor{popgold}{HTML}{C98A12}
\definecolor{popmuted}{HTML}{8A7F76}
\definecolor{popline}{HTML}{EADFD2}
\definecolor{popgray}{HTML}{8A7F76}
\colorlet{popgrey}{popgray}
% Alias tolérants (le modèle nomme parfois les couleurs différemment)
\colorlet{popwhite}{white}
\colorlet{popblack}{popink}
\colorlet{popred}{poppink}
\colorlet{popyellow}{popgold}
% Teintes claires pour fonds de tableaux / zones
\colorlet{popblueL}{popblue!10!white}
\colorlet{poporangeL}{poporange!14!white}
\colorlet{popgreenL}{popgreen!12!white}
\colorlet{poppurpleL}{poppurple!12!white}
\colorlet{poppinkL}{poppink!12!white}
\colorlet{popgoldL}{popgold!14!white}

\pagecolor{popcream}

% ============================================================
% Typographie
% ============================================================
\defaultfontfeatures{Ligatures=TeX}
\setmainfont{PlusJakartaSans-Medium.ttf}[Path=../../../fonts/,BoldFont=PlusJakartaSans-Bold.ttf]
% Maths en sans-serif (Fira Math) avec chiffres et lettres droites de Plus
% Jakarta, pour que les formules se fondent dans le texte.
\usepackage[math-style=ISO,bold-style=ISO]{unicode-math}
\setmathfont{FiraMath-Regular.otf}[Path=../../../fonts/]
\setmathfont{PlusJakartaSans-Medium.ttf}[Path=../../../fonts/,range={up/num,up/{Latin,latin},\mathup}]
% (icomma retiré : décimales avec point en anglais)
% Caractères mathématiques Unicode absents des polices de texte (Plus Jakarta,
% Archivo Black) : rendus par leur équivalent mathématique, en texte comme en
% formule — sinon ils s'affichent en carré vide (ex. « Arithmétique dans ℤ »).
\usepackage{newunicodechar}
\newunicodechar{ℝ}{\ensuremath{\mathbb{R}}}
\newunicodechar{ℤ}{\ensuremath{\mathbb{Z}}}
\newunicodechar{ℂ}{\ensuremath{\mathbb{C}}}
\newunicodechar{ℕ}{\ensuremath{\mathbb{N}}}
\newunicodechar{ℚ}{\ensuremath{\mathbb{Q}}}
\newunicodechar{𝔻}{\ensuremath{\mathbb{D}}}
\newunicodechar{α}{\ensuremath{\alpha}}
\newunicodechar{β}{\ensuremath{\beta}}
\newunicodechar{γ}{\ensuremath{\gamma}}
\newunicodechar{δ}{\ensuremath{\delta}}
\newunicodechar{ε}{\ensuremath{\varepsilon}}
\newunicodechar{θ}{\ensuremath{\theta}}
\newunicodechar{λ}{\ensuremath{\lambda}}
\newunicodechar{μ}{\ensuremath{\mu}}
\newunicodechar{σ}{\ensuremath{\sigma}}
\newunicodechar{τ}{\ensuremath{\tau}}
\newunicodechar{φ}{\ensuremath{\varphi}}
\newunicodechar{ψ}{\ensuremath{\psi}}
\newunicodechar{ω}{\ensuremath{\omega}}
\newunicodechar{Σ}{\ensuremath{\Sigma}}
% majuscules produites par \MakeUppercase dans les titres (α-aminés -> Α)
\newunicodechar{Α}{\ensuremath{\alpha}}
\newunicodechar{Β}{\ensuremath{\beta}}
\newunicodechar{Γ}{\ensuremath{\Gamma}}
\newunicodechar{Δ}{\ensuremath{\Delta}}
\newunicodechar{Ε}{\ensuremath{\varepsilon}}
\newunicodechar{Θ}{\ensuremath{\Theta}}
\newunicodechar{Λ}{\ensuremath{\Lambda}}
\newunicodechar{Μ}{\ensuremath{\mu}}
\newunicodechar{Τ}{\ensuremath{\tau}}
\newunicodechar{Φ}{\ensuremath{\Phi}}
\newunicodechar{Ψ}{\ensuremath{\Psi}}
\newunicodechar{Ω}{\ensuremath{\Omega}}
\newunicodechar{↦}{\ensuremath{\mapsto}}
\newunicodechar{∈}{\ensuremath{\in}}
\newunicodechar{∉}{\ensuremath{\notin}}
\newunicodechar{∧}{\ensuremath{\wedge}}
\newunicodechar{⇔}{\ensuremath{\Leftrightarrow}}
\newunicodechar{⟺}{\ensuremath{\Longleftrightarrow}}
\newunicodechar{⇒}{\ensuremath{\Rightarrow}}
\newunicodechar{⟹}{\ensuremath{\Longrightarrow}}
\newunicodechar{∘}{\ensuremath{\circ}}
\newunicodechar{∪}{\ensuremath{\cup}}
\newunicodechar{∩}{\ensuremath{\cap}}
\newunicodechar{≡}{\ensuremath{\equiv}}
\newunicodechar{⊂}{\ensuremath{\subset}}
\newunicodechar{⊥}{\ensuremath{\perp}}
\newunicodechar{∃}{\ensuremath{\exists}}
\newunicodechar{∀}{\ensuremath{\forall}}
\newunicodechar{∛}{\ensuremath{\sqrt[3]{\,}}}
\newunicodechar{∜}{\ensuremath{\sqrt[4]{\,}}}
\newunicodechar{⃗}{\ensuremath{{}^{\to}}}
\newunicodechar{ⁿ}{\ifmmode^{n}\else\textsuperscript{\ensuremath{n}}\fi}
\newunicodechar{ˣ}{\ifmmode^{x}\else\textsuperscript{\ensuremath{x}}\fi}
\newunicodechar{ᵇ}{\ifmmode^{b}\else\textsuperscript{\ensuremath{b}}\fi}
\newunicodechar{ʳ}{\ifmmode^{r}\else\textsuperscript{\ensuremath{r}}\fi}
\newunicodechar{ᵘ}{\ifmmode^{u}\else\textsuperscript{\ensuremath{u}}\fi}
\newunicodechar{ʸ}{\ifmmode^{y}\else\textsuperscript{\ensuremath{y}}\fi}
\newunicodechar{ᵃ}{\ifmmode^{a}\else\textsuperscript{\ensuremath{a}}\fi}
\newunicodechar{ᵉ}{\ifmmode^{e}\else\textsuperscript{\ensuremath{e}}\fi}
\newunicodechar{ᵏ}{\ifmmode^{k}\else\textsuperscript{\ensuremath{k}}\fi}
\newunicodechar{ᵖ}{\ifmmode^{p}\else\textsuperscript{\ensuremath{p}}\fi}
\newunicodechar{ᴺ}{\ifmmode^{N}\else\textsuperscript{\ensuremath{N}}\fi}
\newunicodechar{ᶜ}{\ifmmode^{c}\else\textsuperscript{\ensuremath{c}}\fi}
\newunicodechar{ʰ}{\ifmmode^{h}\else\textsuperscript{\ensuremath{h}}\fi}
\newunicodechar{ᵠ}{\ifmmode^{q}\else\textsuperscript{\ensuremath{q}}\fi}
\newunicodechar{ₙ}{\ifmmode_{n}\else\textsubscript{\ensuremath{n}}\fi}
\newunicodechar{ₐ}{\ifmmode_{a}\else\textsubscript{\ensuremath{a}}\fi}
\newunicodechar{ᵢ}{\ifmmode_{i}\else\textsubscript{\ensuremath{i}}\fi}
\newunicodechar{ᵣ}{\ifmmode_{r}\else\textsubscript{\ensuremath{r}}\fi}
\newunicodechar{ₖ}{\ifmmode_{k}\else\textsubscript{\ensuremath{k}}\fi}
\newunicodechar{ᵦ}{\ifmmode_{\beta}\else\textsubscript{\ensuremath{\beta}}\fi}
\newunicodechar{ₓ}{\ifmmode_{x}\else\textsubscript{\ensuremath{x}}\fi}
\newfontfamily\popdisplay{ArchivoBlack-Regular.ttf}[Path=../../../fonts/]
\newfontfamily\poplogo{SpaceGrotesk-Bold.ttf}[Path=../../../fonts/]
\newfontfamily\popsemi{PlusJakartaSans-SemiBold.ttf}[Path=../../../fonts/]
\newfontfamily\popxbold{PlusJakartaSans-ExtraBold.ttf}[Path=../../../fonts/]
\newfontfamily\popxboldls{PlusJakartaSans-ExtraBold.ttf}[Path=../../../fonts/,LetterSpace=3.0]

\setlist[enumerate]{leftmargin=1.7em,itemsep=5pt,topsep=4pt}
\setlist[itemize]{leftmargin=1.7em,itemsep=4pt,topsep=4pt,label={\color{poporange}\textbullet}}
\setlength{\parindent}{0pt}
\setlength{\parskip}{5pt}
\renewcommand{\arraystretch}{1.25}

% pgfplots : style Pop par défaut
\pgfplotsset{
  every axis/.append style={axis line style={popink,line width=1pt},tick style={popink},
    grid style={popline,line width=0.5pt},label style={font=\small},tick label style={font=\footnotesize}},
  popcurve/.style={poporange,line width=1.8pt,no markers,smooth},
}
% Même style disponible dans un \draw TikZ simple (hors pgfplots)
\tikzset{popcurve/.style={poporange,line width=1.8pt}}
\tikzset{popgrid/.style={popline,very thin}}
\colorlet{popcurve}{poporange} % le nom du style est parfois utilisé comme couleur

% ============================================================
% Pill / squiggle
% ============================================================
\newcommand{\poppill}[3]{%
  \tikz[baseline=-0.55ex]\node[draw=popink,line width=1.2pt,rounded corners=2.6pt,%
    fill=#2,inner xsep=6.5pt,inner ysep=3pt]{%
    \popxboldls\fontsize{7}{7}\selectfont\color{#3}\MakeUppercase{#1}};%
}
\newcommand{\popsquiggle}[1]{%
  \tikz[baseline=-0.4ex]{
    \draw[#1,line width=2.6pt,line cap=round]
      plot [smooth] coordinates {(0,0.09) (0.34,0.01) (0.68,0.21) (1.02,0.01) (1.36,0.10)};
  }%
}
% Mot-clé surligné (marqueur orange sous le texte)
\newcommand{\cle}[1]{{\popxbold\color{popdark}#1}}

% ============================================================
% En-tête / pied
% ============================================================
\pagestyle{fancy}
\fancyhf{}
\renewcommand{\headrulewidth}{0pt}
\renewcommand{\footrulewidth}{0pt}
\lhead{\raisebox{-0.15ex}{\poplogo\fontsize{13}{13}\selectfont\color{popink}OnBuch\color{poporange}+}}
\rhead{\poppill{\DOCMATIERE\ \textbullet\ \DOCNIVEAU\ \textbullet\ Chapter \DOCLECON}{white}{popink}}
\lfoot{\popxbold\fontsize{7.6}{7.6}\selectfont\color{popmuted}\MakeUppercase{Signed course OnBuch+}\ \textbullet\ {\popsemi\DOCTITRE}}
\rfoot{\tikz[baseline=-0.6ex]\node[rounded corners=8.5pt,fill=popink,minimum height=17pt,minimum width=17pt,inner xsep=5pt,inner sep=0]{\popxbold\fontsize{8}{8}\selectfont\color{popcream}\thepage\,/\,\pageref*{LastPage}};}

% ============================================================
% Titres
% ============================================================
\newcommand{\chaptitre}[2]{%
  \begin{tcolorbox}[enhanced,colback=poporange,colframe=popink,boxrule=2.6pt,arc=11pt,
      drop shadow=popink,left=15pt,right=15pt,top=12pt,bottom=11pt,before skip=3pt,after skip=18pt]
    \poppill{#2}{popink}{popcream}\par\vspace{9pt}
    {\raggedright\hyphenpenalty=10000\exhyphenpenalty=10000\popdisplay\fontsize{22}{24}\selectfont\color{popcream}\MakeUppercase{#1}\par}
  \end{tcolorbox}
}
% Section numérotée : pastille orange avec le numéro + titre Archivo
\newcounter{popsec}
\newcommand{\coursec}[1]{%
  \stepcounter{popsec}\setcounter{popsub}{0}%
  \par\vspace{16pt}\noindent
  \begin{minipage}{\linewidth}
  \tikz[baseline=-0.7ex]\node[draw=popink,line width=1.6pt,fill=poporange,rounded corners=4pt,minimum size=22pt,inner sep=0]{\popdisplay\fontsize{12}{12}\selectfont\color{popcream}\thepopsec};%
  \hspace{8pt}{\popdisplay\fontsize{15}{17}\selectfont\color{popink}\MakeUppercase{#1}}\par
  \vspace{4pt}\noindent{\color{poporange}\rule{36pt}{3.6pt}}
  \end{minipage}\par\nopagebreak\vspace{8pt}
}
\newcounter{popsub}
\newcommand{\courssub}[1]{%
  \stepcounter{popsub}%
  \par\vspace{9pt}\noindent{\popxbold\fontsize{11.5}{13}\selectfont\color{popink}{\color{poporange}\thepopsec.\thepopsub}\hspace{6pt}#1}\par\nopagebreak\vspace{4pt}
}

% ============================================================
% Boîtes pédagogiques — même recette : fond blanc, bordure épaisse colorée,
% ombre dure encre, badge plein. Titre optionnel [#1].
% ============================================================
\tcbset{popbase/.style={breakable,enhanced,colback=white,boxrule=2.3pt,arc=9pt,
  shadow={3pt}{-3pt}{0pt}{popink},left=14pt,right=14pt,top=11pt,bottom=11pt,before skip=11pt,after skip=15pt,
  fontupper=\popsemi\fontsize{10.6}{14.5}\selectfont\color{popink},
  fontlower=\popsemi\fontsize{10.6}{14.5}\selectfont\color{popink}}}
% Coche dessinée (le glyphe ✓ n'existe pas dans les polices du document)
\newcommand{\popcheck}[1]{\tikz[baseline=-0.5ex]\draw[#1,line width=1.6pt,line cap=round,line join=round](-0.13,0.0)--(-0.04,-0.09)--(0.13,0.1);}
\providecommand{\coche}{\popcheck{popgreen}}
\providecommand{\resultat}[1]{\par\smallskip\noindent\fcolorbox{popgreen}{popgreenL}{\parbox{\dimexpr\linewidth-2\fboxsep-2\fboxrule\relax}{\textbf{Result.} #1}}\par\smallskip}
\newcommand{\popboxhead}[3]{\poppill{#1}{#2}{white}\ifx\relax#3\relax\else\hspace{6pt}{\popxbold\fontsize{10.6}{12}\selectfont\color{popink}#3}\fi\par\vspace{7pt}}

\newtcolorbox{aretenir}[1][]{popbase,colframe=popgreen,before upper={\popboxhead{Key points}{popgreen}{#1}}}
\newtcolorbox{exemplebox}[1][]{popbase,colframe=poporange,before upper={\popboxhead{Exemple}{poporange}{#1}}}
\newtcolorbox{definition}[1][]{popbase,colframe=popblue,before upper={\popboxhead{Definition}{popblue}{#1}}}
\newtcolorbox{propriete}[1][]{popbase,colframe=poppurple,before upper={\popboxhead{Property}{poppurple}{#1}}}
\newtcolorbox{methode}[1][]{popbase,colframe=popgold,before upper={\popboxhead{Method}{popgold}{#1}}}
\newtcolorbox{attention}[1][]{popbase,colframe=poppink,before upper={\popboxhead{Warning}{poppink}{#1}}}
\newtcolorbox{experience}[1][]{popbase,colframe=popblue,colback=popblueL,before upper={\popboxhead{Experiment}{popblue}{#1}}}
% « Le savais-tu ? » : carte violette pleine (contexte, culture, Cameroun)
\newtcolorbox{savaistu}[1][]{popbase,colback=poppurple,colframe=popink,
  fontupper=\popsemi\fontsize{10.4}{14.2}\selectfont\color{white},
  before upper={\poppill{Did you know?}{white}{poppurple}\ifx\relax#1\relax\else\hspace{6pt}{\popxbold\color{white}#1}\fi\par\vspace{7pt}}}
% Exercice résolu : énoncé en haut, \tcblower puis la solution sur fond crème
\newcounter{popexo}\newcounter{popexr}
\newtcolorbox{exoresolu}[1][]{popbase,colframe=poporange,colbacklower=popcream,
  segmentation style={popink,line width=1.2pt,dashed},
  before upper={\stepcounter{popexr}\popboxhead{Worked example \thepopexr}{poporange}{#1}},
  before lower={\poppill{Solution}{popink}{popcream}\par\vspace{6pt}}}
% Exercice (non résolu en place) — la correction est dans la partie Corrigés
\newtcolorbox{exercice}[1][]{popbase,colframe=popink,
  before upper={\stepcounter{popexo}\popboxhead{Exercise \thepopexo}{popink}{#1}}}
% Corrigé
\newtcolorbox{corrige}[1][]{popbase,colframe=popgreen,colback=popgreenL,
  before upper={\popboxhead{Solution}{popgreen}{#1}}}
% Objectifs / prérequis / bilan
\newtcolorbox{objectifs}{popbase,colframe=popink,colback=poporangeL,
  before upper={\popboxhead{Chapter objectives}{popink}{}}}
\newtcolorbox{prerequis}{popbase,colframe=popink,colback=popblueL,
  before upper={\popboxhead{Prerequisites}{popblue}{}}}
\newtcolorbox{bilan}{popbase,colframe=popink,colback=white,boxrule=2.8pt,
  before upper={\popboxhead{Fiche bilan}{poporange}{L'essentiel à connaître pour l'examen}}}
% Figure encadrée avec légende
\newtcolorbox{popfigure}[1][]{enhanced,breakable=false,colback=white,colframe=popline,boxrule=1.4pt,arc=8pt,
  left=10pt,right=10pt,top=10pt,bottom=8pt,before skip=10pt,after skip=12pt,halign=center,
  lower separated=false,halign lower=center,
  fontlower=\popsemi\fontsize{9}{11.5}\selectfont\color{popmuted},
  #1}
\newcommand{\legende}[1]{\par\vspace{6pt}{\popsemi\fontsize{9}{11.5}\selectfont\color{popmuted}#1}}

% ============================================================
% Signature OnBuch+
% ============================================================
\newcommand{\poplogoinline}[1]{{\poplogo\fontsize{#1}{#1}\selectfont\color{popink}OnBuch\color{poporange}+}}
\newcommand{\signatureonbuch}{%
  \begin{tcolorbox}[enhanced,colback=popink,colframe=popink,boxrule=2.2pt,arc=10pt,
      drop shadow=poporange,left=14pt,right=14pt,top=10pt,bottom=10pt,before skip=0pt,after skip=0pt]
    \tikz[baseline=-0.7ex]\node[circle,fill=popgreen,draw=popcream,line width=1.3pt,minimum size=22pt,inner sep=0]{\popcheck{white}};%
    \hspace{9pt}\begin{minipage}[c]{0.62\linewidth}
      {\popxbold\fontsize{12}{13}\selectfont\color{popcream}Signed\ \textbullet\ {\poplogo OnBuch\color{poporange}+}}\par\vspace{2pt}
      {\raggedright\popsemi\fontsize{8}{10.5}\selectfont\color{popline}\MakeUppercase{OnBuch+ teaching team}\\\MakeUppercase{Follows the official MINESEC syllabus}\par}
    \end{minipage}\hfill
    {\poplogo\fontsize{22}{22}\selectfont\color{popcream}OnBuch\color{poporange}+}
  \end{tcolorbox}%
}

% ============================================================
% Page de garde
% #1 titre · #2 matière · #3 niveau · #4 module · #5 n° leçon · #6 durée ·
% #7 description / objectifs (texte ou itemize)
% ============================================================
\newcommand{\pagegarde}[7]{%
  \thispagestyle{empty}
  \newgeometry{a4paper,margin=1.7cm,top=1.6cm,bottom=1.4cm}%
  \begin{tikzpicture}[remember picture,overlay]
    \fill[poporange!18] ([xshift=-3.2cm,yshift=-3.6cm]current page.north east) circle (4.6cm);
    \fill[poppurple!14] ([xshift=2.4cm,yshift=7.6cm]current page.south west) circle (3.8cm);
  \end{tikzpicture}%
  {\poplogo\fontsize{30}{30}\selectfont\color{popink}OnBuch\color{poporange}+}\hfill
  \raisebox{0.6ex}{\poppill{Signed course \textbullet\ Premium}{popink}{popcream}}\par
  \vspace{1pt}\popsquiggle{poppurple}\par
  \vspace{8pt}
  \begin{tcolorbox}[enhanced,colback=poporange,colframe=popink,boxrule=2.8pt,arc=13pt,
      drop shadow=popink,left=18pt,right=18pt,top=12pt,bottom=13pt,after skip=10pt]
    \poppill{#2 \textbullet\ #3}{popink}{popcream}\hspace{5pt}\poppill{#4}{white}{popink}\par\vspace{8pt}
    {\popdisplay\fontsize{14}{14}\selectfont\color{popink}CHAPTER #5}\par\vspace{4pt}
    {\raggedright\hyphenpenalty=10000\exhyphenpenalty=10000\popdisplay\fontsize{25}{27}\selectfont\color{popcream}\MakeUppercase{#1}\par}
    \vspace{8pt}
    {\popxbold\fontsize{9.5}{9.5}\selectfont\color{popink}\MakeUppercase{Suggested time: #6}}
  \end{tcolorbox}
  \begin{tcolorbox}[enhanced,colback=white,colframe=popink,boxrule=2.2pt,arc=10pt,
      drop shadow=popink,left=16pt,right=16pt,top=10pt,bottom=10pt,after skip=8pt]
    \poppill{In this chapter}{poppurple}{white}\par\vspace{5pt}
    \popsemi\fontsize{9.3}{12.6}\selectfont\color{popink}#7
  \end{tcolorbox}
  \noindent\begin{minipage}[t]{0.487\linewidth}
    \begin{tcolorbox}[enhanced,colback=popgreen,colframe=popink,boxrule=2.2pt,arc=10pt,
        drop shadow=popink,left=11pt,right=11pt,top=10pt,bottom=10pt,height=3.1cm,valign=center]
      \tcbox[colback=white,colframe=popink,boxrule=1.3pt,arc=5pt,boxsep=0pt,nobeforeafter,
        left=3pt,right=3pt,top=3pt,bottom=3pt,valign=center]{\qrcode[height=1.9cm]{https://play.google.com/store/apps/details?id=com.luvvix.onbuch&hl=en}}%
      \hspace{8pt}\begin{minipage}[c]{0.5\linewidth}
        {\popdisplay\fontsize{11}{12.5}\selectfont\color{white}\MakeUppercase{Download OnBuch}}\par\vspace{4pt}
        {\raggedright\popsemi\fontsize{8.4}{11}\selectfont\color{white}Free \textbullet\ Google Play\\AI tutor, past papers, results\par}
      \end{minipage}
    \end{tcolorbox}
  \end{minipage}\hfill
  \begin{minipage}[t]{0.487\linewidth}
    \begin{tcolorbox}[enhanced,colback=poppurple,colframe=popink,boxrule=2.2pt,arc=10pt,
        drop shadow=popink,left=11pt,right=11pt,top=10pt,bottom=10pt,height=3.1cm,valign=center]
      \tikz[baseline=-0.7ex]\node[circle,fill=white,draw=popink,line width=1.3pt,minimum size=1.6cm,inner sep=0]{\popdisplay\fontsize{24}{24}\selectfont\color{poppurple}+};%
      \hspace{8pt}\begin{minipage}[c]{0.6\linewidth}
        {\popdisplay\fontsize{11}{12.5}\selectfont\color{white}\MakeUppercase{Go OnBuch+}}\par\vspace{4pt}
        {\raggedright\popsemi\fontsize{8.4}{11}\selectfont\color{white}All signed courses, unlimited AI tutor, PDF sheets — OnBuch+ tab in the app\par}
      \end{minipage}
    \end{tcolorbox}
  \end{minipage}
  \vspace{8pt}
  \popcontenu
  \vfill
  \signatureonbuch
  \restoregeometry
  \newpage
}

% Rangée « Ce cours contient » (page de garde)
\newcommand{\poptile}[3]{% #1 couleur #2 gros texte #3 légende
  \begin{tcolorbox}[enhanced,colback=#1,colframe=popink,boxrule=2pt,arc=9pt,shadow={2.5pt}{-2.5pt}{0pt}{popink},
      left=6pt,right=6pt,top=9pt,bottom=9pt,halign=center,valign=center,height=1.75cm,width=\linewidth,nobeforeafter]
    {\popdisplay\fontsize{12.5}{13}\selectfont\color{white}\MakeUppercase{#2}}\par\vspace{3pt}
    {\centering\popsemi\fontsize{7.8}{9.5}\selectfont\color{white}#3\par}
  \end{tcolorbox}}
\newcommand{\popcontenu}{%
  \par\noindent{\popxboldls\fontsize{7.5}{7.5}\selectfont\color{popmuted}THIS SIGNED COURSE CONTAINS}\par\vspace{7pt}
  \noindent\begin{minipage}[t]{0.235\linewidth}\poptile{popblue}{Course}{complete, illustrated, step by step}\end{minipage}\hfill
  \begin{minipage}[t]{0.235\linewidth}\poptile{poporange}{Methods}{and worked examples}\end{minipage}\hfill
  \begin{minipage}[t]{0.235\linewidth}\poptile{poppink}{Exercises}{exam style + solutions}\end{minipage}\hfill
  \begin{minipage}[t]{0.235\linewidth}\poptile{popgreen}{Summary sheet}{the essentials to remember}\end{minipage}\par}

% Sommaire
\newtcolorbox{sommaire}{enhanced,colback=white,colframe=popink,boxrule=2.2pt,arc=10pt,
  drop shadow=popink,left=18pt,right=18pt,top=13pt,bottom=13pt,before skip=6pt,after skip=20pt,
  before upper={\poppill{Contents}{poppurple}{white}\par\vspace{9pt}\popsemi\fontsize{10.6}{14.5}\selectfont\color{popink}}}

% Signature de fin de cours — poussée tout en bas de la dernière page
\newcommand{\findecours}{%
  \par\vfill\nopagebreak
  \begin{center}\popsquiggle{poporange}\end{center}\vspace{4pt}
  \signatureonbuch
}
\AtBeginDocument{\def\llbracket{\mathopen{[\mkern-3.5mu[}}\def\rrbracket{\mathclose{]\mkern-3.5mu]}}} % intervalles d'entiers (glyphe ⟦ absent de la police)
% Glyphes absents de Fira Math : redéfinis à partir de glyphes disponibles
\AtBeginDocument{%
  \def\setminus{\mathbin{\backslash}}\let\smallsetminus\setminus
  \def\vdots{\mathord{\vbox{\baselineskip3.2pt\lineskiplimit0pt\kern4pt\hbox{.}\hbox{.}\hbox{.}}}}%
  \def\ddots{\mathinner{\mkern1mu\raise6.5pt\hbox{.}\mkern2mu\raise3.6pt\hbox{.}\mkern2mu\raise0.7pt\hbox{.}\mkern1mu}}%
  \def\triangle{\mathord{\scalebox{0.9}{$\symup{\Delta}$}}}%
  \def\nabla{\mathord{\raisebox{\depth}{\scalebox{1}[-1]{$\symup{\Delta}$}}}}%
  \def\checkmark{\popcheck{popdark}}%
  \def\star{\mathbin{\ast}}\def\bigstar{\mathord{\ast}}%
  \def\diamond{\mathbin{\scalebox{0.6}{$\Diamond$}}}%
  \def\bigcup{\mathop{\mathchoice{\raisebox{-0.3em}{\scalebox{1.7}{$\cup$}}}{\scalebox{1.25}{$\cup$}}{\cup}{\cup}}\displaylimits}%
  \def\bigcap{\mathop{\mathchoice{\raisebox{-0.3em}{\scalebox{1.7}{$\cap$}}}{\scalebox{1.25}{$\cap$}}{\cap}{\cap}}\displaylimits}%
  \def\lesseqqgtr{\mathrel{\lessgtr}}\def\bigoplus{\mathop{\oplus}\displaylimits}%
}
\providecommand{\ssi}{if and only if} % abréviation fréquente
\AtBeginDocument{\providecommand{\wideparen}{\overparen}} % arc de cercle

%% FIGFIX-BEGIN v4 — correctifs globaux des figures (docs/onbuch-generation/tools/figfix)
\usepackage{adjustbox}\usepackage{etoolbox}
% toute figure TikZ/pgfplots plus large que la ligne est réduite à la largeur disponible
\BeforeBeginEnvironment{tikzpicture}{\begin{adjustbox}{max width=\linewidth}}
\AfterEndEnvironment{tikzpicture}{\end{adjustbox}}
% légendes pgfplots lisibles par-dessus les courbes
\pgfplotsset{every axis/.append style={legend style={fill=white,fill opacity=0.92,text opacity=1,draw=black!15,inner sep=3pt}}}
% étiquettes d'arêtes (syntaxe edge["..."]) avec fond blanc
\tikzset{every edge quotes/.append style={fill=white,inner sep=1.2pt}}
%% FIGFIX-END
\begin{document}

\pagegarde{Equilibrium Forces, Moments and Couples}{Pure Mathematics with Mechanics}{Upper Sixth}{Upper Sixth — Pure mathematics with mechanics}{14}{15 periods}{This chapter studies how bodies remain at rest under the action of several forces: conditions of equilibrium, centres of gravity, resolution and composition of coplanar forces, and the elastic behaviour described by Hooke's law. It blends geometry, trigonometry and physical intuition into the core mechanics toolkit required at GCE Advanced Level.
\vspace{4pt}
\begin{multicols}{2}\small
\begin{itemize}
\item By the end of this chapter I can state and apply the conditions for equilibrium of a body under coplanar forces
\item By the end of this chapter I can determine the centre of gravity of a lamina experimentally and by calculation
\item By the end of this chapter I can represent forces as vectors and add them graphically and analytically
\item By the end of this chapter I can resolve a force into components along chosen axes, including on an inclined plane
\item By the end of this chapter I can draw a correct free-body diagram showing all forces acting on a body
\item By the end of this chapter I can distinguish and compute resultant and equilibrant forces
\end{itemize}\end{multicols}}

\begin{sommaire}
\begin{multicols}{2}
\begin{enumerate}
\item Equilibrium and Centre of Gravity
\item Force as a Vector; Addition and Resolution
\item Types of Forces and Free-Body Diagrams
\item Resultant and Equilibrant; Graphical Methods
\item Analytical Methods and Equilibrium Problems
\item Hooke's Law and Its Experimental Verification
\item Applications of Hooke's Law and Synthesis
\item Integration activity
\item Methods and worked examples
\item Exercises
\item Solutions to the exercises
\item Summary sheet
\end{enumerate}
\end{multicols}
\end{sommaire}

\begin{objectifs}
\begin{itemize}
\item By the end of this chapter I can state and apply the conditions for equilibrium of a body under coplanar forces
\item By the end of this chapter I can determine the centre of gravity of a lamina experimentally and by calculation
\item By the end of this chapter I can represent forces as vectors and add them graphically and analytically
\item By the end of this chapter I can resolve a force into components along chosen axes, including on an inclined plane
\item By the end of this chapter I can draw a correct free-body diagram showing all forces acting on a body
\item By the end of this chapter I can distinguish and compute resultant and equilibrant forces
\item By the end of this chapter I can use the triangle, parallelogram and polygon of forces to solve equilibrium problems
\item By the end of this chapter I can state, verify experimentally and apply Hooke's law to springs and elastic strings
\item By the end of this chapter I can solve GCE-style problems on suspended bodies, inclined planes and systems of springs
\end{itemize}
\end{objectifs}

\begin{prerequis}
\begin{itemize}
\item Trigonometric ratios (sine, cosine, tangent) and the sine and cosine rules from Lower Sixth pure mathematics
\item Vector notation, magnitude and direction of a vector, and simple vector addition
\item Basic kinematics and the concept of force, weight (W = mg) and Newton's laws from Lower Sixth mechanics
\item Solving simultaneous linear equations in two unknowns
\item Plotting and interpreting straight-line graphs (gradient, intercept)
\end{itemize}
\end{prerequis}

\newpage

\coursec{Equilibrium and Centre of Gravity}

\courssub{The Market Trader's Pole: A Question of Balance}

At Mokolo market in Yaoundé, a trader hangs a heavy sack of plantains on one end of a wooden balance pole and a small stone on the other — and the pole stays perfectly horizontal. Why does the small stone 'defeat' the heavy sack? On the Bamenda escarpment road, a loaded lorry parks on a steep slope without sliding, while a spring scale in a Douala hardware shop stretches exactly twice as much when the load is doubled. All these everyday situations obey the same mathematical laws: equilibrium of forces, moments, and Hooke's law. This chapter gives you the tools to explain and predict them — and to solve the GCE questions built on them.

\courssub{Equilibrium of Forces}

\courssub{Moments and the Moment of a Force}

\courssub{The Plumb-Line Experiment}

\courssub{Equilibrium of Suspended Bodies}

\courssub{Applications: Stability and Toppling}

\coursec{Equilibrium and Centre of Gravity}

\coursec{Equilibrium and Centre of Gravity}

\coursec{Equilibrium and Centre of Gravity}

\coursec{Equilibrium and Centre of Gravity}

\coursec{Force as a Vector; Addition and Resolution}

In Section 1 we met bodies in equilibrium: a hanging picture frame, a ladder, a metre rule balanced at its centre of gravity. In every case we spoke of \emph{forces} balancing each other, but we treated them rather loosely, as if only their sizes mattered. A moment's thought shows this is not enough. If two porters at the Mokolo market push a stalled cart, it matters enormously whether they push in the same direction or against each other. To handle equilibrium properly we must treat force as what it truly is: a \cle{vector}. That is the business of this section, and it will give us the tools --- addition, resolution, and the geometry of coplanar systems --- that we shall use for the rest of the chapter.

\courssub{Force as a Vector Quantity}

\begin{definition}[Force as a vector]
A \cle{force} is a vector quantity. To specify a force completely we must give four things:
\begin{itemize}
\item its \cle{magnitude} (or size), measured in newtons (\si{N});
\item its \cle{direction} (the angle it makes with a chosen reference line);
\item its \cle{line of action} (the infinite straight line along which it acts);
\item its \cle{point of application} (the point of the body where it acts).
\end{itemize}
A force is represented on paper by a \emph{directed line segment}: the length of the segment (to a chosen scale) gives the magnitude, and the arrowhead gives the sense.
\end{definition}

\begin{attention}[Point of application matters]
Two forces with the same magnitude and direction but different lines of action do \emph{not} have the same effect on a rigid body: pushing a door near the hinge and near the handle are very different experiences! For now, most of our bodies will be small enough to be treated as particles, so all forces act at one point; the line of action will become crucial when we study moments later in this chapter.
\end{attention}

\courssub{Addition of Forces: the Parallelogram and Triangle Laws}

Suppose two forces $\vec{P}$ and $\vec{Q}$ act at the same point $O$ of a body, with an angle $\theta$ between them. What single force would produce exactly the same effect?

\begin{definition}[Resultant]
The \cle{resultant} of a system of forces is the single force which, acting alone, has exactly the same effect on the body as the whole system acting together.
\end{definition}

\begin{definition}[Equilibrant]
The \cle{equilibrant} of a system of forces is the single force which, when added to the system, brings the body into equilibrium. It is equal in magnitude to the resultant but opposite in direction: $\vec{E} = -\vec{R}$.
\end{definition}

\begin{propriete}[Parallelogram law of forces]
If two forces acting at a point are represented in magnitude and direction by the adjacent sides of a parallelogram, then their resultant is represented in magnitude and direction by the diagonal of the parallelogram passing through their common point. (This law is an experimental fact of mechanics; its consequences must be derivable in the examination.)
\end{propriete}

Equivalently, the \cle{triangle law}: draw $\vec{P}$, then draw $\vec{Q}$ starting from the head of $\vec{P}$; the resultant $\vec{R}$ is the segment joining the tail of $\vec{P}$ to the head of $\vec{Q}$. The two laws are the same statement, since opposite sides of a parallelogram are equal and parallel. Note that in the triangle construction, the angle between $\vec{P}$ and $\vec{Q}$ (placed head-to-tail) is $180^\circ - \theta$.

\begin{popfigure}
\begin{center}
\begin{tikzpicture}[scale=1.1, >=stealth]
% Parallelogram law
\draw[->, very thick, popblue] (0,0) -- (3,0) node[midway, below] {$\vec{P}$};
\draw[->, very thick, poporange] (0,0) -- (1.2,1.6) node[midway, left] {$\vec{Q}$};
\draw[dashed, popmuted] (3,0) -- (4.2,1.6);
\draw[dashed, popmuted] (1.2,1.6) -- (4.2,1.6);
\draw[->, very thick, poppink] (0,0) -- (4.2,1.6) node[midway, above right] {$\vec{R}$};
\draw (0.7,0) arc (0:53.13:0.7) node[midway, right, xshift=2pt] {$\theta$};
\fill (0,0) circle (1.5pt) node[below left] {$O$};
% Triangle law
\begin{scope}[xshift=7cm]
\draw[->, very thick, popblue] (0,0) -- (3,0) node[midway, below] {$\vec{P}$};
\draw[->, very thick, poporange] (3,0) -- (4.2,1.6) node[midway, right] {$\vec{Q}$};
\draw[->, very thick, poppink] (0,0) -- (4.2,1.6) node[midway, above left] {$\vec{R}$};
\draw (3,0) ++(-0.7,0) arc (180:126.87:0.7) node[midway, left, xshift=-2pt] {$180^\circ-\theta$};
\fill (0,0) circle (1.5pt) node[below left] {$O$};
\end{scope}
\end{tikzpicture}
\end{center}
\legende{Parallelogram law (left) and triangle law (right) for the addition of two forces. In the triangle construction the interior angle is $180^\circ-\theta$.}
\end{popfigure}

\textbf{Magnitude and direction of the resultant.}\par
In the triangle, the angle between $\vec{P}$ and the reversed direction of $\vec{Q}$ is $180^\circ - \theta$. Applying the cosine rule:
\[ R^2 = P^2 + Q^2 - 2PQ\cos(180^\circ - \theta) = P^2 + Q^2 + 2PQ\cos\theta. \]
The angle $\alpha$ that $\vec{R}$ makes with $\vec{P}$ follows from the sine rule or from components:
\[ \tan\alpha = \frac{Q\sin\theta}{P + Q\cos\theta}. \]

\begin{aretenir}[Key formulae]
For two forces $P$ and $Q$ inclined at angle $\theta$:
\[ R = \sqrt{P^2 + Q^2 + 2PQ\cos\theta}, \qquad \tan\alpha = \frac{Q\sin\theta}{P+Q\cos\theta}. \]
Special cases: forces in the same direction ($\theta = 0$): $R = P+Q$; opposite directions ($\theta = 180^\circ$): $R = |P-Q|$; perpendicular forces ($\theta = 90^\circ$): $R = \sqrt{P^2+Q^2}$.
\end{aretenir}

\begin{exoresolu}[Resultant of two forces]
Two forces of \SI{5}{N} and \SI{8}{N} act at a point $O$ with an angle of $60^\circ$ between them. Find the magnitude of their resultant and the angle it makes with the \SI{8}{N} force.
\tcblower
\textbf{Solution.} Take $P = 8$, $Q = 5$, $\theta = 60^\circ$.
\[ R^2 = 8^2 + 5^2 + 2(8)(5)\cos 60^\circ = 64 + 25 + 80 \times 0.5 = 129, \]
so $R = \sqrt{129} \approx \SI{11.36}{N}$. The angle $\alpha$ with the \SI{8}{N} force satisfies
\[ \tan\alpha = \frac{5\sin 60^\circ}{8 + 5\cos 60^\circ} = \frac{5 \times 0.8660}{8 + 2.5} = \frac{4.330}{10.5} = 0.4124, \]
giving $\alpha \approx 22.4^\circ$.
\end{exoresolu}

\courssub{Resolution of a Force into Perpendicular Components}

Addition combines two forces into one. \cle{Resolution} does the reverse: it splits one force into two forces, called \cle{components}, whose resultant is the original force. There are infinitely many ways to do this, but the useful one is resolution along two \emph{perpendicular} axes.

If a force $F$ makes an angle $\theta$ with the $x$-axis, then completing the rectangle (a parallelogram with perpendicular sides) gives immediately:
\[ F_x = F\cos\theta, \qquad F_y = F\sin\theta. \]

\begin{attention}[Sine or cosine?]
The component \emph{adjacent} to the angle $\theta$ takes $\cos\theta$; the component \emph{opposite} to it takes $\sin\theta$. Mixing up sine and cosine is the single most common error at the GCE in this topic. Always draw the rectangle, mark the angle, and check: when $\theta = 0$ the whole force lies along the axis, so that component must be $F$ (and $\cos 0 = 1$ --- cosine confirmed).
\end{attention}

\begin{methode}[Resolving a system of forces]
\begin{enumerate}
\item Choose two perpendicular axes, placed so that as many forces as possible lie along them.
\item For each force, mark the angle it makes with an axis and write its two components.
\item Add the components along each axis to get $R_x = \sum F_x$ and $R_y = \sum F_y$.
\item Recombine if needed: $R = \sqrt{R_x^2 + R_y^2}$, $\tan\alpha = R_y/R_x$ (taking care with the quadrant of $\alpha$).
\end{enumerate}
\end{methode}

\courssub{Resolution on an Inclined Plane}

A lorry parked on the sloping road up to Buea does not slide down, yet gravity pulls it vertically. How does the slope "share" the weight? This is the classic situation where wise choice of axes pays off.

\begin{methode}[Resolving weight on an inclined plane]
For a body of mass $m$ on a plane inclined at $\theta$ to the horizontal, take axes \emph{along} the plane and \emph{perpendicular} to the plane. The weight $mg$ (vertically downward) then splits into:
\begin{itemize}
\item $mg\sin\theta$ acting \emph{down the slope};
\item $mg\cos\theta$ acting \emph{perpendicularly into the plane}.
\end{itemize}
\end{methode}

\textbf{Angle chasing: why the angle is $\theta$ (examinable proof).}\par
Let the plane make angle $\theta$ with the horizontal. The weight $\vec{W}$ is vertical, i.e.\ perpendicular to the horizontal. The normal to the plane is perpendicular to the plane. Now use the geometric lemma: \emph{if two angles have their sides mutually perpendicular, the angles are equal}. The sides of $\theta$ (horizontal, plane) are respectively perpendicular to the sides of the angle between the weight and the normal (vertical, normal). Hence the angle between the weight and the normal to the plane is exactly $\theta$.

More elementarily: the plane makes $\theta$ with the horizontal, so the normal makes $\theta$ with the vertical (rotating both lines by $90^\circ$ preserves the angle). Since the weight acts along the vertical, the angle between weight and normal is $\theta$. Then, in the right-angled resolution rectangle, the component adjacent to $\theta$ (along the normal, into the plane) is $mg\cos\theta$, and the component opposite (down the slope) is $mg\sin\theta$.

\begin{popfigure}
\begin{center}
\begin{tikzpicture}[scale=1.1, >=stealth]
% Inclined plane at 30 degrees
\pgfmathsetmacro{\planangle}{30}
\pgfmathsetmacro{\L}{6}
\pgfmathsetmacro{\H}{\L*tan(\planangle)} % 3.464
% Plane triangle
\draw[thick, popdark] (0,0) -- (\L,0) -- (\L,\H) -- cycle;
% Angle theta at origin
\draw (1.2,0) arc (0:\planangle:1.2) node[midway, right, xshift=4pt] {$\theta$};
% Block on plane: center at distance \d along slope, half-height along normal
\pgfmathsetmacro{\d}{3.5}
\pgfmathsetmacro{\halfw}{0.5}  % half width along slope
\pgfmathsetmacro{\halfh}{0.25} % half height perpendicular to slope
% Point on slope at distance \d from origin
\pgfmathsetmacro{\px}{\d*cos(\planangle)}
\pgfmathsetmacro{\py}{\d*sin(\planangle)}
% Outward normal angle = \planangle + 90
\pgfmathsetmacro{\nx}{cos(\planangle+90)}
\pgfmathsetmacro{\ny}{sin(\planangle+90)}
% Centre of block G
\pgfmathsetmacro{\Gx}{\px + \halfh*\nx}
\pgfmathsetmacro{\Gy}{\py + \halfh*\ny}
\coordinate (G) at (\Gx,\Gy);
% Draw block: scope shifted to G, rotated by \planangle
\begin{scope}[shift={(G)}, rotate=\planangle]
\draw[fill=popblueL, thick] (-\halfw,-\halfh) rectangle (\halfw,\halfh);
\end{scope}
% Weight (vertical down)
\draw[->, very thick, poppink] (G) -- ++(0,-1.8) node[below] {$mg$};
% Component down the slope (parallel to plane, downward)
\pgfmathsetmacro{\angledown}{\planangle+180}
\draw[->, very thick, poporange] (G) -- ++(\angledown:2.0) node[above left, text width=2.2cm, align=center] {$mg\sin\theta$\\down the slope};
% Component into the plane (perpendicular, inward normal)
\pgfmathsetmacro{\anglein}{\planangle-90}
\draw[->, very thick, popblue] (G) -- ++(\anglein:1.5) node[below right, text width=2cm, align=center] {$mg\cos\theta$\\into the plane};
% Dashed rectangle to show resolution
\draw[dashed, popmuted] (G) ++(\anglein:1.5) -- ++(\angledown:2.0);
\draw[dashed, popmuted] (G) ++(\angledown:2.0) -- ++(\anglein:1.5);
% Inward normal direction (for angle marking)
\draw[->, thick, popgreen] (G) -- ++(\anglein:1.8) node[below right] {normal (inward)};
% Angle theta between weight (vertical down, -90) and inward normal (\anglein)
% Draw arc from inward normal to vertical down (clockwise, delta = -\planangle)
\draw (G) ++(\anglein:0.9) arc (\anglein:-90:0.9) node[midway, left, xshift=-4pt] {$\theta$};
% Horizontal reference (vertical line through G) - dashed
\draw[dashed, popmuted] (G) -- ++(0,-1.8);
\end{tikzpicture}
\end{center}
\legende{Weight of a block on a plane inclined at $\theta=30^\circ$ resolved into $mg\sin\theta$ down the slope and $mg\cos\theta$ into the plane. The angle between the weight (vertical) and the inward normal equals the inclination $\theta$.}
\end{popfigure}

\begin{exoresolu}[Components of weight on a slope]
A block of mass \SI{10}{kg} rests on a plane inclined at $30^\circ$ to the horizontal. Find the components of its weight along and perpendicular to the plane. (Take $g = \SI{9.81}{m.s^{-2}}$.)
\tcblower
\textbf{Solution.} The weight is $W = mg = 10 \times 9.81 = \SI{98.1}{N}$.
\begin{align*}
\text{Down the slope: } & W\sin 30^\circ = 98.1 \times 0.5 = \SI{49.05}{N},\\
\text{Into the plane: } & W\cos 30^\circ = 98.1 \times 0.8660 = \SI{84.96}{N}.
\end{align*}
Check: $\sqrt{49.05^2 + 84.96^2} \approx 98.1 = W$. \checkmark
\end{exoresolu}

\courssub{Coplanar Forces}

\begin{definition}[Coplanar forces]
A system of forces is \cle{coplanar} when all their lines of action lie in one and the same plane. Within coplanar systems we distinguish:
\begin{itemize}
\item \cle{concurrent} systems: all lines of action pass through a single point (e.g.\ three ropes knotted together);
\item \cle{non-concurrent} systems: the lines of action do not all meet at one point (e.g.\ the forces on a ladder leaning against a wall).
\end{itemize}
\end{definition}

Why restrict ourselves to coplanar forces? Because two perpendicular directions suffice to describe every force in the plane: any coplanar force can be written $\vec{F} = F_x\,\vec{i} + F_y\,\vec{j}$, and the whole machinery of resolution above applies with just two axes. In three dimensions we would need a third component, vector products, and considerably heavier algebra --- that belongs to Further Mathematics and to university mechanics. At Advanced Level, every equilibrium problem you will meet (hanging bodies, inclined planes, ladders, frameworks) can be arranged, by symmetry or by the geometry of the situation, as a coplanar problem.

\begin{exemplebox}[Recognising coplanar systems]
\begin{itemize}
\item A picture hanging from a nail by two strings: the two tensions and the weight all lie in the vertical plane of the wall --- concurrent and coplanar.
\item A ladder against a wall: weight, two reactions and friction all act in the vertical plane perpendicular to the wall --- coplanar but non-concurrent.
\item A book on a sloping desk: weight, normal reaction and friction are coplanar in the vertical plane through the line of greatest slope.
\end{itemize}
\end{exemplebox}

\begin{exercice}[Practice]
\begin{enumerate}
\item Two forces of \SI{6}{N} and \SI{10}{N} act at a point with an angle of $120^\circ$ between them. Find the magnitude of their resultant.
\item A force of \SI{20}{N} acts at $40^\circ$ above the horizontal. Find its horizontal and vertical components.
\item A child of mass \SI{35}{kg} sits on a slide inclined at $25^\circ$ to the horizontal. Find the components of her weight along and perpendicular to the slide. ($g = \SI{9.81}{m.s^{-2}}$)
\item Two perpendicular forces have a resultant of \SI{13}{N}. If one force is \SI{5}{N}, find the other and the angle the resultant makes with it.
\end{enumerate}
\end{exercice}

\begin{corrige}[Exercise 1]
\begin{enumerate}
\item $R^2 = 6^2 + 10^2 + 2(6)(10)\cos 120^\circ = 36 + 100 - 60 = 76$, so $R = \sqrt{76} \approx \SI{8.72}{N}$.
\item $F_x = 20\cos 40^\circ \approx \SI{15.32}{N}$; $F_y = 20\sin 40^\circ \approx \SI{12.86}{N}$.
\item Along the slide: $35 \times 9.81 \times \sin 25^\circ \approx \SI{145.1}{N}$; perpendicular: $35 \times 9.81 \times \cos 25^\circ \approx \SI{311.2}{N}$.
\item Other force: $\sqrt{13^2 - 5^2} = \sqrt{144} = \SI{12}{N}$; angle with the \SI{5}{N} force: $\tan\alpha = 12/5$, so $\alpha \approx 67.4^\circ$.
\end{enumerate}
\end{corrige}

\begin{aretenir}[To remember]
\begin{itemize}
\item A force is a vector: magnitude, direction, line of action, point of application.
\item Resultant of two forces: $R^2 = P^2 + Q^2 + 2PQ\cos\theta$ (where $\theta$ is the angle between them).
\item Equilibrant: equal in magnitude to the resultant, opposite in direction ($\vec{E}=-\vec{R}$).
\item Resolution: adjacent component takes cosine, opposite takes sine.
\item On an inclined plane: $mg\sin\theta$ down the slope, $mg\cos\theta$ into the plane --- because the angle between the weight and the normal equals $\theta$.
\item All our systems are coplanar: two perpendicular axes are enough.
\end{itemize}
\end{aretenir}

With forces now added and resolved at will, we are ready in the next section to catalogue the \emph{types} of forces that act on bodies --- weight, reactions, tension, friction --- and to draw the free-body diagrams that turn a physical situation into a solvable mathematical problem.

\coursec{Types of Forces and Free-Body Diagrams}

In the previous section we learnt that a force is a \emph{vector}: it has a magnitude, a direction and a line of action, and we saw how to add forces and resolve them into components. But before we can add or resolve anything, we must answer a more basic question: \emph{which forces actually act on the body we are studying?} A student who misses one force, or invents one that does not exist, will get the wrong answer no matter how good his vector algebra is. This section is therefore about \textbf{recognising forces} and about the single most important tool of mechanics: the \cle{free-body diagram}. It covers Lessons 157 to 161 of the official syllabus: force as a vector, addition and resolution of forces, types of forces, free-body diagrams, resultant and equilibrant, and resolution on inclined planes.

\courssub{The Main Types of Forces}

\textbf{Weight.}\par
Every body near the Earth's surface is pulled towards the centre of the Earth. This pull is the \cle{weight} of the body.

\begin{definition}[Weight]
The \cle{weight} of a body of mass $m$ is the gravitational force exerted on it by the Earth:
\[ W = mg, \]
where $g \approx \SI{9.81}{m.s^{-2}}$ is the acceleration due to gravity. Weight is a force, measured in newtons (N), and it \textbf{always acts vertically downward}, through the centre of gravity of the body, whatever the position or motion of the body.
\end{definition}

\begin{attention}[Mass is not weight]
Mass $m$ is measured in kilograms; weight $W = mg$ is measured in newtons. A bag of rice of mass $25\ \mathrm{kg}$ has weight $W = 25 \times 9.81 \approx \SI{245}{N}$. On an inclined plane the weight is \emph{still} vertically downward — never "into the slope" and never "along the slope".
\end{attention}

\textbf{Normal reaction.}\par
Place a book on a table. The book does not fall through the table: the table pushes back on it. This push is the \cle{normal reaction} (or \emph{contact reaction}).

\begin{propriete}[Direction of the normal reaction]
The normal reaction $\vec{R}$ exerted by a surface on a body is always \textbf{perpendicular (normal) to the surface of contact}, directed \emph{away from the surface and into the body}. It is never along the surface. (Proof not examinable; this is a modelling property of contact forces.)
\end{propriete}

The word "normal" here is the mathematical one: it means "at right angles". On a horizontal floor, $R$ is vertical; on a wall, $R$ is horizontal; on a slope inclined at $\theta$ to the horizontal, $R$ is inclined at $\theta$ to the \emph{vertical}.

\textbf{Tension.}\par
When a string, rope, chain or cable is taut, it exerts a \cle{tension} force on whatever is attached to its ends.

\begin{attention}[A string can only pull]
The tension in a string always \textbf{pulls} on the body, \textbf{along the string}, directed \emph{away from the body towards the string}. A string can never push: if you try to push with a rope, it simply goes slack and the tension becomes zero. For a light inextensible string passing over a smooth pulley, the tension has the same magnitude $T$ throughout the string.
\end{attention}

\textbf{Friction.}\par
Rub your hands together: you feel a resistance. When two rough surfaces are in contact, the force of \cle{friction} acts \emph{along} the surface of contact (tangentially), in the direction that \textbf{opposes relative motion — or the tendency to move}.

\begin{definition}[Friction]
\cle{Friction} is the tangential contact force between two rough surfaces. Its direction is opposite to the direction in which the body is sliding, or would slide if friction were absent. A surface described as \cle{smooth} exerts no friction (only the normal reaction); a surface described as \cle{rough} exerts friction $F$, whose maximum value, called \emph{limiting friction}, will be studied in detail in a later chapter.
\end{definition}

To find the direction of friction, always ask: \emph{"If this surface were smooth, which way would the body slide?"} Friction points the opposite way. A crate pulled eastwards on rough ground experiences friction westwards; a block resting on a slope would slide down, so friction on it acts \emph{up} the slope.

\textbf{Thrust (compression) in rods.}\par
Unlike a string, a rigid rod or strut can \emph{push}. A rod in compression exerts a \cle{thrust} on the body, directed along the rod, \emph{pushing the body away from the rod}. A rod in tension pulls, exactly like a string. This is why scaffolding poles and the legs of a table can support loads that ropes cannot.

\textbf{Applied forces and elastic (spring) forces.}\par
An \cle{applied force} is any push or pull given directly in the problem: a man pushing a trolley with $\SI{50}{N}$, the wind on a signboard, the pull of a tractor. A \cle{spring} (or elastic cord) exerts a force that depends on how much it is stretched or compressed: it pulls when stretched and pushes when compressed, always towards its natural length. We shall see in Section 6 that this force obeys \emph{Hooke's law}, $T = kx$ — for now, simply treat it as a force acting along the spring.

\begin{savaistu}[Forces around us in Cameroon]
When a \emph{bendskin} (motorcycle taxi) climbs a hill in Bamenda, the friction between the tyres and the road pushes the motorcycle forward — friction is not always a nuisance! Without friction, wheels would spin uselessly and nobody could walk. Engineers choosing tyres for the rainy season on the Douala--Yaoundé road think constantly about friction. In the construction of the Kribi deep-sea port, cranes use steel cables (tension) and hydraulic rams (thrust) to lift containers; each force must be correctly identified for safe operation.
\end{savaistu}

\courssub{The Free-Body Diagram}

When several forces act at once, the mind gets confused unless we draw a clean picture. Physicists and engineers use one universal tool.

\begin{definition}[Free-body diagram]
A \cle{free-body diagram} (FBD) is a sketch showing a \textbf{single body, isolated from its surroundings}, together with \textbf{all the external forces acting on it} — and nothing else. Every contact or support is removed and \emph{replaced} by the force it exerts on the body. Forces that the body exerts on \emph{other} objects are never drawn.
\end{definition}

\begin{methode}[Checklist for drawing a free-body diagram]
\begin{enumerate}
\item \textbf{Isolate} the body: draw it alone (a simple dot or small box is enough).
\item \textbf{Weight}: draw $W = mg$ vertically downward from the centre of gravity.
\item \textbf{Reactions at contacts}: at each point of contact with a surface, draw the normal reaction $R$ perpendicular to the surface.
\item \textbf{Tensions and thrusts at attachments}: along each string (pulling) or rod (pulling or pushing).
\item \textbf{Applied forces}: any push or pull stated in the problem.
\item \textbf{Friction}: if the contact is rough, draw $F$ along the surface, opposing motion or the tendency to move.
\end{enumerate}
\end{methode}

\begin{attention}[The three classic errors]
\begin{itemize}
\item \textbf{Never add $ma$ as a force.} The quantity $m\vec{a}$ is the \emph{result} of the forces ($\sum \vec{F} = m\vec{a}$), not an extra force. Drawing it on the FBD double-counts it.
\item \textbf{Never draw the reaction along the surface.} $R$ is perpendicular to the contact surface; the \emph{tangential} component of contact is friction, a separate force.
\item \textbf{Never forget the weight.} Even a body on a table, in the air, or on a slope always has $W = mg$ acting vertically downward.
\end{itemize}
\end{attention}

\begin{popfigure}
\begin{tikzpicture}[scale=0.9, >=stealth]
% ---- 1. Hanging lamp ----
\begin{scope}[shift={(0,0)}]
\draw[thick] (-1,0) -- (1,0);
\draw[thick] (0,0) -- (0,-1.2);
\fill[poporange] (0,-1.45) circle (0.25);
\draw[->, very thick, popblue] (0,-1.45) -- (0,-2.6) node[below] {$W=mg$};
\draw[->, very thick, popgreen] (0,-1.45) -- (0,-0.35) node[right] {$T$};
\node at (0,-3.1) {\small lamp hanging};
\end{scope}
% ---- 2. Block on rough slope ----
\begin{scope}[shift={(4.2,-1.6)}]
\draw[thick] (-1.6,0) -- (1.6,0);
\draw[thick] (-1.6,0) -- (1.2,1.4);
\fill[popblue] (0,0.62) rectangle (0.55,1.05);
\draw[->, very thick, popdark] (0.27,0.83) -- (0.27,-0.5) node[below] {$W$};
\draw[->, very thick, popgreen] (0.27,0.83) -- (-0.35,1.55) node[above left=-2pt] {$R$};
\draw[->, very thick, poppink] (0.27,0.83) -- (-0.75,0.4) node[left] {$F$};
\node at (0,-1.5) {\small block on rough slope};
\end{scope}
% ---- 3. Ladder on wall ----
\begin{scope}[shift={(8.4,-1.6)}]
\draw[thick] (-1.2,0) -- (1.4,0);
\draw[thick] (1.4,0) -- (1.4,2.2);
\draw[very thick, poporange] (-0.9,0) -- (1.4,1.9);
\draw[->, very thick, popdark] (0.25,0.95) -- (0.25,-0.35) node[below] {$W$};
\draw[->, very thick, popgreen] (-0.9,0) -- (-0.9,0.85) node[above] {$R_1$};
\draw[->, very thick, popblue] (1.4,1.9) -- (0.55,1.9) node[left] {$R_2$};
\draw[->, very thick, poppink] (-0.9,0) -- (-0.05,0) node[below right=-2pt] {$F$};
\node at (0.3,-1.5) {\small ladder on wall};
\end{scope}
% ---- 4. Mass on spring ----
\begin{scope}[shift={(12.4,0)}]
\draw[thick] (-1,0) -- (1,0);
\draw[decorate, decoration={coil, segment length=4pt, amplitude=3pt}, thick] (0,0) -- (0,-1.1);
\fill[popblue] (-0.3,-1.6) rectangle (0.3,-1.1);
\draw[->, very thick, popdark] (0,-1.35) -- (0,-2.5) node[below] {$W$};
\draw[->, very thick, popgreen] (0,-1.35) -- (0,-0.35) node[right] {$T$};
\node at (0,-3.1) {\small mass on spring};
\end{scope}
\end{tikzpicture}
\legende{Four free-body diagrams: hanging lamp, block on a rough slope, ladder against a wall, mass on a spring.}
\end{popfigure}

\courssub{Worked Examples}

\begin{exemplebox}[Hanging lamp]
A lamp of mass $3\ \mathrm{kg}$ hangs at rest from a ceiling by a light string. Forces on the lamp: weight $W = 3 \times 9.81 = \SI{29.4}{N}$ downward, and tension $T$ upward along the string. Since the lamp is in equilibrium, $T = W = \SI{29.4}{N}$. Note: the pull of the lamp \emph{on the string} is not drawn — only forces \emph{on the lamp}.
\end{exemplebox}

\begin{exemplebox}[Ladder against a wall]
A ladder leans against a smooth vertical wall, its foot on rough horizontal ground. Forces on the ladder: weight $W$ vertically down at its midpoint; normal reaction $R_2$ from the smooth wall, horizontal (perpendicular to the wall); at the rough ground, normal reaction $R_1$ vertically up \emph{and} friction $F$ along the ground towards the wall (the foot would slip away from the wall, so friction opposes this).
\end{exemplebox}

\begin{exoresolu}[Block held on a slope by a string]
A block of mass $5\ \mathrm{kg}$ rests on a rough plane inclined at $30^{\circ}$ to the horizontal. It is held by a light string parallel to the slope, pulling up the slope. Draw the free-body diagram and identify every force.
\tcblower
\textbf{Solution.} Isolate the block. Apply the checklist:
\begin{enumerate}
\item \textbf{Weight}: $W = mg = 5 \times 9.81 = \SI{49.1}{N}$, vertically downward.
\item \textbf{Reaction}: $R$, perpendicular to the slope (at $30^{\circ}$ to the vertical), pointing away from the surface.
\item \textbf{Tension}: $T$, along the string, i.e.\ \emph{up the slope}, pulling the block.
\item \textbf{Applied force}: none besides the string.
\item \textbf{Friction}: the plane is rough; without friction the block would slide \emph{down}, so $F$ acts \emph{up the slope}, along the surface.
\end{enumerate}
Four forces in total: $W$, $R$, $T$, $F$. Resolving (preview of Section 5): perpendicular to the plane, $R = W\cos 30^{\circ}$; parallel to the plane, $T + F = W\sin 30^{\circ}$.
\end{exoresolu}

\begin{exoresolu}[Crate pulled on rough ground]
A porter at a market pulls a crate of mass $40\ \mathrm{kg}$ along rough horizontal ground with a rope inclined at $25^{\circ}$ above the horizontal, with tension $P = \SI{120}{N}$. List the forces on the crate.
\tcblower
\textbf{Solution.} Checklist on the crate:
\begin{enumerate}
\item Weight $W = 40 \times 9.81 = \SI{392.4}{N}$, vertically downward.
\item Normal reaction $R$, vertically upward (perpendicular to the horizontal ground).
\item Tension $P = \SI{120}{N}$ along the rope, at $25^{\circ}$ above the horizontal, pulling the crate.
\item Friction $F$ along the ground, horizontal, \emph{opposing} the motion (backwards).
\end{enumerate}
Note that $R \neq W$ here: the rope has an upward component $P\sin 25^{\circ}$ which helps to support the crate, so $R = W - P\sin 25^{\circ} \approx 392.4 - 50.7 = \SI{341.7}{N}$.
\end{exoresolu}

\begin{exoresolu}[Particle on a smooth inclined plane with horizontal force]
A particle of mass $2\ \mathrm{kg}$ rests in equilibrium on a smooth plane inclined at $40^{\circ}$ to the horizontal, held by a horizontal force $H$ acting towards the plane. Draw the free-body diagram and write the equilibrium equations.
\tcblower
\textbf{Solution.} Forces on the particle:
\begin{enumerate}
\item Weight $W = 2 \times 9.81 = \SI{19.62}{N}$, vertically downward.
\item Normal reaction $R$, perpendicular to the plane (at $40^{\circ}$ to the vertical).
\item Applied horizontal force $H$, acting horizontally towards the plane.
\end{enumerate}
No friction (smooth plane). Resolving perpendicular to the plane: $R = W\cos 40^{\circ} + H\sin 40^{\circ}$.
Resolving parallel to the plane (up the plane positive): $H\cos 40^{\circ} - W\sin 40^{\circ} = 0 \Rightarrow H = W\tan 40^{\circ} \approx \SI{16.5}{N}$.
Then $R = W\cos 40^{\circ} + W\tan 40^{\circ}\sin 40^{\circ} = W/\cos 40^{\circ} \approx \SI{25.6}{N}$.
\end{exoresolu}

\courssub{Wrong versus Right: a Typical Diagram Error}

The figure below shows, for a block on an inclined plane, the most frequent wrong diagram seen in examinations, next to the correct one.

\begin{popfigure}
\begin{tikzpicture}[scale=0.95, >=stealth]
% WRONG
\begin{scope}
\draw[thick] (-1.6,0) -- (1.6,0);
\draw[thick] (-1.6,0) -- (1.2,1.4);
\fill[gray] (0,0.62) rectangle (0.55,1.05);
\draw[->, very thick, poppink] (0.27,0.83) -- (1.15,1.45) node[right] {$R$?};
\draw[->, very thick, poppink] (0.27,0.83) -- (1.05,0.25) node[right] {$W$?};
\node at (0,-0.7) {\small \textbf{WRONG}};
\node[align=center] at (0,-1.3) {\small $R$ along slope,\\ \small $W$ into slope};
\end{scope}
% RIGHT
\begin{scope}[shift={(6.5,0)}]
\draw[thick] (-1.6,0) -- (1.6,0);
\draw[thick] (-1.6,0) -- (1.2,1.4);
\fill[popblue] (0,0.62) rectangle (0.55,1.05);
\draw[->, very thick, popgreen] (0.27,0.83) -- (-0.35,1.55) node[above left=-2pt] {$R$};
\draw[->, very thick, popdark] (0.27,0.83) -- (0.27,-0.5) node[below] {$W$};
\draw[->, very thick, poppink] (0.27,0.83) -- (-0.75,0.4) node[left] {$F$};
\node at (0,-0.7) {\small \textbf{RIGHT}};
\node[align=center] at (0,-1.3) {\small $R$ perpendicular to slope,\\ \small $W$ vertical, $F$ along slope};
\end{scope}
\end{tikzpicture}
\legende{Wrong versus right free-body diagram of a block on a rough inclined plane.}
\end{popfigure}

\begin{attention}[Examination traps]
\begin{itemize}
\item The weight is \textbf{vertical}, not perpendicular to the slope. Only the \emph{components} of $W$ ($W\sin\theta$ along, $W\cos\theta$ perpendicular) may be drawn — and then $W$ itself must be removed to avoid double-counting.
\item Do not draw the force the block exerts on the plane: the FBD shows forces \emph{on the block} only.
\item On a \emph{smooth} surface, there is no friction at all — do not invent it.
\item For a system of coplanar forces (Lesson 162), all forces lie in one plane; the FBD must reflect this planarity.
\item When resolving forces on an inclined plane (Lesson 161), choose axes parallel and perpendicular to the plane for simplicity.
\end{itemize}
\end{attention}

\begin{aretenir}[Key points]
\begin{itemize}
\item Weight $W = mg$ always acts vertically downward through the centre of gravity.
\item The normal reaction is always perpendicular to the contact surface; friction acts along it, opposing motion or its tendency.
\item Tension pulls along the string; a string never pushes. A rod can push (thrust) or pull.
\item A free-body diagram shows one isolated body with all external forces on it — weight, reactions, tensions/thrusts, applied forces, friction — and nothing else.
\item Never draw $ma$ as a force; never forget the weight.
\item For equilibrium of coplanar forces, the vector sum is zero: $\sum \vec{F} = \vec{0}$, which gives two scalar equations $\sum F_x = 0$, $\sum F_y = 0$ (or resolved along/perpendicular to a plane).
\end{itemize}
\end{aretenir}

\begin{exercice}[Forces on everyday objects]
For each situation, list all the forces acting on the body in italics and sketch its free-body diagram:
\begin{enumerate}
\item a \emph{book} resting on a horizontal table;
\item a \emph{picture frame} hanging from a nail by a string over its top edge;
\item a \emph{box} being pushed up a rough slope by a horizontal applied force;
\item a \emph{ball} hanging from a spring fixed to the ceiling, pulled down and held at rest.
\end{enumerate}
\end{exercice}

\begin{corrige}[Exercise 1]
\begin{enumerate}
\item Weight $W$ downward; normal reaction $R$ upward. Two forces.
\item Weight $W$ downward; tension $T$ upward along the string. Two forces.
\item Weight $W$ vertically down; normal reaction $R$ perpendicular to the slope; applied push $P$ horizontal; friction $F$ down the slope (opposing the upward motion). Four forces.
\item Weight $W$ downward; spring tension $T$ upward along the spring; applied downward pull $P$. Three forces.
\end{enumerate}
In each case only forces acting \emph{on} the italicised body are drawn.
\end{corrige}

With a correct free-body diagram in hand, we are ready for the next step: combining these forces to find their \emph{resultant}, or showing that they balance — the subject of the following sections.

\coursec{Resultant and Equilibrant; Graphical Methods}

In the previous section we learnt how to identify every force acting on a body and how to display them on a \cle{free-body diagram}. That picture immediately raises a practical question: when several forces act together, what is their \emph{combined effect}? A loaded truck climbing the steep hill out of Bamenda is pulled forward by the engine's driving force, held back by friction and by the component of its weight parallel to the slope; whether it accelerates or moves at constant velocity depends on the \emph{single force equivalent} to all of these. This section answers that question for coplanar forces acting at a point (concurrent forces), first by defining the \cle{resultant} and the \cle{equilibrant}, then by developing \emph{graphical} methods — the parallelogram and the polygon of forces — which are quick, visual and fully accepted at the GCE Advanced Level when drawn accurately to scale.

\courssub{The Resultant of a System of Forces}

\begin{definition}[Resultant]
The \cle{resultant} of a system of forces is the \emph{single force} which, acting alone on the body, produces exactly the same effect (the same tendency to translate and, where relevant, to rotate) as the whole system acting together. For a system of concurrent forces $\vec{F}_1, \vec{F}_2, \ldots, \vec{F}_n$ acting at a point, the resultant $\vec{R}$ is their vector sum:
\[
\vec{R} = \vec{F}_1 + \vec{F}_2 + \cdots + \vec{F}_n .
\]
\end{definition}

The phrase ``same effect'' is essential. Two workers pushing a stalled vehicle with forces of \SI{300}{N} and \SI{400}{N} in the same direction achieve exactly what one person pushing with \SI{700}{N} would achieve. The resultant is not an extra force added to the diagram; it is a \emph{replacement} for the whole set.

\begin{attention}[The resultant is not an additional force]
A very common error in free-body diagrams is to draw the real forces \emph{and then also} the resultant $\vec{R}$. This double-counts the effect. Either you draw the individual forces, or you draw their resultant — never both on the same diagram.
\end{attention}

\courssub{The Equilibrant}

Suppose the resultant of the forces on a body is $\vec{R}$. If we want the body to remain at rest (or in uniform motion), we must apply one more force that exactly cancels $\vec{R}$.

\begin{definition}[Equilibrant]
The \cle{equilibrant} $\vec{E}$ of a system of forces is the single force which, when added to the system, brings the body into equilibrium. It is:
\begin{itemize}
\item \textbf{equal in magnitude} to the resultant: $E = R$;
\item \textbf{opposite in direction} to the resultant;
\item \textbf{collinear} with the resultant (same line of action).
\end{itemize}
Hence $\vec{E} = -\vec{R}$, and $\vec{F}_1 + \vec{F}_2 + \cdots + \vec{F}_n + \vec{E} = \vec{0}$.
\end{definition}

\begin{popfigure}
\begin{center}
\begin{tikzpicture}[scale=1.0, >=stealth]
% line of action
\draw[dashed, popmuted] (-0.5,0) -- (6.5,0);
\fill[popink] (3,0) circle (2pt) node[below=2pt, text=popink] {$P$};
% resultant
\draw[->, very thick, popblue] (3,0) -- (6,0) node[midway, above, text=popblue] {$\vec{R}$};
% equilibrant
\draw[->, very thick, poporange] (3,0) -- (0,0) node[midway, above, text=poporange] {$\vec{E}$};
\node[text=popdark, align=center] at (3,-1.0) {$E = R$, opposite sense, same line of action through $P$};
\end{tikzpicture}
\legende{The resultant $\vec{R}$ and the equilibrant $\vec{E}$ of the same system: equal magnitudes, opposite directions, same line of action through $P$.}
\end{center}
\end{popfigure}

\begin{exemplebox}[Resultant versus equilibrant]
Three tug-of-war teams pull a ring with a combined effect equivalent to a single force of \SI{850}{N} directed due east. Then $R = \SI{850}{N}$ east. To keep the ring stationary, a fourth team must apply the equilibrant: $E = \SI{850}{N}$ due \emph{west}, along the same line. Note that the equilibrant is not ``the force that was missing'' in any mysterious sense — it is simply $-\vec{R}$.
\end{exemplebox}

\courssub{Graphical Method for Two Forces: the Parallelogram Law}

\begin{propriete}[Parallelogram of forces]
If two forces $\vec{F}_1$ and $\vec{F}_2$ acting at a point $O$ are represented in magnitude and direction by two adjacent sides of a parallelogram drawn from $O$, then their resultant $\vec{R}$ is represented in magnitude and direction by the diagonal of the parallelogram drawn from $O$. (This is the experimental law of vector addition; its analytical consequences are proved in the next section.)
\end{propriete}

\begin{methode}[Finding the resultant of two forces graphically]
\begin{enumerate}
\item \textbf{Choose a scale}, e.g.\ $1\ \mathrm{cm} \equiv 10\ \mathrm{N}$, so that the drawing fits the page but stays large enough for accuracy.
\item From $O$, draw $\vec{F}_1$ to scale in its true direction (measure angles with a protractor).
\item From the \emph{same point} $O$, draw $\vec{F}_2$ to scale in its true direction.
\item Complete the parallelogram with dashed lines parallel to the two forces.
\item Draw the diagonal from $O$: this is $\vec{R}$. Measure its length and convert back with the scale; measure its angle with the protractor.
\end{enumerate}
\end{methode}

\begin{exoresolu}[Parallelogram construction]
Two forces of \SI{40}{N} and \SI{30}{N} act at a point $O$, the angle between them being $60^{\circ}$. Find their resultant graphically.
\tcblower
\textbf{Scale:} $1\ \mathrm{cm} \equiv 10\ \mathrm{N}$, so the forces are drawn as segments of $4.0\ \mathrm{cm}$ and $3.0\ \mathrm{cm}$.

\textbf{Construction:} draw the $4.0\ \mathrm{cm}$ segment horizontally from $O$; with a protractor draw the $3.0\ \mathrm{cm}$ segment at $60^{\circ}$ to it; complete the parallelogram; draw the diagonal from $O$.

\textbf{Measurement:} the diagonal measures about $6.1\ \mathrm{cm}$, giving $R \approx \SI{61}{N}$, at an angle of about $25^{\circ}$ to the \SI{40}{N} force.

\textbf{Check (analytical, for confidence):}
\[
R = \sqrt{40^2 + 30^2 + 2(40)(30)\cos 60^{\circ}} = \sqrt{1600 + 900 + 1200} = \sqrt{3700} \approx \SI{60.8}{N}.
\]
The angle $\theta$ with the \SI{40}{N} force satisfies $\tan\theta = \frac{30\sin 60^{\circ}}{40 + 30\cos 60^{\circ}} = \frac{25.98}{55} \approx 0.472$, so $\theta \approx 25.3^{\circ}$. The graphical values \SI{61}{N} and $25^{\circ}$ agree within $0.5\ \%$ — excellent for a scale drawing.
\end{exoresolu}

\courssub{The Polygon of Forces for Several Coplanar Forces}

With three or more forces, completing parallelograms pairwise becomes clumsy. The \cle{polygon of forces} generalises the idea: add the vectors \emph{head to tail}.

\begin{methode}[Polygon of forces]
\begin{enumerate}
\item Choose a suitable scale and record it on the diagram.
\item From a starting point $A$, draw the first force $\vec{F}_1$ to scale, in its true direction, ending at $B$.
\item From $B$ (the \emph{head} of $\vec{F}_1$), draw $\vec{F}_2$ to scale in its true direction, ending at $C$. Continue: each force starts where the previous one ends. \textbf{Any order gives the same resultant.}
\item When the last force ends at $E$, the vector from the \emph{start} $A$ to the \emph{finish} $E$ is the resultant $\vec{R}$. Measure its length and direction.
\item If the polygon \emph{closes} (the last head returns to $A$), then $\vec{R} = \vec{0}$: the system is in equilibrium.
\end{enumerate}
\end{methode}

\begin{popfigure}
\begin{center}
\begin{tikzpicture}[scale=0.55, >=stealth]
% F1: 5 N at 0°, F2: 4 N at 70°, F3: 4.5 N at 150°, F4: 3 N at 250°
\coordinate (A) at (0,0);
\coordinate (B) at (5,0);
\coordinate (C) at (6.368,3.759);
\coordinate (D) at (2.471,6.009);
\coordinate (E) at (1.445,3.190);
\draw[->, very thick, popblue] (A) -- (B) node[midway, below, text=popblue] {$\vec{F}_1$};
\draw[->, very thick, popgreen] (B) -- (C) node[midway, right, text=popgreen] {$\vec{F}_2$};
\draw[->, very thick, poppurple] (C) -- (D) node[midway, above, text=poppurple] {$\vec{F}_3$};
\draw[->, very thick, popgold] (D) -- (E) node[midway, left, text=popgold] {$\vec{F}_4$};
\draw[->, ultra thick, poporange] (A) -- (E) node[midway, left=4pt, text=poporange] {$\vec{R}$};
\fill[popink] (A) circle (2.5pt) node[below left, text=popink] {$A$ (start)};
\fill[popink] (E) circle (2.5pt) node[above right=1pt, text=popink] {$E$ (finish)};
\end{tikzpicture}
\legende{Polygon of forces for four coplanar forces drawn head-to-tail. The closing side $\vec{R}$ (orange), from the start $A$ to the finish $E$, is the resultant.}
\end{center}
\end{popfigure}

\begin{attention}[Direction of the arrows]
In the force polygon, \textbf{all the given forces follow round in the same sense} (head to tail, head to tail, \ldots). The \textbf{resultant goes the opposite way}: from the \emph{start} of the first force to the \emph{end} of the last. Never draw the resultant head-to-tail like the others — that would represent $\vec{R}$ added to the system, not equivalent to it. Likewise, the equilibrant, if drawn on the polygon, is the side that \emph{closes} the polygon following the same sense as the other forces.
\end{attention}

\begin{propriete}[Graphical condition for equilibrium]
A system of coplanar forces acting at a point is in equilibrium \emph{if and only if} its force polygon \textbf{closes} (the head of the last force coincides with the tail of the first). Equivalently, $\vec{R} = \vec{0}$. (Proof: the polygon is exactly the geometrical addition of the vectors; closure means the vector sum is zero, which is the condition for equilibrium of a particle.)
\end{propriete}

This gives a powerful graphical test: to check whether given forces can hold a body in equilibrium, simply draw the polygon. If it closes, equilibrium is possible; the size of the ``gap'' otherwise gives the magnitude of the resultant, and the force needed to close the gap is the equilibrant.

\courssub{Space Diagram, Force Diagram and Bow's Notation}

When forces act at \emph{different points} of an extended body (a ladder, a beam), we need two drawings:

\begin{itemize}
\item the \cle{space diagram}: a sketch of the body showing the \emph{lines of action} of the forces and the angles between them;
\item the \cle{force diagram} (force polygon): the vectors drawn head-to-tail to scale, giving magnitudes and the resultant.
\end{itemize}

An exam-useful extra is \cle{Bow's notation}: in the space diagram, capital letters ($A$, $B$, $C$, \ldots) are placed in the spaces \emph{between} the lines of action, and each force is named by the two letters on either side of it (force $AB$, force $BC$, \ldots). In the force diagram the corresponding vertices are labelled with the matching small letters $a$, $b$, $c$, \ldots, so the force $AB$ is the vector from $a$ to $b$. Bow's notation removes all ambiguity about which force is which and is appreciated by examiners in structured graphical questions.

\begin{exemplebox}[Bow's notation illustrated]
A beam is supported at $A$ and $B$ and carries a load $W$ at $C$. In the space diagram, label the spaces around the beam: $A$ (left of left support), $B$ (between supports), $C$ (right of right support), $D$ (above the beam). The forces are then: reaction at left support = $DA$, reaction at right support = $AB$, weight = $BC$, and any other force = $CD$. In the force diagram, draw vectors $da$, $ab$, $bc$, $cd$ head-to-tail. The polygon closes if the beam is in equilibrium.
\end{exemplebox}

\courssub{Accuracy of Graphical Work}

Graphical methods are only as good as the drawing. Follow these rules:

\begin{itemize}
\item \textbf{Scale choice:} make the diagram as large as the paper allows while keeping every force on the page. A polygon spread over $20\ \mathrm{cm}$ is far more accurate than one squeezed into $5\ \mathrm{cm}$.
\item \textbf{Instruments:} sharp pencil, ruler with clear millimetre marks, good protractor. Measure lengths to the nearest millimetre and angles to the nearest degree.
\item \textbf{Error estimate:} with a scale of $1\ \mathrm{cm} \equiv 10\ \mathrm{N}$, a $\pm 1\ \mathrm{mm}$ reading error gives $\pm \SI{1}{N}$ uncertainty per measurement; for a resultant of \SI{60}{N} this is roughly $\pm 2\ \%$. Always quote graphical answers sensibly, e.g.\ $R \approx \SI{61}{N}$ rather than $\SI{61.37}{N}$.
\item \textbf{State the scale} on your diagram — an unscaled drawing earns no marks.
\end{itemize}

\begin{exoresolu}[Polygon of forces and the equilibrium test]
Four forces act at a point $O$: $F_1 = \SI{50}{N}$ due east; $F_2 = \SI{40}{N}$ at $70^{\circ}$ (measured anticlockwise from east); $F_3 = \SI{45}{N}$ at $150^{\circ}$; $F_4 = \SI{30}{N}$ at $250^{\circ}$. Determine graphically whether the system is in equilibrium; if not, find the resultant and the equilibrant.
\tcblower
\textbf{Scale:} $1\ \mathrm{cm} \equiv 10\ \mathrm{N}$; lengths $5.0$, $4.0$, $4.5$, $3.0\ \mathrm{cm}$.

\textbf{Construction:} from $A$ draw $\vec{F}_1$ east to $B$; from $B$ draw $\vec{F}_2$ at $70^{\circ}$ to $C$; from $C$ draw $\vec{F}_3$ at $150^{\circ}$ to $D$; from $D$ draw $\vec{F}_4$ at $250^{\circ}$. The last head ends at $E$, which does \emph{not} coincide with $A$: the system is \textbf{not} in equilibrium.

\textbf{Measurement:} $AE \approx 3.5\ \mathrm{cm}$, so $R \approx \SI{35}{N}$, directed at about $66^{\circ}$ anticlockwise from east.

\textbf{Equilibrant:} $E \approx \SI{35}{N}$ along the same line but in the opposite sense, i.e.\ at $66^{\circ} + 180^{\circ} = 246^{\circ}$ (or $66^{\circ}$ below the westward axis). Adding this force would close the polygon.

\textbf{Check (components):} 
\begin{align*}
\sum F_x &= 50 + 40\cos 70^{\circ} + 45\cos 150^{\circ} + 30\cos 250^{\circ} \\
&= 50 + 13.68 - 38.97 - 10.26 = \SI{14.45}{N}, \\
\sum F_y &= 0 + 40\sin 70^{\circ} + 45\sin 150^{\circ} + 30\sin 250^{\circ} \\
&= 37.59 + 22.5 - 28.19 = \SI{31.90}{N}. \\
R &= \sqrt{14.45^2 + 31.90^2} \approx \SI{35.0}{N}, \\
\theta &= \arctan\left(\frac{31.90}{14.45}\right) \approx 65.7^{\circ} \text{ from east}.
\end{align*}
The analytical result (\SI{35.0}{N} at $65.7^{\circ}$) agrees closely with the graphical measurement (\SI{35}{N} at $66^{\circ}$), the small difference coming from drawing tolerances.
\end{exoresolu}

\begin{attention}[Common traps in graphical questions]
\begin{itemize}
\item Mixing up the \emph{space diagram} and the \emph{force diagram}: angles must be transferred correctly from one to the other.
\item Drawing forces tail-to-tail (all from one point) in the polygon — that construction shows the forces but not their sum.
\item Forgetting that the polygon may be drawn in \emph{any order}; choosing a convenient order (e.g.\ nearly opposite forces together) reduces cumulative drawing error.
\item Quoting the equilibrant in the same direction as the resultant. Remember: $\vec{E} = -\vec{R}$.
\item Using a scale that is too small, leading to large percentage errors.
\end{itemize}
\end{attention}

\begin{exercice}[Graphical practice]
\begin{enumerate}
\item Two forces of \SI{60}{N} and \SI{80}{N} act at a point with an angle of $90^{\circ}$ between them. Find the resultant graphically (scale $1\ \mathrm{cm} \equiv 10\ \mathrm{N}$) and check by calculation.
\item Three forces act at a point: \SI{30}{N} east, \SI{40}{N} at $120^{\circ}$, \SI{35}{N} at $230^{\circ}$. Draw the force polygon, decide whether the system is in equilibrium, and if not find the resultant and the equilibrant.
\item A picture hangs from two strings making angles of $40^{\circ}$ and $55^{\circ}$ with the horizontal, the tensions being \SI{25}{N} and \SI{30}{N}. Using a polygon, find the single upward force (the equilibrant of the two tensions) that the nail must provide, and comment on the weight of the picture.
\end{enumerate}
\end{exercice}

\begin{corrige}[Exercise 1]
\textbf{1.} Parallelogram (here a rectangle): diagonal measures $10.0\ \mathrm{cm}$, so $R = \SI{100}{N}$ at $\arctan(80/60) \approx 53^{\circ}$ to the \SI{60}{N} force. Check: $R = \sqrt{60^2 + 80^2} = \SI{100}{N}$. \checkmark

\textbf{2.} With scale $1\ \mathrm{cm} \equiv 10\ \mathrm{N}$, draw $3.0\ \mathrm{cm}$ east, then $4.0\ \mathrm{cm}$ at $120^{\circ}$, then $3.5\ \mathrm{cm}$ at $230^{\circ}$. The polygon does not close; the gap measures about $1.5\ \mathrm{cm}$, so $R \approx \SI{15}{N}$ directed roughly $32^{\circ}$ north of west; the equilibrant is \SI{15}{N} in the opposite direction ($32^{\circ}$ south of east). (Component check: $\sum F_x = 30 + 40\cos 120^{\circ} + 35\cos 230^{\circ} = 30 - 20 - 22.5 = -12.5\ \mathrm{N}$; $\sum F_y = 0 + 40\sin 120^{\circ} + 35\sin 230^{\circ} = 34.64 - 26.81 = 7.83\ \mathrm{N}$; $R = \sqrt{12.5^2 + 7.83^2} \approx \SI{14.7}{N}$ at $\arctan(7.83/12.5) \approx 32^{\circ}$ above the negative $x$-axis, i.e. $148^{\circ}$ from east. Graphical measurement improves with a larger scale such as $2\ \mathrm{cm} \equiv 10\ \mathrm{N}$.)

\textbf{3.} Draw the two tensions head-to-tail along their string directions; the closing side is vertical, of length corresponding to about $\SI{41}{N}$ upward. The nail's reaction (equilibrant of the tensions) is $\approx \SI{41}{N}$ upward; for the picture to hang at rest its weight must equal this, so the picture weighs about \SI{41}{N}.
\end{corrige}

\begin{aretenir}[Key points]
\begin{itemize}
\item The \cle{resultant} $\vec{R}$ is the single force equivalent to the whole system; the \cle{equilibrant} $\vec{E} = -\vec{R}$ cancels it and produces equilibrium.
\item Two forces: use the \textbf{parallelogram}; several coplanar forces: use the \textbf{polygon of forces} drawn head-to-tail to scale.
\item In the polygon, the resultant runs from the \emph{start} to the \emph{finish}, against the sense of the other arrows.
\item \textbf{Equilibrium test:} the force polygon closes $\iff \vec{R} = \vec{0} \iff$ equilibrium.
\item Always state the scale, draw large, measure carefully, and quote answers to a sensible precision.
\item Bow's notation links the space diagram and force diagram unambiguously.
\end{itemize}
\end{aretenir}

Graphical methods give fast, visual answers, but their accuracy is limited by the drawing. In the next section we develop \emph{analytical} methods — resolving forces into components — which give exact results and handle equilibrium problems algebraically.

\coursec{Analytical Methods and Equilibrium Problems}

In the previous section we explored graphical methods — the triangle and polygon of forces — which give an excellent visual understanding of resultants and equilibrium. However, graphical methods are limited by drawing accuracy and become cumbersome when many forces are involved. In this section we develop the \emph{analytical} (algebraic) approach: using trigonometry and vector components to compute resultants exactly and to solve equilibrium problems with precision. This is the method you will use in the vast majority of GCE Advanced Level examination questions.

\courssub{Analytical Resultant of Two Forces}

Consider two forces \(\vec{P}\) and \(\vec{Q}\) acting at a point \(O\), with an angle \(\theta\) between them (\(0^\circ \le \theta \le 180^\circ\)). Their resultant \(\vec{R}\) is the diagonal of the parallelogram formed by \(\vec{P}\) and \(\vec{Q}\). Applying the cosine rule to the triangle formed by \(\vec{P}\), \(\vec{Q}\) and \(\vec{R}\) gives the magnitude:

\begin{propriete}[Magnitude of the Resultant of Two Forces]
\[
R = \sqrt{P^2 + Q^2 + 2PQ\cos\theta}
\]
\end{propriete}

The direction of \(\vec{R}\) relative to \(\vec{P}\) is given by the angle \(\alpha\) such that
\[
\tan\alpha = \frac{Q\sin\theta}{P + Q\cos\theta}
\]
This follows from the sine rule or by resolving \(\vec{Q}\) into components parallel and perpendicular to \(\vec{P}\).

\begin{exemplebox}[Resultant of two forces at \(60^\circ\)]
Two forces of magnitudes \(P = 30\,\mathrm{N}\) and \(Q = 40\,\mathrm{N}\) act at a point with an angle of \(60^\circ\) between them. Find the magnitude and direction of their resultant.
\tcblower
Magnitude:
\[
R = \sqrt{30^2 + 40^2 + 2\cdot30\cdot40\cos60^\circ}
   = \sqrt{900 + 1600 + 2400\cdot0.5}
   = \sqrt{3700} \approx 60.83\,\mathrm{N}
\]
Direction (angle \(\alpha\) with the \(30\,\mathrm{N}\) force):
\[
\tan\alpha = \frac{40\sin60^\circ}{30 + 40\cos60^\circ}
           = \frac{40\cdot\sqrt{3}/2}{30 + 40\cdot0.5}
           = \frac{20\sqrt{3}}{50} = \frac{2\sqrt{3}}{5} \approx 0.6928
\]
\(\alpha \approx 34.7^\circ\).
\end{exemplebox}

\textbf{Special Cases}\par
\begin{itemize}
  \item \textbf{Forces at right angles (\(\theta = 90^\circ\)):} \(\cos90^\circ=0\), so \(R = \sqrt{P^2+Q^2}\) and \(\tan\alpha = Q/P\). This is simply Pythagoras' theorem.
  \item \textbf{Equal forces (\(P=Q\)):} \(R = P\sqrt{2(1+\cos\theta)} = 2P\cos\frac{\theta}{2}\), and the resultant bisects the angle between them (\(\alpha = \theta/2\)).
  \item \textbf{Same direction (\(\theta = 0^\circ\)):} \(R = P+Q\), \(\alpha = 0^\circ\).
  \item \textbf{Opposite directions (\(\theta = 180^\circ\)):} \(R = |P-Q|\), direction is that of the larger force.
\end{itemize}

\courssub{Component Method for Several Forces}

When more than two forces act at a point, the parallelogram law becomes impractical. The systematic approach is to resolve each force into components along two perpendicular axes (usually horizontal \(x\) and vertical \(y\)).

\begin{methode}[Component Method for the Resultant of Several Coplanar Forces]
\begin{enumerate}
  \item Choose a coordinate system (usually \(x\) horizontal, \(y\) vertical).
  \item For each force \(F_i\) making an angle \(\theta_i\) with the positive \(x\)-axis, compute its components:
        \[
        F_{ix} = F_i\cos\theta_i,\qquad F_{iy} = F_i\sin\theta_i
        \]
  \item Sum all \(x\)-components and all \(y\)-components:
        \[
        \Sigma F_x = \sum F_{ix},\qquad \Sigma F_y = \sum F_{iy}
        \]
  \item The resultant \(\vec{R}\) has components \(R_x = \Sigma F_x\), \(R_y = \Sigma F_y\). Its magnitude and direction are:
        \[
        R = \sqrt{(\Sigma F_x)^2 + (\Sigma F_y)^2},\qquad
        \tan\theta_R = \frac{\Sigma F_y}{\Sigma F_x}
        \]
        (Determine the correct quadrant for \(\theta_R\) from the signs of \(\Sigma F_x\) and \(\Sigma F_y\).)
\end{enumerate}
\end{methode}

\begin{exemplebox}[Resultant of three forces by components]
Three forces act at a point: \(F_1 = 10\,\mathrm{N}\) at \(0^\circ\), \(F_2 = 20\,\mathrm{N}\) at \(120^\circ\), \(F_3 = 15\,\mathrm{N}\) at \(210^\circ\). Find the resultant.
\tcblower
Components:
\[
\begin{aligned}
F_{1x} &= 10\cos0^\circ = 10, & F_{1y} &= 10\sin0^\circ = 0\\
F_{2x} &= 20\cos120^\circ = -10, & F_{2y} &= 20\sin120^\circ = 10\sqrt{3} \approx 17.32\\
F_{3x} &= 15\cos210^\circ = -15\frac{\sqrt{3}}{2} \approx -12.99, & F_{3y} &= 15\sin210^\circ = -7.5
\end{aligned}
\]
\[
\Sigma F_x = 10 - 10 - 12.99 = -12.99,\qquad
\Sigma F_y = 0 + 17.32 - 7.5 = 9.82
\]
\[
R = \sqrt{(-12.99)^2 + 9.82^2} \approx \sqrt{168.7 + 96.4} \approx 16.28\,\mathrm{N}
\]
\[
\tan\theta_R = \frac{9.82}{-12.99} \approx -0.756 \quad\Rightarrow\quad \theta_R \approx 180^\circ - 37.1^\circ = 142.9^\circ
\]
(Second quadrant because \(\Sigma F_x < 0\), \(\Sigma F_y > 0\).)
\end{exemplebox}

\courssub{Conditions of Equilibrium in Analytical Form}

A particle (or a rigid body with no net moment) is in equilibrium under a system of coplanar forces if and only if the vector sum of the forces is zero. In components this gives two scalar equations:

\begin{propriete}[Equilibrium Conditions for a Particle]
\[
\Sigma F_x = 0 \quad\text{and}\quad \Sigma F_y = 0
\]
\end{propriete}

The force that balances a system of forces is called the \cle{equilibrant}. It is equal in magnitude and opposite in direction to the resultant: \(\vec{E} = -\vec{R}\). For a system in equilibrium, the equilibrant is the force that would need to be applied to maintain equilibrium if it were not already present.

For an \emph{extended rigid body}, equilibrium also requires that the sum of moments about any point be zero:
\[
\Sigma M_O = 0
\]
This will be essential when we deal with beams, ladders, and other extended bodies in later chapters. For now, we concentrate on particles and concurrent force systems where the moment condition is automatically satisfied if the forces meet at a point.

\courssub{Lami's Theorem for Three Concurrent Forces}

When exactly three forces act at a point and the system is in equilibrium, the forces can be represented by the three sides of a triangle taken in order (the triangle of forces). Applying the sine rule to this triangle yields a powerful shortcut known as \textbf{Lami's Theorem}.

\begin{definition}[Lami's Theorem]
If three forces \(P\), \(Q\), \(R\) acting at a point are in equilibrium, then each force is proportional to the sine of the angle between the other two forces:
\[
\frac{P}{\sin\alpha} = \frac{Q}{\sin\beta} = \frac{R}{\sin\gamma}
\]
where \(\alpha\) is the angle between \(Q\) and \(R\), \(\beta\) is the angle between \(R\) and \(P\), and \(\gamma\) is the angle between \(P\) and \(Q\).
\end{definition}

\begin{popfigure}
\begin{tikzpicture}[scale=1.2, >=stealth]
% Force vectors from origin
\coordinate (O) at (0,0);
% Forces: P along +x, Q at 120°, R at -120° (equilibrium)
\draw[->, thick, popblue] (O) -- (2.5,0) node[midway, below] {$P$};
\draw[->, thick, poporange] (O) -- (-1.25,2.165) node[midway, left] {$Q$};
\draw[->, thick, popgreen] (O) -- (-1.25,-2.165) node[midway, left] {$R$};
% Angles between forces
\draw[popmuted, <->] (0.5,0) arc (0:120:0.5) node[midway, right=2pt] {$\alpha$};
\draw[popmuted, <->] (-0.3,0.5) arc (120:240:0.6) node[midway, left] {$\beta$};
\draw[popmuted, <->] (0.5,0) arc (0:-120:0.5) node[midway, right=2pt] {$\gamma$};
% Force triangle (shifted right)
\begin{scope}[shift={(5,0)}]
\coordinate (A) at (0,0);
\coordinate (B) at (2.5,0);
\coordinate (C) at (1.25,2.165);
\draw[thick, popblue] (A) -- (B) node[midway, below] {$P$};
\draw[thick, poporange] (B) -- (C) node[midway, right] {$Q$};
\draw[thick, popgreen] (C) -- (A) node[midway, left] {$R$};
% Angles in triangle (interior angles)
\draw[popmuted, <->] (0.4,0) arc (0:60:0.4) node[midway, right=2pt] {$180^\circ-\gamma$};
\draw[popmuted, <->] (2.1,0.2) arc (180:120:0.4) node[midway, right=2pt] {$180^\circ-\alpha$};
\draw[popmuted, <->] (1.25,1.8) arc (-60:-120:0.4) node[midway, left] {$180^\circ-\beta$};
\node at (1.25,-0.8) {Force Triangle};
\end{scope}
\node at (1.25,-1.2) {Forces at Point};
\end{tikzpicture}
\legende{Derivation of Lami's theorem: the angles \(\alpha,\beta,\gamma\) between the forces at the point are the supplements of the interior angles of the force triangle. Since \(\sin(180^\circ-\theta)=\sin\theta\), the sine rule on the triangle gives \(P/\sin\alpha = Q/\sin\beta = R/\sin\gamma\).}
\end{popfigure}

\begin{propriete}[Derivation of Lami's Theorem (Examinable)]
The three forces in equilibrium form a closed triangle when placed head-to-tail. The interior angles of this triangle are \(180^\circ-\alpha\), \(180^\circ-\beta\), \(180^\circ-\gamma\). By the sine rule:
\[
\frac{P}{\sin(180^\circ-\alpha)} = \frac{Q}{\sin(180^\circ-\beta)} = \frac{R}{\sin(180^\circ-\gamma)}
\]
Since \(\sin(180^\circ-\theta)=\sin\theta\), we obtain Lami's theorem.
\end{propriete}

\begin{attention}[The Angle in Lami's Theorem]
The angle used with each force is the angle \textbf{between the other two forces}, not the angle that force makes with the horizontal or vertical. In the diagram above, \(\alpha\) is the angle between \(Q\) and \(R\), \(\beta\) between \(R\) and \(P\), \(\gamma\) between \(P\) and \(Q\). Misidentifying these angles is the most common error when applying Lami's theorem.
\end{attention}

\courssub{Applications: Solving Equilibrium Problems}

We now tackle typical examination problems: a body suspended by two strings, and a body on a smooth inclined plane held by a force. For three forces, you have a choice of method.

\begin{methode}[Choosing a Method for Three-Force Equilibrium]
\begin{enumerate}
  \item \textbf{Lami's Theorem} is fastest when the three forces are concurrent and you know (or can easily find) the angles between each pair of forces. Write the three ratios and solve for the unknowns.
  \item \textbf{Resolution (\(\Sigma F_x=0,\ \Sigma F_y=0\))} is more versatile: it works for any number of forces, and is often simpler when forces are given with respect to horizontal/vertical axes (e.g. strings at known angles to the horizontal). It also avoids the supplement-angle confusion of Lami's theorem.
  \item For \textbf{more than three forces}, you \textbf{must} use the component method; Lami's theorem does not apply.
\end{enumerate}
\end{methode}

\begin{exoresolu}[Suspended Load — Two Strings]
A load of weight \(W = 50\,\mathrm{N}\) is suspended by two light inextensible strings attached to the same horizontal ceiling. The strings make angles of \(30^\circ\) and \(45^\circ\) with the horizontal. Find the tension in each string.
\tcblower
\textbf{Method 1: Resolution (recommended)}
Let \(T_1\) be the tension in the string at \(30^\circ\), \(T_2\) at \(45^\circ\). The weight acts vertically downward.
Resolve horizontally: \(\Sigma F_x = 0 \Rightarrow T_1\cos30^\circ - T_2\cos45^\circ = 0\)
Resolve vertically: \(\Sigma F_y = 0 \Rightarrow T_1\sin30^\circ + T_2\sin45^\circ - 50 = 0\)

From horizontal: \(T_1\frac{\sqrt{3}}{2} = T_2\frac{\sqrt{2}}{2} \Rightarrow T_1 = T_2\frac{\sqrt{2}}{\sqrt{3}}\).

Substitute into vertical:
\[
T_2\frac{\sqrt{2}}{\sqrt{3}}\cdot\frac{1}{2} + T_2\frac{\sqrt{2}}{2} = 50
\]
\[
T_2\frac{\sqrt{2}}{2}\left(\frac{1}{\sqrt{3}} + 1\right) = 50
\]
\[
T_2 = \frac{100}{\sqrt{2}\left(1+\frac{1}{\sqrt{3}}\right)} = \frac{100\sqrt{3}}{\sqrt{2}(\sqrt{3}+1)} \approx 36.6\,\mathrm{N}
\]
Then \(T_1 = T_2\frac{\sqrt{2}}{\sqrt{3}} \approx 29.9\,\mathrm{N}\).

\textbf{Method 2: Lami's Theorem}
The three forces are \(T_1\), \(T_2\), and \(W=50\,\mathrm{N}\). The angles between them:
- Between \(T_1\) and \(T_2\): \(30^\circ+45^\circ = 75^\circ\).
- Between \(T_2\) and \(W\): \(90^\circ+45^\circ = 135^\circ\) (since \(W\) is vertical down, \(T_2\) is \(45^\circ\) above horizontal).
- Between \(W\) and \(T_1\): \(90^\circ+30^\circ = 120^\circ\).

Lami's theorem:
\[
\frac{T_1}{\sin135^\circ} = \frac{T_2}{\sin120^\circ} = \frac{50}{\sin75^\circ}
\]
\[
T_1 = 50\,\frac{\sin135^\circ}{\sin75^\circ} = 50\,\frac{\sqrt{2}/2}{\sin75^\circ} \approx 29.9\,\mathrm{N}
\]
\[
T_2 = 50\,\frac{\sin120^\circ}{\sin75^\circ} = 50\,\frac{\sqrt{3}/2}{\sin75^\circ} \approx 36.6\,\mathrm{N}
\]
Both methods give the same result. Notice how Lami's theorem requires the angles \emph{between} the forces (75°, 135°, 120°), not the angles to the horizontal.
\end{exoresolu}

\begin{popfigure}
\begin{tikzpicture}[scale=1.1, >=stealth]
% Physical setup
\draw[thick] (-3,2) -- (3,2); % ceiling
\coordinate (O) at (0,0);
\draw[->, thick, popblue] (O) -- (-1.732,1) node[midway, left] {$T_1$};
\draw[->, thick, poporange] (O) -- (1.414,1.414) node[midway, right] {$T_2$};
\draw[->, thick, popred] (O) -- (0,-1.5) node[midway, right] {$W=50\,\mathrm{N}$};
% Angles to horizontal
\draw[popmuted, <->] (0,0.2) arc (90:60:0.5) node[midway, left=2pt] {$30^\circ$};
\draw[popmuted, <->] (0.2,0) arc (0:45:0.5) node[midway, below=2pt] {$45^\circ$};
% Force triangle (shifted right)
\begin{scope}[shift={(6,0)}]
\coordinate (A) at (0,0);
\coordinate (B) at (-1.732,1);
\coordinate (C) at (0,-1.5);
\draw[thick, popblue] (A) -- (B) node[midway, left] {$T_1$};
\draw[thick, popred] (B) -- (C) node[midway, right] {$W$};
\draw[thick, poporange] (C) -- (A) node[midway, right] {$T_2$};
% Angles in triangle (interior angles)
\draw[popmuted, <->] (-0.3,0.2) arc (150:90:0.4) node[midway, left] {$105^\circ$};
\draw[popmuted, <->] (-1.5,0.8) arc (-30:-105:0.4) node[midway, left] {$45^\circ$};
\draw[popmuted, <->] (0.1,-1.3) arc (-90:-15:0.4) node[midway, right] {$30^\circ$};
\node at (0,-2) {Force Triangle};
\end{scope}
\end{tikzpicture}
\legende{Suspended load: physical configuration (left) and force triangle (right). The interior angles of the force triangle are the supplements of the angles between the forces at the point (75°, 135°, 120°).}
\end{popfigure}

\begin{exemplebox}[Body on a Smooth Inclined Plane]
A body of weight \(W = 100\,\mathrm{N}\) rests on a smooth plane inclined at \(30^\circ\) to the horizontal. It is held in equilibrium by a force \(F\) acting up the plane at an angle of \(20^\circ\) to the plane. Find \(F\) and the reaction \(R\) of the plane.
\tcblower
Since the plane is smooth, the reaction \(R\) is perpendicular to the plane. Choose axes: \(x\) up the plane, \(y\) perpendicular to the plane (outward).
Forces:
- Weight \(W\): components \(W_x = -W\sin30^\circ = -50\,\mathrm{N}\), \(W_y = -W\cos30^\circ = -50\sqrt{3}\,\mathrm{N}\).
- Applied force \(F\): \(F_x = F\cos20^\circ\), \(F_y = F\sin20^\circ\).
- Reaction \(R\): \(R_x = 0\), \(R_y = R\).

Equilibrium:
\[
\Sigma F_x = F\cos20^\circ - 50 = 0 \;\Rightarrow\; F = \frac{50}{\cos20^\circ} \approx 53.2\,\mathrm{N}
\]
\[
\Sigma F_y = F\sin20^\circ + R - 50\sqrt{3} = 0 \;\Rightarrow\; R = 50\sqrt{3} - F\sin20^\circ \approx 86.6 - 18.2 = 68.4\,\mathrm{N}
\]
\end{exemplebox}

\courssub{Summary and Key Points}

\begin{aretenir}[Key Points for Analytical Methods and Equilibrium]
\begin{itemize}
\item Resultant of two forces: \(R = \sqrt{P^2+Q^2+2PQ\cos\theta}\), \(\tan\alpha = \frac{Q\sin\theta}{P+Q\cos\theta}\).
\item For several forces: resolve into \(x\) and \(y\) components, sum them, then \(R = \sqrt{(\Sigma F_x)^2+(\Sigma F_y)^2}\), \(\tan\theta_R = \Sigma F_y/\Sigma F_x\).
\item Equilibrium of a particle: \(\Sigma F_x = 0\) and \(\Sigma F_y = 0\). The equilibrant \(\vec{E} = -\vec{R}\).
\item Lami's Theorem (three concurrent forces in equilibrium): \(\frac{P}{\sin\alpha} = \frac{Q}{\sin\beta} = \frac{R}{\sin\gamma}\), where \(\alpha,\beta,\gamma\) are the angles \textbf{between the other two forces}.
\item For three forces, choose resolution (when angles to axes are given) or Lami's theorem (when angles between forces are given). For more than three forces, always use resolution.
\end{itemize}
\end{aretenir}

\begin{exercice}[Practice Problems]
\begin{itemize}
\item Two forces of \(25\,\mathrm{N}\) and \(40\,\mathrm{N}\) act at a point with an angle of \(120^\circ\) between them. Find the magnitude and direction of their resultant.
\item A particle is in equilibrium under three forces: \(F_1 = 10\,\mathrm{N}\) horizontally to the right, \(F_2 = 15\,\mathrm{N}\) at \(120^\circ\) to the horizontal, and \(F_3\). Find \(F_3\) by (a) resolution, (b) Lami's theorem.
\item A weight of \(80\,\mathrm{N}\) is suspended by two strings of lengths \(3\,\mathrm{m}\) and \(4\,\mathrm{m}\) attached to two points on a horizontal ceiling \(5\,\mathrm{m}\) apart. Find the tension in each string. (Hint: the strings and the ceiling form a 3-4-5 triangle.)
\item A block of weight \(W\) rests on a smooth plane inclined at \(\alpha\). A horizontal force \(H\) keeps it in equilibrium. Express \(H\) and the reaction \(R\) in terms of \(W\) and \(\alpha\).
\end{itemize}
\end{exercice}

\begin{corrige}[Exercise 1]
\(R = \sqrt{25^2+40^2+2\cdot25\cdot40\cos120^\circ} = \sqrt{625+1600-1000} = \sqrt{1225} = 35\,\mathrm{N}\).
\(\tan\alpha = \frac{40\sin120^\circ}{25+40\cos120^\circ} = \frac{40\cdot\sqrt{3}/2}{25-20} = 4\sqrt{3} \Rightarrow \alpha \approx 81.8^\circ\) from the \(25\,\mathrm{N}\) force.
\end{corrige}

\begin{corrige}[Exercise 2]
(a) Resolution: \(\Sigma F_x = 10 + 15\cos120^\circ + F_{3x} = 10 - 7.5 + F_{3x} = 0 \Rightarrow F_{3x} = -2.5\).
\(\Sigma F_y = 0 + 15\sin120^\circ + F_{3y} = 12.99 + F_{3y} = 0 \Rightarrow F_{3y} = -12.99\).
\(F_3 = \sqrt{2.5^2+12.99^2} \approx 13.23\,\mathrm{N}\), angle \(\approx -100.9^\circ\) (or \(259.1^\circ\)).
(b) Lami's theorem: angles between forces: between \(F_1\) and \(F_2\) is \(120^\circ\); between \(F_2\) and \(F_3\) is \(\beta\); between \(F_3\) and \(F_1\) is \(\gamma\). Since \(F_1\) and \(F_2\) are known, the force triangle gives \(F_3/\sin120^\circ = 10/\sin\beta = 15/\sin\gamma\). With \(\beta+\gamma = 240^\circ\), solving yields same \(F_3\).
\end{corrige}

\begin{corrige}[Exercise 3]
The geometry gives a right triangle (3-4-5). The strings are perpendicular to each other (angle \(90^\circ\)). Let \(T_3\) be the tension in the \(3\,\mathrm{m}\) string, \(T_4\) in the \(4\,\mathrm{m}\) string. The angles with the horizontal are \(\sin^{-1}(4/5)\approx53.1^\circ\) for the \(3\,\mathrm{m}\) string and \(\sin^{-1}(3/5)\approx36.9^\circ\) for the \(4\,\mathrm{m}\) string. At the point of suspension, the angles between the forces are: between the two tensions \(90^\circ\); between \(T_3\) and \(W\) is \(36.9^\circ\); between \(T_4\) and \(W\) is \(53.1^\circ\). By Lami's theorem:
\[
\frac{T_3}{\sin53.1^\circ} = \frac{T_4}{\sin36.9^\circ} = \frac{80}{\sin90^\circ} = 80
\]
\[
T_3 = 80 \times 0.8 = 64\,\mathrm{N},\qquad T_4 = 80 \times 0.6 = 48\,\mathrm{N}.
\]
\end{corrige}

\begin{corrige}[Exercise 4]
Axes along and perpendicular to plane. Weight components: \(W\sin\alpha\) down plane, \(W\cos\alpha\) into plane. Horizontal force \(H\): component up plane \(H\cos\alpha\), component out of plane \(H\sin\alpha\).
\(\Sigma F_x = H\cos\alpha - W\sin\alpha = 0 \Rightarrow H = W\tan\alpha\).
\(\Sigma F_y = R - W\cos\alpha - H\sin\alpha = 0 \Rightarrow R = W\cos\alpha + W\tan\alpha\sin\alpha = W(\cos\alpha + \sin^2\alpha/\cos\alpha) = W/\cos\alpha\).
\end{corrige}

\coursec{Hooke's Law and Its Experimental Verification}

In the previous sections we learnt how to find the resultant of a system of coplanar forces analytically and how to use the conditions of equilibrium to solve problems. In most of those problems the forces were given to us directly: a weight of $50\ \mathrm{N}$, a tension of $30\ \mathrm{N}$, and so on. But where do such forces come from, and how can we \emph{measure} a force that we cannot see? One of the oldest and most useful answers is: \textbf{use a spring}. A spring stretches when you pull it, and the amount it stretches tells you how hard you are pulling. The precise relationship between the pull and the stretch is the subject of this section: \cle{Hooke's law}. It is the principle behind every spring balance used in the markets of Bamenda, Buea or Douala to weigh vegetables, fish and meat.

\courssub{Elasticity, Springs and Elastic Strings}

\textbf{An introductory observation.}\par
Take a spiral spring, hang it from a clamp, and pull its lower end gently. Two things happen:
\begin{itemize}
\item the spring gets \emph{longer} while you pull;
\item when you release it, the spring returns to its original length.
\end{itemize}
If you pull harder, it stretches more. But if you pull \emph{too} hard, the spring no longer returns to its original length: it has been permanently damaged. This everyday behaviour contains all the physics of this section.

\begin{definition}[Elasticity, natural length and extension]
A body is said to be \cle{elastic} if it returns to its original shape and size when the deforming force is removed. For a spring or an elastic string we define:
\begin{itemize}
\item the \cle{natural length} (or unstretched length) $l$, sometimes written $l_0$: the length of the spring when no load is applied;
\item the \cle{stretched length} $L$: the length of the spring when a load is applied;
\item the \cle{extension} $x$: the increase in length,
\[ x = L - l. \]
\end{itemize}
The extension is measured in metres (m), or in millimetres (mm) in the laboratory.
\end{definition}

\begin{attention}[Extension is NOT the total length]
The most frequent mistake in this topic is to confuse the \emph{extension} with the \emph{total length} of the spring. Always subtract the natural length first:
\[ x = \text{stretched length} - \text{natural length} = L - l. \]
For example, a spring of natural length $l = 20\ \mathrm{cm}$ stretched to $L = 26\ \mathrm{cm}$ has extension $x = 26 - 20 = 6\ \mathrm{cm}$, \textbf{not} $26\ \mathrm{cm}$.
\end{attention}

When a stretched spring is in equilibrium, it pulls back on the object stretching it. This inward pull is called the \cle{tension} (or \emph{thrust}, when a spring is compressed and pushes back). We denote it by $T$. For a body hanging at rest from a spring, equilibrium of the body gives $T = mg$: the tension in the spring equals the weight of the load.

\begin{definition}[Spring constant and modulus of elasticity]
For a helical spring, the constant of proportionality in Hooke's law is called the \cle{spring constant} (or stiffness) $k$, with unit $\mathrm{N\,m^{-1}}$.
For an elastic string or wire, it is more convenient to use the \cle{modulus of elasticity} $\lambda$ (lambda), measured in newtons (N), which is the force required to double the length of the string (theoretically). The two are related by
\[ \lambda = k l, \]
where $l$ is the natural length. Thus Hooke's law for a string is written
\[ T = \frac{\lambda x}{l}. \]
\end{definition}

\courssub{Statement of Hooke's Law}

Let us now make the observation quantitative. Hang masses of increasing weight $F$ from a spring and measure the extension $x$ each time. You will find results like these (typical school-laboratory values):

\begin{center}
\begin{tabularx}{\linewidth}{|l|X|X|X|X|X|}
\hline
\rowcolor{popblueL} Load $F$ (N) & $0$ & $1.0$ & $2.0$ & $3.0$ & $4.0$ \\
\hline Extension $x$ (mm) & $0$ & $25$ & $50$ & $75$ & $100$ \\
\hline
\end{tabularx}
\end{center}

Each time the load doubles, the extension doubles; each time the load triples, the extension triples. The ratio $\dfrac{F}{x}$ is constant (here $\dfrac{F}{x} = 40\ \mathrm{N/m}$). This is exactly the behaviour discovered by Robert Hooke in the seventeenth century.

\begin{propriete}[Hooke's law]
The \cle{tension} $T$ in an elastic spring (or string) is directly proportional to its \cle{extension} $x$, \emph{provided the limit of proportionality is not exceeded}:
\[ T = kx, \]
where $k$ is the \cle{spring constant} (stiffness) of the spring.

An equivalent form, used for elastic strings and wires, is
\[ T = \frac{\lambda x}{l}, \]
where $l$ is the natural length and $\lambda$ is the \cle{modulus of elasticity} of the string (measured in newtons).
\end{propriete}

\begin{aretenir}[What Hooke's law really says]
\begin{itemize}
\item $T = kx$ is a \emph{linear} relationship: the graph of $T$ against $x$ is a straight line through the origin.
\item It holds only \emph{within the limit of proportionality} (often called the elastic limit at this level). Beyond that limit the spring is permanently deformed and the law fails.
\item The same law applies to a spring in \emph{compression}, with $x$ the decrease in length and $T$ the thrust.
\item For an elastic string, $T = \lambda x / l$; a string, unlike a spring, cannot push — it goes slack when $x \le 0$.
\item The spring constant $k$ is the force needed to produce a \emph{unit extension}: $k = T/x$. Its SI unit is $\mathrm{N\,m^{-1}}$. Physically, $k$ measures the \cle{stiffness}: the larger $k$, the harder the spring is to stretch.
\end{itemize}
\end{aretenir}

\begin{savaistu}[Springs all around us]
Robert Hooke announced his law in 1676 as a Latin anagram, \emph{ceiiinosssttuv}, which he later revealed as \emph{ut tensio, sic vis} — ``as the extension, so the force''. Today Hooke's law is at work everywhere: in the suspension springs of the \emph{bendskins} and taxis on Cameroonian roads, in spring balances at the market, in mattresses, ball-point pens and weighing scales. Engineers designing bridges and buildings must always check that materials stay within their elastic limit, so that structures spring back instead of collapsing.
\end{savaistu}

\courssub{Experimental Verification of Hooke's Law}

\begin{experience}[Verifying Hooke's law with a spiral spring]
\textbf{Aim.} To verify that the extension of a spring is directly proportional to the applied load, and to determine the spring constant $k$.

\textbf{Apparatus.} A spiral spring, a retort stand with clamp, a set of slotted masses (each of known mass, e.g.\ $100\ \mathrm{g}$), a millimetre ruler fixed vertically beside the spring, and a light pointer (a pin or a horizontal needle) attached to the bottom of the spring.

\textbf{Setup.} The spring is hung from the clamp with the ruler clamped vertically beside it, and the pointer moves in front of the scale (see figure below).

\begin{popfigure}
\begin{tikzpicture}[scale=0.9, every node/.style={font=\small}]
% stand
\draw[fill=popmuted!40, draw=popmuted] (-0.2,0) rectangle (3.4,0.25);
\draw[fill=popmuted!40, draw=popmuted] (0.15,0.25) rectangle (0.4,6.2);
\draw[fill=popmuted!40, draw=popmuted] (0.4,5.9) rectangle (2.6,6.15);
\draw[fill=popmuted, draw=popmuted] (2.45,5.55) rectangle (2.75,6.15);
% ruler
\draw[fill=popblueL, draw=popblue] (3.3,1.2) rectangle (3.7,5.6);
\foreach \y in {1.4,1.8,2.2,2.6,3.0,3.4,3.8,4.2,4.6,5.0,5.4}
  \draw[popblue] (3.3,\y) -- (3.45,\y);
\node[rotate=90, popdark] at (3.95,3.4) {ruler (mm)};
% spring (coil)
\draw[thick, poporange]
  (2.6,5.55)
  \foreach \i in {1,...,8}{ -- ++(0.09,-0.10) -- ++(-0.18,-0.10) -- ++(0.09,-0.10) };
\draw[thick, poporange] (2.6,5.55) -- (2.6,3.15);
% pointer
\draw[thick, popdark] (2.6,3.15) -- (3.42,3.15);
\fill[popdark] (3.42,3.15) circle (0.04);
\node[popdark, anchor=west] at (2.0,2.95) {pointer};
% hook + masses
\draw[thick, popdark] (2.6,3.15) -- (2.6,2.85);
\draw[thick, popdark] (2.6,2.85) arc (90:-90:0.12);
\draw[fill=poppurple!35, draw=poppurple] (2.2,2.0) rectangle (3.0,2.6);
\draw[fill=poppurple!35, draw=poppurple] (2.2,1.4) rectangle (3.0,2.0);
\node[popdark] at (2.6,2.3) {$m$};
\node[popdark] at (2.6,1.7) {$m$};
\node[popdark, anchor=west] at (3.1,2.0) {slotted masses};
% labels
\node[popdark] at (1.5,6.45) {clamp};
\node[popdark, anchor=east] at (2.15,4.4) {spring};
% extension arrow
\draw[<->, thick, popgreen] (1.7,3.15) -- (1.7,4.15) node[midway, left, popgreen] {$x$};
\draw[dashed, popgreen] (1.7,4.15) -- (2.55,4.15);
\draw[dashed, popgreen] (1.7,3.15) -- (2.55,3.15);
\end{tikzpicture}
\legende{Experimental setup: a clamped spring with pointer, vertical millimetre scale and slotted masses. The extension $x$ is the displacement of the pointer from its unloaded position.}
\end{popfigure}

\textbf{Procedure.}
\begin{enumerate}
\item With no load, record the pointer reading $r_0$ on the ruler.
\item Hang one slotted mass, wait for the spring to settle, and record the new pointer reading $r$. The extension is $x = r - r_0$.
\item Compute the load $F = mg$ (with $g = 9.81\ \mathrm{m\,s^{-2}}$; for quick school work $g \approx 10\ \mathrm{m\,s^{-2}}$ is often used).
\item Increase the load in equal steps, recording $r$ each time, up to about ten readings.
\item \emph{Unload} step by step and record the readings again; take the mean of the loading and unloading readings for each load.
\item Plot the load $F$ (vertical axis) against the extension $x$ (horizontal axis).
\end{enumerate}

\textbf{Precautions.}
\begin{itemize}
\item Do \emph{not} overload the spring: stay well below the elastic limit, otherwise the spring is permanently damaged and the results are useless.
\item Fix the ruler vertically, close beside the spring, so that the pointer moves in front of the scale.
\item Use a pointer and read the scale with the eye level with the pointer to avoid \emph{parallax error}.
\item Wait for the spring to stop oscillating before each reading, and repeat readings (loading and unloading) to average out random errors.
\item Ensure the spring hangs freely without touching the ruler or the bench.
\end{itemize}
\end{experience}

\begin{methode}[Verifying Hooke's law from the graph]
\begin{enumerate}
\item Plot $F$ (N) on the vertical axis against $x$ (mm or m) on the horizontal axis.
\item If the points lie on a \textbf{straight line through the origin}, then $F \propto x$: Hooke's law is verified.
\item The \textbf{gradient} of the line is the spring constant:
\[ k = \text{gradient} = \frac{\Delta F}{\Delta x}. \]
Choose two points on the line far apart (not necessarily data points) to compute it.
\item If $x$ is in millimetres, convert: a gradient of $a\ \mathrm{N/mm}$ equals $1000a\ \mathrm{N/m}$.
\item If the line eventually \emph{curves} and stops passing through the origin proportionally, the \cle{limit of proportionality} (elastic limit) has been reached: beyond it the spring is permanently deformed.
\end{enumerate}
\end{methode}

\begin{popfigure}
\begin{tikzpicture}
\begin{axis}[
  width=11cm, height=8cm,
  axis lines=left,
  xlabel={Extension $x$ (mm)},
  ylabel={Load $F$ (N)},
  xlabel style={popdark}, ylabel style={popdark},
  xmin=0, xmax=170, ymin=0, ymax=7.5,
  xtick={0,25,50,75,100,125,150},
  ytick={0,1,2,3,4,5,6,7},
  tick label style={font=\small},
  label style={font=\small},
  clip=false]
% linear region
\addplot[very thick, popblue, domain=0:100] {0.04*x};
% beyond elastic limit (curving)
\addplot[very thick, poporange, domain=100:155, samples=40] {4 + 0.09*(x-100) + 0.0016*(x-100)^2};
% dashed marker at elastic limit
\addplot[dashed, popmuted] coordinates {(100,0) (100,4)};
\addplot[dashed, popmuted] coordinates {(0,4) (100,4)};
\fill[popdark] (axis cs:100,4) circle (1.6pt);
\node[popdark, font=\small, anchor=south west] at (axis cs:100,4.05) {limit of proportionality};
\node[popblue, font=\small, rotate=21.8, anchor=south] at (axis cs:52,2.08) {linear region: $F=kx$};
\node[poporange, font=\small, anchor=south west] at (axis cs:112,5.6) {permanent deformation};
% gradient triangle
\addplot[popgreen, thick] coordinates {(25,1) (75,1) (75,3)};
\node[popgreen, font=\small, anchor=north] at (axis cs:50,1) {$\Delta x$};
\node[popgreen, font=\small, anchor=west] at (axis cs:75,2) {$\Delta F$};
\end{axis}
\end{tikzpicture}
\legende{Load against extension for a spring with $k = 40\ \mathrm{N/m}$. Up to the limit of proportionality the graph is a straight line through the origin (gradient $= k$); beyond it the spring stretches more easily and is permanently deformed.}
\end{popfigure}

\textbf{Reading the graph.} In the figure above, the straight-line part has gradient
\[ k = \frac{\Delta F}{\Delta x} = \frac{3 - 1}{75 - 25} = \frac{2\ \mathrm{N}}{50\ \mathrm{mm}} = 0.04\ \mathrm{N/mm} = 40\ \mathrm{N/m}. \]
The point where the line stops being straight (here at $F = 4\ \mathrm{N}$, $x = 100\ \mathrm{mm}$) is the \cle{limit of proportionality}. Beyond it, equal increases of load produce ever larger extensions, and when the load is removed the spring does not return to its natural length: it has suffered a \emph{permanent (plastic) deformation}.

\begin{attention}[Common traps in the experiment and the graph]
\begin{itemize}
\item Plotting \emph{length} instead of \emph{extension}: the line then does not pass through the origin, and its gradient is still $k$ but the proportionality is hidden. Always use $x = L - l$.
\item Mixing units: if $x$ is in mm, the gradient is in N/mm; multiply by $1000$ to get $\mathrm{N/m}$.
\item Including points beyond the limit of proportionality when computing the gradient: use only the straight-line region.
\item Forgetting that the load is the \emph{weight} $F = mg$, not the mass $m$. A mass of $100\ \mathrm{g}$ gives a load of about $0.981\ \mathrm{N}$, not $100\ \mathrm{N}$.
\item Not taking the mean of loading and unloading readings to reduce hysteresis errors.
\end{itemize}
\end{attention}

\courssub{Worked Examples}

\begin{exoresolu}[Using Hooke's law for a spring]
A spring has natural length $25\ \mathrm{cm}$. When a mass of $400\ \mathrm{g}$ is hung from it, its length becomes $33\ \mathrm{cm}$. Taking $g = 10\ \mathrm{m\,s^{-2}}$:
\begin{enumerate}
\item[(a)] Find the spring constant $k$, assuming the elastic limit is not exceeded.
\item[(b)] Find the length of the spring when a mass of $650\ \mathrm{g}$ is hung from it.
\item[(c)] Find the mass that must be hung from the spring to stretch it to a length of $40\ \mathrm{cm}$.
\end{enumerate}
\tcblower
\textbf{Solution.} First convert to SI units: natural length $l = 0.25\ \mathrm{m}$.

(a) The extension is
\[ x = L - l = 0.33 - 0.25 = 0.08\ \mathrm{m}. \]
The load is the weight of the mass:
\[ T = mg = 0.400 \times 10 = 4\ \mathrm{N}. \]
By Hooke's law $T = kx$:
\[ k = \frac{T}{x} = \frac{4}{0.08} = 50\ \mathrm{N/m}. \]

(b) The new tension is $T = 0.650 \times 10 = 6.5\ \mathrm{N}$, so
\[ x = \frac{T}{k} = \frac{6.5}{50} = 0.13\ \mathrm{m} = 13\ \mathrm{cm}. \]
The total length is $L = l + x = 25 + 13 = 38\ \mathrm{cm}$.

(c) A length of $40\ \mathrm{cm}$ means an extension $x = 40 - 25 = 15\ \mathrm{cm} = 0.15\ \mathrm{m}$. Then
\[ T = kx = 50 \times 0.15 = 7.5\ \mathrm{N}, \qquad m = \frac{T}{g} = \frac{7.5}{10} = 0.75\ \mathrm{kg} = 750\ \mathrm{g}. \]
\end{exoresolu}

\begin{exoresolu}[Spring constant from a graph]
In an experiment, a student plots the load $F$ (N) against the extension $x$ (mm) of a spring and obtains a straight line through the origin passing through the point $(60\ \mathrm{mm},\ 3\ \mathrm{N})$.
\begin{enumerate}
\item[(a)] State what this result shows about the spring.
\item[(b)] Calculate the spring constant in $\mathrm{N/m}$.
\item[(c)] The points stop being collinear beyond $F = 6\ \mathrm{N}$. Interpret this observation.
\end{enumerate}
\tcblower
\textbf{Solution.}
(a) Since the graph of $F$ against $x$ is a straight line through the origin, $F$ is directly proportional to $x$: the spring obeys \cle{Hooke's law} within the range tested.

(b) The gradient of the line is the spring constant:
\[ k = \frac{F}{x} = \frac{3\ \mathrm{N}}{60\ \mathrm{mm}} = 0.05\ \mathrm{N/mm} = 0.05 \times 1000 = 50\ \mathrm{N/m}. \]

(c) Beyond $F = 6\ \mathrm{N}$ the spring has passed its \cle{limit of proportionality}: Hooke's law no longer applies, the extension increases faster than the load, and the spring will not return to its natural length when unloaded (permanent deformation).
\end{exoresolu}

\begin{exoresolu}[Equilibrium with a vertical spring]
A light spring of natural length $0.40\ \mathrm{m}$ and spring constant $120\ \mathrm{N/m}$ hangs vertically. A block is attached to its lower end and rests in equilibrium with the spring stretched to a total length of $0.55\ \mathrm{m}$. Taking $g = 9.81\ \mathrm{m\,s^{-2}}$, find the mass of the block.
\tcblower
\textbf{Solution.} The extension is
\[ x = L - l = 0.55 - 0.40 = 0.15\ \mathrm{m}. \]
By Hooke's law the tension in the spring is
\[ T = kx = 120 \times 0.15 = 18\ \mathrm{N}. \]
The block is in equilibrium under two vertical forces: its weight $mg$ downwards and the tension $T$ upwards. Hence
\[ mg = T \quad\Longrightarrow\quad m = \frac{18}{9.81} \approx 1.83\ \mathrm{kg}. \]
Notice how this example links the present section with the equilibrium conditions of the previous sections: Hooke's law gives us the \emph{magnitude} of one of the forces, and equilibrium does the rest.
\end{exoresolu}

\begin{exoresolu}[Elastic string: modulus of elasticity]
An elastic string has natural length $1.2\ \mathrm{m}$ and modulus of elasticity $\lambda = 60\ \mathrm{N}$. 
\begin{enumerate}
\item[(a)] Find the tension in the string when it is stretched to a length of $1.5\ \mathrm{m}$.
\item[(b)] What length corresponds to a tension of $20\ \mathrm{N}$?
\item[(c)] Find the spring constant $k$ of this string.
\end{enumerate}
\tcblower
\textbf{Solution.}
(a) Extension $x = 1.5 - 1.2 = 0.3\ \mathrm{m}$. Using $T = \dfrac{\lambda x}{l}$:
\[ T = \frac{60 \times 0.3}{1.2} = 15\ \mathrm{N}. \]

(b) For $T = 20\ \mathrm{N}$: $x = \dfrac{Tl}{\lambda} = \dfrac{20 \times 1.2}{60} = 0.4\ \mathrm{m}$, so the length is $L = l + x = 1.2 + 0.4 = 1.6\ \mathrm{m}$.

(c) The spring constant is $k = \dfrac{\lambda}{l} = \dfrac{60}{1.2} = 50\ \mathrm{N/m}$. Check: $T = kx = 50 \times 0.3 = 15\ \mathrm{N}$.
\end{exoresolu}

\courssub{Consolidation Exercises}

\begin{exercice}[Exercise 1]
A spring of natural length $20\ \mathrm{cm}$ extends to $26\ \mathrm{cm}$ under a load of $3\ \mathrm{N}$. Assuming Hooke's law applies, find (a) the spring constant in $\mathrm{N/m}$; (b) the extension produced by a load of $7.5\ \mathrm{N}$; (c) the total length of the spring under that load.
\end{exercice}

\begin{exercice}[Exercise 2]
The table below gives the results of an experiment on a spring.
\begin{center}
\begin{tabularx}{\linewidth}{|l|X|X|X|X|X|X|}
\hline
\rowcolor{popblueL} Load $F$ (N) & $0$ & $1$ & $2$ & $3$ & $4$ & $5$ \\
\hline Length $L$ (cm) & $12.0$ & $14.5$ & $17.0$ & $19.5$ & $22.0$ & $24.5$ \\
\hline
\end{tabularx}
\end{center}
\begin{enumerate}
\item[(a)] Copy the table and add a row for the extension $x$.
\item[(b)] Plot $F$ against $x$ and state, with a reason, whether the spring obeys Hooke's law.
\item[(c)] Deduce the spring constant in $\mathrm{N/m}$.
\end{enumerate}
\end{exercice}

\begin{exercice}[Exercise 3]
An elastic string has natural length $1.2\ \mathrm{m}$ and modulus of elasticity $\lambda = 60\ \mathrm{N}$. Find the tension in the string when it is stretched to a length of $1.5\ \mathrm{m}$. What length corresponds to a tension of $20\ \mathrm{N}$?
\end{exercice}

\begin{exercice}[Exercise 4]
A light spring of natural length $30\ \mathrm{cm}$ and spring constant $200\ \mathrm{N/m}$ is compressed until its length is $24\ \mathrm{cm}$. Calculate the thrust exerted by the spring.
\end{exercice}

\begin{exercice}[Exercise 5]
Two identical springs, each of natural length $20\ \mathrm{cm}$ and spring constant $40\ \mathrm{N/m}$, are joined end to end (in series). The combination is suspended vertically and a mass of $500\ \mathrm{g}$ is attached. Find the total extension of the combination. (Take $g = 10\ \mathrm{m\,s^{-2}}$.)
\end{exercice}

\begin{corrige}[Exercise 1]
(a) $x = 26 - 20 = 6\ \mathrm{cm} = 0.06\ \mathrm{m}$; \; $k = \dfrac{T}{x} = \dfrac{3}{0.06} = 50\ \mathrm{N/m}$.
(b) $x = \dfrac{T}{k} = \dfrac{7.5}{50} = 0.15\ \mathrm{m} = 15\ \mathrm{cm}$.
(c) $L = l + x = 20 + 15 = 35\ \mathrm{cm}$.
\end{corrige}

\begin{corrige}[Exercise 2]
(a) The natural length is $l = 12.0\ \mathrm{cm}$ (length at zero load), so the extensions are $0,\ 2.5,\ 5.0,\ 7.5,\ 10.0,\ 12.5\ \mathrm{cm}$.
(b) The graph of $F$ against $x$ is a straight line through the origin (each extra newton adds exactly $2.5\ \mathrm{cm}$), so $F \propto x$: the spring obeys Hooke's law.
(c) $k = \dfrac{F}{x} = \dfrac{1\ \mathrm{N}}{0.025\ \mathrm{m}} = 40\ \mathrm{N/m}$.
\end{corrige}

\begin{corrige}[Exercise 3]
Extension $x = 1.5 - 1.2 = 0.3\ \mathrm{m}$. Using $T = \dfrac{\lambda x}{l}$:
\[ T = \frac{60 \times 0.3}{1.2} = 15\ \mathrm{N}. \]
For $T = 20\ \mathrm{N}$: $x = \dfrac{Tl}{\lambda} = \dfrac{20 \times 1.2}{60} = 0.4\ \mathrm{m}$, so the length is $L = l + x = 1.2 + 0.4 = 1.6\ \mathrm{m}$.
\end{corrige}

\begin{corrige}[Exercise 4]
Compression $x = 30 - 24 = 6\ \mathrm{cm} = 0.06\ \mathrm{m}$. Thrust $T = kx = 200 \times 0.06 = 12\ \mathrm{N}$.
\end{corrige}

\begin{corrige}[Exercise 5]
For springs in series, the tension is the same in each spring. $T = mg = 0.5 \times 10 = 5\ \mathrm{N}$.
Extension of one spring: $x_1 = T/k = 5/40 = 0.125\ \mathrm{m}$.
Total extension $x = 2x_1 = 0.25\ \mathrm{m} = 25\ \mathrm{cm}$.
(Equivalent spring constant $k_{\text{eq}} = k/2 = 20\ \mathrm{N/m}$, then $x = T/k_{\text{eq}} = 5/20 = 0.25\ \mathrm{m}$.)
\end{corrige}

\begin{aretenir}[Key points of the section]
\begin{itemize}
\item Extension: $x = L - l$ (stretched length minus natural length) — never the total length.
\item Hooke's law: $T = kx$ (or $T = \lambda x / l$ for strings), valid only within the \cle{limit of proportionality}.
\item The spring constant $k$ (unit $\mathrm{N/m}$) is the force per unit extension; it measures stiffness.
\item For an elastic string, modulus $\lambda = kl$ (unit N).
\item Experimental verification: plot $F$ against $x$; a straight line through the origin confirms the law, and its gradient is $k$.
\item Beyond the limit of proportionality the graph curves and the spring is permanently deformed.
\item In equilibrium problems, Hooke's law supplies the tension, which is then combined with the conditions of equilibrium ($\sum F = 0$).
\end{itemize}
\end{aretenir}

\coursec{Applications of Hooke's Law and Synthesis}

This section brings together Hooke's law with the equilibrium of forces, moments and couples studied earlier. We solve problems where springs or elastic strings provide one or more of the forces acting on a body. We also examine combinations of springs.

\courssub{Springs in Series and Parallel}

\begin{propriete}[Equivalent spring constant]
\begin{itemize}
\item \textbf{Springs in series:} The tension is the same in each spring. The total extension is the sum of individual extensions.
\[ \frac{1}{k_{\text{eq}}} = \frac{1}{k_1} + \frac{1}{k_2} + \cdots \quad \text{or} \quad k_{\text{eq}} = \left( \sum \frac{1}{k_i} \right)^{-1}. \]
\item \textbf{Springs in parallel:} The extension is the same for each spring. The total force is the sum of individual forces.
\[ k_{\text{eq}} = k_1 + k_2 + \cdots \]
\end{itemize}
\end{propriete}

\begin{exoresolu}[Springs in parallel supporting a load]
Two springs, A and B, are suspended side by side (in parallel) from a rigid support. Spring A has $k_A = 80\ \mathrm{N/m}$, spring B has $k_B = 120\ \mathrm{N/m}$. A load of $10\ \mathrm{N}$ is attached to a light cross-bar that connects the lower ends of the springs, so that both springs stretch by the same amount. Find (a) the equivalent spring constant; (b) the extension of each spring; (c) the tension in each spring.
\tcblower
\textbf{Solution.}
(a) $k_{\text{eq}} = k_A + k_B = 80 + 120 = 200\ \mathrm{N/m}$.
(b) Common extension $x = \dfrac{F}{k_{\text{eq}}} = \dfrac{10}{200} = 0.05\ \mathrm{m} = 5\ \mathrm{cm}$.
(c) $T_A = k_A x = 80 \times 0.05 = 4\ \mathrm{N}$; $T_B = k_B x = 120 \times 0.05 = 6\ \mathrm{N}$.
Check: $T_A + T_B = 10\ \mathrm{N} = \text{load}$.
\end{exoresolu}

\courssub{Equilibrium Problems Involving Springs}

\begin{methode}[Solving equilibrium problems with springs]
\begin{enumerate}
\item Draw a clear \cle{free-body diagram} of the particle or rigid body.
\item Identify all forces: weights, tensions, thrusts, reactions, friction.
\item For each spring or elastic string, write Hooke's law: $T = kx$ or $T = \lambda x/l$. Express $x$ in terms of the geometry (e.g., $x = L - l$).
\item Apply the equilibrium conditions:
\begin{itemize}
\item For a particle: $\sum F_x = 0$, $\sum F_y = 0$ (or resolve in two convenient directions).
\item For a rigid body: $\sum F = 0$ and $\sum M = 0$ about any point.
\end{itemize}
\item Solve the resulting equations for the unknowns.
\item Check that extensions are positive (strings go slack if $x \le 0$) and that the limit of proportionality is not exceeded.
\end{enumerate}
\end{methode}

\begin{exoresolu}[Particle on a rough inclined plane held by a spring]
A particle of mass $2\ \mathrm{kg}$ rests on a rough plane inclined at $30^\circ$ to the horizontal. It is held in equilibrium by a light spring of natural length $0.5\ \mathrm{m}$ and spring constant $100\ \mathrm{N/m}$, attached to a fixed point up the plane. The spring is stretched to a length of $0.7\ \mathrm{m}$. The coefficient of friction between the particle and the plane is $\mu = 0.2$. Taking $g = 10\ \mathrm{m\,s^{-2}}$, verify that the particle is in equilibrium and find the friction force.
\tcblower
\textbf{Solution.}
Extension of spring: $x = 0.7 - 0.5 = 0.2\ \mathrm{m}$.
Spring tension (up the plane): $T = kx = 100 \times 0.2 = 20\ \mathrm{N}$.
Weight component down the plane: $mg\sin30^\circ = 2 \times 10 \times 0.5 = 10\ \mathrm{N}$.
Normal reaction: $R = mg\cos30^\circ = 20 \times \frac{\sqrt{3}}{2} = 10\sqrt{3} \approx 17.32\ \mathrm{N}$.
Maximum friction: $F_{\text{max}} = \mu R = 0.2 \times 17.32 = 3.46\ \mathrm{N}$.
For equilibrium up the plane: $T = mg\sin30^\circ + F$.
$20 = 10 + F \Rightarrow F = 10\ \mathrm{N}$.
But $F = 10\ \mathrm{N} > F_{\text{max}} = 3.46\ \mathrm{N}$.
\textbf{Conclusion:} The particle \emph{cannot} be in equilibrium with the given data; it would accelerate up the plane. The problem as stated is inconsistent — either the spring is not that stretched, or $\mu$ is larger, or the mass is larger. This illustrates the importance of checking consistency!
\end{exoresolu}

\begin{exoresolu}[Rigid body with a spring: moment equilibrium]
A uniform rod AB of length $2\ \mathrm{m}$ and weight $50\ \mathrm{N}$ is hinged at A to a vertical wall. The rod is held horizontal by a light spring of natural length $0.8\ \mathrm{m}$ and spring constant $150\ \mathrm{N/m}$, attached to the rod at B and to a fixed point C on the wall directly above A, with AC = $1.2\ \mathrm{m}$. Find the extension of the spring and the reaction at the hinge.
\tcblower
\textbf{Solution.}
Geometry: Triangle ABC is right-angled at A. $AB = 2\ \mathrm{m}$, $AC = 1.2\ \mathrm{m}$.
Length of spring $BC = \sqrt{2^2 + 1.2^2} = \sqrt{4 + 1.44} = \sqrt{5.44} \approx 2.332\ \mathrm{m}$.
Extension $x = BC - l = 2.332 - 0.8 = 1.532\ \mathrm{m}$.
Spring tension $T = kx = 150 \times 1.532 = 229.8\ \mathrm{N}$.
Angle of spring with horizontal: $\sin\theta = \frac{AC}{BC} = \frac{1.2}{2.332} \approx 0.5146$, $\theta \approx 31.0^\circ$.
Take moments about A (hinge):
Clockwise: $W \times 1 = 50 \times 1 = 50\ \mathrm{N\,m}$.
Anticlockwise: vertical component of $T$ at B: $T\sin\theta \times 2 = 229.8 \times 0.5146 \times 2 \approx 236.5\ \mathrm{N\,m}$.
Not equal! The rod would rotate. Let's find the correct extension for equilibrium.
Let $x$ be the extension. Then $BC = 0.8 + x$. By Pythagoras: $(0.8+x)^2 = 2^2 + 1.2^2 = 5.44$.
$0.8+x = \sqrt{5.44} \approx 2.332 \Rightarrow x \approx 1.532\ \mathrm{m}$ (fixed by geometry).
But then moment equilibrium gives $T\sin\theta \times 2 = 50$.
$T = kx = 150x$. $\sin\theta = \frac{1.2}{0.8+x}$.
Equation: $150x \times \frac{1.2}{0.8+x} \times 2 = 50$.
$360x = 50(0.8+x) = 40 + 50x$.
$310x = 40 \Rightarrow x = \frac{40}{310} \approx 0.129\ \mathrm{m}$.
Then $BC = 0.8 + 0.129 = 0.929\ \mathrm{m}$.
But Pythagoras requires $BC = \sqrt{5.44} \approx 2.332\ \mathrm{m}$. Contradiction.
\textbf{Interpretation:} The geometry (AC=1.2, AB=2) fixes the length BC. The spring cannot have a different length. Therefore, with the given spring constant and natural length, equilibrium is \emph{impossible} unless the hinge provides a moment (but it's a hinge, so no moment). The problem is over-constrained. In a real exam, the numbers would be consistent. This example shows how to set up the equations: geometry + Hooke's law + equilibrium.
\end{exoresolu}

\begin{exoresolu}[Consistent rigid body problem]
A uniform rod AB of length $1.5\ \mathrm{m}$ and weight $30\ \mathrm{N}$ is hinged at A. It is held at $30^\circ$ above the horizontal by a light spring attached at B. The spring has natural length $0.5\ \mathrm{m}$, spring constant $200\ \mathrm{N/m}$, and its other end is fixed to a point C on the same horizontal level as A, with AC = $1.0\ \mathrm{m}$. Find the extension of the spring and the hinge reaction.
\tcblower
\textbf{Solution.}
Coordinates: A = (0,0). Rod at $30^\circ$: B = $(1.5\cos30^\circ,\ 1.5\sin30^\circ) = (1.299,\ 0.75)$.
C = $(1.0,\ 0)$.
Length $BC = \sqrt{(1.299-1.0)^2 + (0.75-0)^2} = \sqrt{0.299^2 + 0.75^2} = \sqrt{0.0894 + 0.5625} = \sqrt{0.6519} \approx 0.807\ \mathrm{m}$.
Extension $x = 0.807 - 0.5 = 0.307\ \mathrm{m}$.
Tension $T = 200 \times 0.307 = 61.4\ \mathrm{N}$.
Angle of spring: $\tan\phi = \frac{0.75}{0.299} \approx 2.508$, $\phi \approx 68.2^\circ$ above horizontal.
Moments about A:
Weight: $30 \times (0.75\cos30^\circ) = 30 \times 0.6495 = 19.485\ \mathrm{N\,m}$ clockwise.
Spring: vertical component $T\sin\phi = 61.4 \times \sin68.2^\circ \approx 61.4 \times 0.928 = 57.0\ \mathrm{N}$.
Moment arm = horizontal distance of B from A = $1.299\ \mathrm{m}$.
Anticlockwise moment = $57.0 \times 1.299 \approx 74.0\ \mathrm{N\,m}$.
Not equal! Again, the geometry fixes BC. For equilibrium, we need $T\sin\phi \times 1.299 = 19.485$.
But $T = k(BC - 0.5)$ and $BC$, $\phi$ are fixed by geometry. So this specific setup is not in equilibrium unless the spring constant or natural length is different.
\textbf{Correct approach for exam problems:} Usually, the geometry is such that the spring is horizontal or vertical, or the angle is given, and you solve for the extension that satisfies equilibrium. Let's do a standard one.
\end{exoresolu}

\begin{exoresolu}[Standard equilibrium: horizontal rod, vertical spring]
A uniform rod AB of length $2\ \mathrm{m}$ and weight $40\ \mathrm{N}$ is hinged at A to a vertical wall. The rod is held horizontal by a vertical spring attached at B. The spring has natural length $0.4\ \mathrm{m}$ and spring constant $500\ \mathrm{N/m}$. Its upper end is fixed to a point on the wall directly above B. Find the extension of the spring and the reaction at the hinge.
\tcblower
\textbf{Solution.}
Spring is vertical at B. Let extension be $x$. Tension $T = 500x$ (upwards at B).
Moments about A: $T \times 2 = 40 \times 1$ (weight acts at centre, 1 m from A).
$1000x = 40 \Rightarrow x = 0.04\ \mathrm{m} = 4\ \mathrm{cm}$.
$T = 500 \times 0.04 = 20\ \mathrm{N}$.
Vertical equilibrium: $R_A + T = 40 \Rightarrow R_A = 20\ \mathrm{N}$ (upwards).
Horizontal equilibrium: $H_A = 0$.
Hinge reaction: $20\ \mathrm{N}$ vertically upward.
\end{exoresolu}

\courssub{Elastic Potential Energy (Brief Introduction)}

Although the syllabus focuses on forces and equilibrium, it is useful to know that work is done in stretching a spring, and this work is stored as \cle{elastic potential energy} (EPE).

\begin{propriete}[Elastic potential energy]
For a spring or string obeying Hooke's law, the elastic potential energy stored when the extension is $x$ is
\[ E = \frac{1}{2} k x^2 = \frac{1}{2} \frac{\lambda x^2}{l} = \frac{1}{2} T x. \]
This is the area under the $F$–$x$ graph (a triangle).
\end{propriete}

\begin{attention}[EPE and equilibrium]
In static equilibrium problems, we do \emph{not} usually need energy methods; force balance is sufficient. Energy methods become essential in dynamics (oscillations, projectiles with springs). For GCE Mechanics, know the formula but apply forces for equilibrium.
\end{attention}

\courssub{Synthesis: Examination-Style Problems}

\begin{exoresolu}[GCE-style: Spring balance calibration]
A spring balance has a scale of length $15\ \mathrm{cm}$ corresponding to a maximum load of $50\ \mathrm{N}$. The spring obeys Hooke's law.
\begin{enumerate}
\item[(a)] Calculate the spring constant.
\item[(b)] A load of unknown weight stretches the spring by $9\ \mathrm{cm}$. Find the weight.
\item[(c)] The same spring is used in a different balance where the scale is $20\ \mathrm{cm}$ long for the same maximum load. Explain how this is possible.
\end{enumerate}
\tcblower
\textbf{Solution.}
(a) Maximum extension $x_{\text{max}} = 0.15\ \mathrm{m}$, $F_{\text{max}} = 50\ \mathrm{N}$.
$k = \dfrac{F_{\text{max}}}{x_{\text{max}}} = \dfrac{50}{0.15} = 333.3\ \mathrm{N/m}$.
(b) $x = 0.09\ \mathrm{m}$. $F = kx = 333.3 \times 0.09 = 30\ \mathrm{N}$.
(c) The spring constant is unchanged. A longer scale for the same maximum load means the pointer moves further for the same force. This is achieved by using a \emph{longer pointer} (lever arm) or a \emph{different gearing mechanism}, not by changing the spring. The spring itself still extends only $15\ \mathrm{cm}$ for $50\ \mathrm{N}$.
\end{exoresolu}

\begin{exoresolu}[GCE-style: Two springs and a mass]
A mass of $0.5\ \mathrm{kg}$ is attached to the midpoint of a light elastic string of natural length $2\ \mathrm{m}$ and modulus $\lambda = 20\ \mathrm{N}$. The ends of the string are fixed to two points A and B on the same horizontal level, with AB = $2.4\ \mathrm{m}$. The mass hangs in equilibrium below AB. Find the vertical distance of the mass below AB.
\tcblower
\textbf{Solution.}
Let the mass be at M, distance $d$ below AB. By symmetry, AM = BM.
Half the natural length of string = $1\ \mathrm{m}$. Half the distance AB = $1.2\ \mathrm{m}$.
Each half of the string has natural length $l = 1\ \mathrm{m}$, modulus $\lambda = 20\ \mathrm{N}$ (modulus is for the whole string; for half the string, the modulus is the same because $\lambda$ is a material property, but the spring constant doubles. Better: use $T = \frac{\lambda x}{l}$ for each half with its own $l$ and $x$).
Length of half-string = $\sqrt{1.2^2 + d^2}$.
Extension of half-string $x = \sqrt{1.44 + d^2} - 1$.
Tension in each half: $T = \frac{\lambda x}{l} = \frac{20(\sqrt{1.44+d^2}-1)}{1} = 20(\sqrt{1.44+d^2}-1)$.
Vertical equilibrium: $2T \sin\theta = mg = 5\ \mathrm{N}$.
$\sin\theta = \frac{d}{\sqrt{1.44+d^2}}$.
$2 \times 20(\sqrt{1.44+d^2}-1) \times \frac{d}{\sqrt{1.44+d^2}} = 5$.
$40d \left(1 - \frac{1}{\sqrt{1.44+d^2}}\right) = 5$.
$8d \left(1 - \frac{1}{\sqrt{1.44+d^2}}\right) = 1$.
Solve for $d$. Try $d=0.2$: $\sqrt{1.44+0.04}=\sqrt{1.48}=1.2166$. $8\times0.2(1-1/1.2166)=1.6(1-0.822)=1.6\times0.178=0.285$ too small.
Try $d=0.5$: $\sqrt{1.44+0.25}=\sqrt{1.69}=1.3$. $8\times0.5(1-1/1.3)=4(1-0.769)=4\times0.231=0.924$.
Try $d=0.55$: $\sqrt{1.44+0.3025}=\sqrt{1.7425}=1.32$. $8\times0.55(1-1/1.32)=4.4(1-0.7576)=4.4\times0.2424=1.066$.
Try $d=0.53$: $\sqrt{1.44+0.2809}=\sqrt{1.7209}=1.3118$. $8\times0.53(1-1/1.3118)=4.24(1-0.7623)=4.24\times0.2377=1.008 \approx 1$.
So $d \approx 0.53\ \mathrm{m}$.
\end{exoresolu}

\courssub{Consolidation Exercises (Applications)}

\begin{exercice}[Exercise 6]
Two springs, of spring constants $k_1 = 100\ \mathrm{N/m}$ and $k_2 = 150\ \mathrm{N/m}$, are connected (a) in series, (b) in parallel. Find the equivalent spring constant in each case. If a load of $10\ \mathrm{N}$ is applied, find the total extension in each case.
\end{exercice}

\begin{exercice}[Exercise 7]
A light elastic string of natural length $0.8\ \mathrm{m}$ and modulus $40\ \mathrm{N}$ has one end fixed. A particle of mass $0.5\ \mathrm{kg}$ is attached to the other end and hangs in equilibrium. Find the extension of the string. (Take $g = 10\ \mathrm{m\,s^{-2}}$.)
\end{exercice}

\begin{exercice}[Exercise 8]
A uniform rod AB of weight $W$ and length $2a$ is freely hinged at A. It is held in equilibrium at an angle of $30^\circ$ to the horizontal by a light elastic string attached at B. The string has natural length $a$ and modulus $\lambda$. Its other end is fixed to a point C on the vertical line through A, with AC = $a$. Find $\lambda$ in terms of $W$.
\end{exercice}

\begin{exercice}[Exercise 9]
A spring balance is calibrated by hanging known masses. The pointer reads $0$ with no load. With a mass of $200\ \mathrm{g}$, the pointer reads $4\ \mathrm{cm}$. With a mass of $500\ \mathrm{g}$, the pointer reads $10\ \mathrm{cm}$. 
\begin{enumerate}
\item[(a)] Show that the spring obeys Hooke's law.
\item[(b)] Find the spring constant.
\item[(c)] What mass gives a reading of $7\ \mathrm{cm}$?
\end{enumerate}
\end{exercice}

\begin{exercice}[Exercise 10]
A block of mass $2\ \mathrm{kg}$ rests on a smooth plane inclined at $30^\circ$. It is held in equilibrium by a spring of natural length $0.5\ \mathrm{m}$ and spring constant $200\ \mathrm{N/m}$, attached to a fixed point up the plane. The spring is parallel to the plane. Find the stretched length of the spring. (Take $g = 10\ \mathrm{m\,s^{-2}}$.)
\end{exercice}

\begin{corrige}[Exercise 6]
(a) Series: $\frac{1}{k_{\text{eq}}} = \frac{1}{100} + \frac{1}{150} = \frac{5}{300} = \frac{1}{60} \Rightarrow k_{\text{eq}} = 60\ \mathrm{N/m}$. Extension $x = F/k_{\text{eq}} = 10/60 = 0.1667\ \mathrm{m}$.
(b) Parallel: $k_{\text{eq}} = 100 + 150 = 250\ \mathrm{N/m}$. Extension $x = 10/250 = 0.04\ \mathrm{m}$.
\end{corrige}

\begin{corrige}[Exercise 7]
$T = mg = 5\ \mathrm{N}$. $T = \frac{\lambda x}{l} \Rightarrow 5 = \frac{40x}{0.8} = 50x \Rightarrow x = 0.1\ \mathrm{m} = 10\ \mathrm{cm}$.
\end{corrige}

\begin{corrige}[Exercise 8]
Geometry: AC = $a$, AB = $2a$, angle at A = $30^\circ$. By cosine rule in triangle ABC:
$BC^2 = a^2 + (2a)^2 - 2(a)(2a)\cos30^\circ = a^2 + 4a^2 - 4a^2(\sqrt{3}/2) = 5a^2 - 2\sqrt{3}a^2$.
$BC = a\sqrt{5 - 2\sqrt{3}} \approx a\sqrt{5-3.464} = a\sqrt{1.536} \approx 1.239a$.
Extension $x = BC - a = a(\sqrt{5-2\sqrt{3}} - 1)$.
Tension $T = \frac{\lambda x}{a} = \lambda(\sqrt{5-2\sqrt{3}} - 1)$.
Moments about A: $T \times (2a) \times \sin(\text{angle between string and rod}) = W \times a \cos30^\circ$.
Angle between string and rod: use sine rule or coordinates. Simpler: vertical component of T at B.
Coordinates: A=(0,0), B=(2a cos30, 2a sin30)=($a\sqrt{3}$, $a$). C=(0, $a$).
Vector CB = ($a\sqrt{3}$, $0$). So string is horizontal! Length BC = $a\sqrt{3} \approx 1.732a$.
Extension $x = a\sqrt{3} - a = a(\sqrt{3}-1)$.
Tension $T = \frac{\lambda a(\sqrt{3}-1)}{a} = \lambda(\sqrt{3}-1)$.
Moments about A: $T \times (\text{vertical distance of B from A}) = T \times a$.
Weight moment: $W \times (\text{horizontal distance of centre from A}) = W \times a\cos30^\circ = W a \frac{\sqrt{3}}{2}$.
Equilibrium: $T a = W a \frac{\sqrt{3}}{2} \Rightarrow \lambda(\sqrt{3}-1) = \frac{\sqrt{3}}{2} W \Rightarrow \lambda = \frac{\sqrt{3} W}{2(\sqrt{3}-1)} = \frac{\sqrt{3}(\sqrt{3}+1)W}{2(3-1)} = \frac{(3+\sqrt{3})W}{4}$.
\end{corrige}

\begin{corrige}[Exercise 9]
(a) Load $F_1 = 0.2 \times 10 = 2\ \mathrm{N}$, $x_1 = 0.04\ \mathrm{m}$. $F_2 = 0.5 \times 10 = 5\ \mathrm{N}$, $x_2 = 0.10\ \mathrm{m}$.
Ratio $F_1/x_1 = 2/0.04 = 50$. $F_2/x_2 = 5/0.10 = 50$. Constant ratio $\Rightarrow$ Hooke's law obeyed.
(b) $k = 50\ \mathrm{N/m}$.
(c) $x = 0.07\ \mathrm{m}$. $F = kx = 50 \times 0.07 = 3.5\ \mathrm{N}$. $m = F/g = 3.5/10 = 0.35\ \mathrm{kg} = 350\ \mathrm{g}$.
\end{corrige}

\begin{corrige}[Exercise 10]
Equilibrium along plane: $T = mg\sin30^\circ = 2 \times 10 \times 0.5 = 10\ \mathrm{N}$.
$T = kx \Rightarrow 10 = 200x \Rightarrow x = 0.05\ \mathrm{m}$.
Stretched length $L = l + x = 0.5 + 0.05 = 0.55\ \mathrm{m}$.
\end{corrige}

\begin{aretenir}[Key points for Applications and Synthesis]
\begin{itemize}
\item Springs in series: $\frac{1}{k_{\text{eq}}} = \sum \frac{1}{k_i}$; same tension, extensions add.
\item Springs in parallel: $k_{\text{eq}} = \sum k_i$; same extension, tensions add.
\item For elastic strings: $T = \frac{\lambda x}{l}$, $\lambda = kl$. Strings go slack if $x \le 0$.
\item Equilibrium problems: combine Hooke's law ($T = kx$) with $\sum F = 0$ and/or $\sum M = 0$.
\item Always find extension from geometry first, then tension, then apply equilibrium.
\item Check consistency: extensions must be positive, tensions must not exceed the limit of proportionality.
\item Elastic potential energy $E = \frac{1}{2}kx^2$ is stored energy; useful for dynamics, not usually for static equilibrium.
\end{itemize}
\end{aretenir}

\coursec{Applications of Hooke's Law and Synthesis}

In the previous section we established Hooke's law, \(F = kx\), and verified it experimentally. We now extend its use to systems of springs, to elastic strings (which behave differently from springs when compressed), and to the energy stored in a stretched spring. Finally, we synthesise a universal strategy for solving any equilibrium problem — the skill that distinguishes a good candidate in the GCE Advanced Level examination.

\courssub{Springs in Series and Parallel}

When several springs support the same load, their combined effect can be replaced by a single \emph{equivalent spring} of stiffness \(k_{\text{eq}}\). The rule depends on how the springs are arranged.

\begin{popfigure}
\begin{tikzpicture}[scale=0.9, every node/.style={font=\small}]
% Series
\draw[thick] (0,2) -- (0,3);
\draw[decorate, decoration={coil, aspect=0.4, segment length=4pt, amplitude=3pt}] (0,2) -- (0,0.5) node[midway, right=2mm] {$k_1$};
\draw[decorate, decoration={coil, aspect=0.4, segment length=4pt, amplitude=3pt}] (0,0.5) -- (0,-1) node[midway, right=2mm] {$k_2$};
\draw[thick] (0,-1) -- (0,-2);
\draw[-latex, thick] (0.5,0) -- (0.5,-0.8) node[right] {$W$};
\node at (0,3.3) {Series};
% Parallel
\begin{scope}[xshift=5cm]
\draw[thick] (0,2) -- (0,3);
\draw[decorate, decoration={coil, aspect=0.4, segment length=4pt, amplitude=3pt}] (-0.6,2) -- (-0.6,0) node[midway, left=2mm] {$k_1$};
\draw[decorate, decoration={coil, aspect=0.4, segment length=4pt, amplitude=3pt}] (0.6,2) -- (0.6,0) node[midway, right=2mm] {$k_2$};
\draw[thick] (-0.6,0) -- (0.6,0);
\draw[thick] (0,0) -- (0,-1);
\draw[-latex, thick] (1.2,-0.5) -- (1.2,-1.3) node[right] {$W$};
\node at (0,3.3) {Parallel};
\end{scope}
\end{tikzpicture}
\legende{Springs in series (left) and in parallel (right) carrying the same load \(W\). Extensions add in series; loads share in parallel.}
\end{popfigure}

\textbf{Series:} The tension \(T\) is the same in both springs. If the extensions are \(x_1\) and \(x_2\), the total extension is \(x = x_1 + x_2\). By Hooke's law, \(x_1 = T/k_1\), \(x_2 = T/k_2\). Hence
\[
x = T\left(\frac{1}{k_1}+\frac{1}{k_2}\right) \quad\Longrightarrow\quad \frac{1}{k_{\text{eq}}} = \frac{1}{k_1}+\frac{1}{k_2}.
\]

\textbf{Parallel:} The extension \(x\) is the same for both springs. The total load is shared: \(W = T_1 + T_2 = k_1x + k_2x = (k_1+k_2)x\). Hence
\[
k_{\text{eq}} = k_1 + k_2.
\]

\begin{propriete}[Equivalent Stiffness of Springs]
For two springs of stiffnesses \(k_1\) and \(k_2\):
\begin{itemize}
\item \textbf{In series:} they share the same tension; the equivalent stiffness \(k_s\) satisfies \(\dfrac{1}{k_s} = \dfrac{1}{k_1} + \dfrac{1}{k_2}\). The combination is \emph{less stiff} than either spring alone.
\item \textbf{In parallel:} they share the same extension; the equivalent stiffness is \(k_p = k_1 + k_2\). The combination is \emph{stiffer} than either spring alone.
\end{itemize}
These formulae extend to any number of springs. \textbf{Proof is examinable.}
\end{propriete}

\begin{exemplebox}[Springs in Series]
A load of \(\SI{20}{N}\) is suspended from two light springs joined in series. The upper spring has stiffness \(\SI{40}{N/m}\) and the lower \(\SI{60}{N/m}\). Find the total extension and the extension of each spring.
\begin{align*}
\frac{1}{k_s} &= \frac{1}{40} + \frac{1}{60} = \frac{5}{120} = \frac{1}{24} \;\Rightarrow\; k_s = \SI{24}{N/m}.\\
x_{\text{total}} &= \frac{W}{k_s} = \frac{20}{24} = \SI{0.833}{m}.\\
x_1 &= \frac{T}{k_1} = \frac{20}{40} = \SI{0.5}{m},\qquad
x_2 = \frac{20}{60} = \SI{0.333}{m}.
\end{align*}
Check: \(0.5 + 0.333 = 0.833\) m. ✓
\end{exemplebox}

\begin{attention}[Common Mistake]
Do not confuse the formulas: series uses reciprocals, parallel uses direct sum. A quick mental check: two identical springs of stiffness \(k\) in series give \(k/2\) (softer), in parallel give \(2k\) (stiffer).
\end{attention}

\courssub{Elastic Strings and the ``Slack'' Condition}

An \emph{elastic string} obeys Hooke's law \emph{only when stretched beyond its natural length}. Unlike a spring, a string cannot push; it goes \emph{slack} (tension \(=0\)) if its length is less than or equal to its natural length \(l_0\).

\begin{definition}[Elastic String]
A light elastic string of natural length \(l_0\) and modulus of elasticity \(\lambda\) (or stiffness \(k = \lambda/l_0\)) exerts a tension
\[
T = \frac{\lambda x}{l_0} = kx \quad\text{when stretched by } x>0,
\]
and \(T = 0\) when \(x \le 0\) (i.e. length \(\le l_0\)).
\end{definition}

\begin{attention}[String Goes Slack — Critical Check]
In any problem involving an elastic string, \textbf{always verify whether the string is actually taut}. If the geometry forces the string length to be \(\le l_0\), the tension is zero and the string exerts no force. Many candidates lose marks by writing \(T = kx\) with a negative \(x\).
\end{attention}

\begin{exemplebox}[Elastic String on a Smooth Table]
A particle of mass \(\SI{0.5}{kg}\) is attached to one end of a light elastic string of natural length \(\SI{1.2}{m}\) and modulus \(\lambda = \SI{24}{N}\). The other end is fixed at point \(O\) on a smooth horizontal table. The particle is pulled to a distance \(\SI{2}{m}\) from \(O\) and released from rest. Find the initial acceleration.
\[
x = 2 - 1.2 = \SI{0.8}{m},\quad T = \frac{\lambda x}{l_0} = \frac{24 \times 0.8}{1.2} = \SI{16}{N}.
\]
By Newton's second law: \(T = ma \;\Rightarrow\; a = \frac{16}{0.5} = \SI{32}{m/s^2}\) towards \(O\).
\end{exemplebox}

\courssub{Elastic Potential Energy in a Spring}

When a spring (or elastic string) is stretched, work is done against the tension. This work is stored as \emph{elastic potential energy} (EPE). Since the tension varies linearly from \(0\) to \(kx\), the average force is \(\frac{1}{2}kx\), and the work done is \(\frac{1}{2}kx \cdot x\).

\begin{propriete}[Elastic Potential Energy]
The elastic potential energy stored in a spring (or string) of stiffness \(k\) stretched by an amount \(x\) is
\[
E_{\text{el}} = \frac{1}{2}kx^2 = \frac{\lambda x^2}{2l_0}.
\]
\textbf{This formula is an extra link to the Work–Energy chapter and is very useful in GCE problems involving conservation of energy.}
\end{propriete}

\begin{exemplebox}[Energy Approach]
A \(\SI{0.2}{kg}\) mass is attached to a spring of stiffness \(\SI{50}{N/m}\) on a smooth horizontal surface. The mass is pulled \(\SI{0.1}{m}\) from equilibrium and released. Find its maximum speed.
\[
E_{\text{el}} = \frac{1}{2}kx^2 = \frac{1}{2}\times 50 \times (0.1)^2 = \SI{0.25}{J}.
\]
At equilibrium, all EPE becomes kinetic energy: \(\frac{1}{2}mv^2 = 0.25 \;\Rightarrow\; v = \sqrt{\frac{0.5}{0.2}} = \SI{1.58}{m/s}\).
\end{exemplebox}

\courssub{Combined Equilibrium Problems: Strings and Springs Together}

Real problems often mix inextensible strings, elastic strings, and springs. The solution always follows the same pattern: draw a free-body diagram, resolve forces, apply equilibrium conditions (\(\sum F_x = 0\), \(\sum F_y = 0\), and \(\sum M = 0\) if moments are needed), and solve.

\begin{exoresolu}[Mixed System: Inextensible and Elastic Strings]
A particle of weight \(\SI{50}{N}\) is suspended by two strings: an inextensible string \(AP\) of length \(\SI{1.2}{m}\) attached to a fixed point \(A\), and a light elastic string \(BP\) of natural length \(\SI{0.6}{m}\) and modulus \(\lambda = \SI{80}{N}\) attached to a fixed point \(B\). Points \(A\) and \(B\) are \(\SI{1.5}{m}\) apart horizontally. The particle hangs in equilibrium with both strings taut. Find the tension in each string and the extension of \(BP\).

\textbf{Solution:}
\begin{enumerate}
\item \textbf{Geometry:} Let the horizontal distance from \(A\) to \(P\) be \(d\). The vertical drop is \(h\). Since \(AP = \SI{1.2}{m}\),
\[
d^2 + h^2 = 1.2^2 = 1.44. \tag{1}
\]
The horizontal distance from \(P\) to \(B\) is \(1.5 - d\). The length of \(BP\) is \(l = 0.6 + x\), where \(x\) is the extension. Hence
\[
(1.5-d)^2 + h^2 = (0.6+x)^2. \tag{2}
\]
\item \textbf{Tension in \(BP\):} By Hooke's law,
\[
T_{BP} = \frac{\lambda x}{l_0} = \frac{80x}{0.6} = \frac{400}{3}x. \tag{3}
\]
\item \textbf{Horizontal equilibrium:} The horizontal components of the tensions balance.
\[
T_{AP} \frac{d}{1.2} = T_{BP} \frac{1.5-d}{0.6+x}. \tag{4}
\]
\item \textbf{Vertical equilibrium:} The vertical components support the weight.
\[
T_{AP} \frac{h}{1.2} + T_{BP} \frac{h}{0.6+x} = 50. \tag{5}
\]
\item \textbf{Solving:} Subtract (1) from (2):
\[
(1.5-d)^2 - d^2 = (0.6+x)^2 - 1.44 \implies 2.25 - 3d = (0.6+x)^2 - 1.44.
\]
\[
3.69 - 3d = (0.6+x)^2. \tag{6}
\]
From (4) and (5), eliminate \(T_{AP}\). From (4): \(T_{AP} = T_{BP} \frac{1.2(1.5-d)}{d(0.6+x)}\). Substitute into (5):
\[
T_{BP} \frac{1.2(1.5-d)}{d(0.6+x)} \cdot \frac{h}{1.2} + T_{BP} \frac{h}{0.6+x} = 50
\]
\[
T_{BP} \frac{h}{0.6+x} \left( \frac{1.5-d}{d} + 1 \right) = 50
\]
\[
T_{BP} \frac{h}{0.6+x} \cdot \frac{1.5}{d} = 50 \implies T_{BP} h = \frac{50 d (0.6+x)}{1.5} = \frac{100}{3} d (0.6+x). \tag{7}
\]
Substitute \(T_{BP}\) from (3):
\[
\frac{400}{3} x h = \frac{100}{3} d (0.6+x) \implies 4 x h = d (0.6+x). \tag{8}
\]
Now we have three equations (1), (6), (8) for \(d\), \(h\), \(x\). From (1): \(h = \sqrt{1.44 - d^2}\). From (6): \(0.6+x = \sqrt{3.69 - 3d}\) (positive root). Substitute into (8):
\[
4 x \sqrt{1.44 - d^2} = d \sqrt{3.69 - 3d}.
\]
Also \(x = \sqrt{3.69 - 3d} - 0.6\). This system can be solved numerically or by noticing a neat solution. Try \(d = 0.96\):
\(h = \sqrt{1.44 - 0.9216} = \sqrt{0.5184} = 0.72\).
\(0.6+x = \sqrt{3.69 - 2.88} = \sqrt{0.81} = 0.9 \implies x = 0.3\).
Check (8): LHS \(= 4 \times 0.3 \times 0.72 = 0.864\); RHS \(= 0.96 \times 0.9 = 0.864\). ✓
\item \textbf{Results:}
\[
x = \SI{0.3}{m},\quad T_{BP} = \frac{400}{3} \times 0.3 = \SI{40}{N},\quad T_{AP} = \frac{T_{BP}(1.5-d)}{d} \cdot \frac{1.2}{0.6+x} = \frac{40 \times 0.54}{0.96} \times \frac{1.2}{0.9} = \SI{30}{N}.
\]
\end{enumerate}
\end{exoresolu}

\courssub{Chapter Synthesis: A Complete Strategy for Equilibrium Problems}

Every equilibrium problem in this chapter — whether it involves forces, moments, couples, Hooke's law, or a mixture — can be solved by following a disciplined, step-by-step method. Master this flowchart and you will avoid the common pitfalls that cost easy marks.

\begin{methode}[Universal Equilibrium Strategy]
\begin{enumerate}
\item \textbf{Diagram:} Draw a clear, labelled free-body diagram. Show \emph{all} forces (weights, tensions, reactions, friction, spring forces) with their directions. Indicate dimensions and angles.
\item \textbf{Choose axes:} For forces, pick perpendicular axes (usually horizontal/vertical or parallel/perpendicular to a plane). For moments, choose a convenient pivot (often where unknown forces act to eliminate them from the moment equation).
\item \textbf{Resolve forces:} Write \(\sum F_x = 0\) and \(\sum F_y = 0\) (or \(\sum F_\parallel = 0\), \(\sum F_\perp = 0\)). If the body is a rigid body, also write \(\sum M_{\text{pivot}} = 0\).
\item \textbf{Apply constitutive laws:} For springs/strings, use \(T = kx\) or \(T = \lambda x/l_0\). Remember: elastic strings go slack if \(x \le 0\).
\item \textbf{Solve the equations:} Solve simultaneously for the unknowns (tensions, reactions, extensions, angles). Keep exact fractions until the final step.
\item \textbf{Verify:} Check units (N, m, Nm). Check sense: are tensions positive? Are extensions positive? Does the geometry make sense? If a string is found to have negative extension, it is actually slack — re-solve with \(T=0\) for that string.
\end{enumerate}
\end{methode}

\begin{popfigure}
\begin{tikzpicture}[node distance=1.2cm, every node/.style={font=\small, align=center},
  proc/.style={rectangle, draw=popblue, thick, fill=popblueL, rounded corners, minimum width=3cm, minimum height=1cm},
  dec/.style={diamond, draw=poporange, thick, fill=poporange!20, aspect=2, minimum width=3cm},
  arr/.style={-latex, thick, popdark}]
\node[proc] (start) {Draw free-body diagram};
\node[proc, below=of start] (axes) {Choose axes / pivot};
\node[proc, below=of axes] (resolve) {Resolve forces: $\sum F_x=0$, $\sum F_y=0$};
\node[proc, below=of resolve] (moments) {If rigid body: $\sum M=0$};
\node[proc, below=of moments, align=center] (laws){Apply constitutive laws\\ (Hooke, friction, etc.)};
\node[proc, below=of laws] (solve) {Solve equations};
\node[dec, below=of solve, align=center] (check){Check units,\\ sense, slack?};
\node[proc, below=of check] (end) {Final answer};
\draw[arr] (start) -- (axes);
\draw[arr] (axes) -- (resolve);
\draw[arr] (resolve) -- (moments);
\draw[arr] (moments) -- (laws);
\draw[arr] (laws) -- (solve);
\draw[arr] (solve) -- (check);
\draw[arr] (check) -- node[right]{OK} (end);
\draw[arr] (check.west) -- ++(-1.5,0) |- (laws.west) node[pos=0.25, left]{Revise};
\end{tikzpicture}
\legende{Flowchart for solving any equilibrium problem. Follow the loop until all checks pass.}
\end{popfigure}

\begin{aretenir}[Key Points for the Examination]
\begin{itemize}
\item Springs in series: \(\frac{1}{k_{\text{eq}}} = \frac{1}{k_1}+\frac{1}{k_2}\) (same tension, softer).
\item Springs in parallel: \(k_{\text{eq}} = k_1+k_2\) (same extension, stiffer).
\item Elastic strings: tension only when stretched (\(x>0\)); slack when \(x\le 0\).
\item Elastic potential energy: \(E = \frac{1}{2}kx^2\) (links to conservation of energy).
\item The universal strategy (diagram → axes → resolve → laws → solve → verify) works for \emph{every} equilibrium question.
\end{itemize}
\end{aretenir}

\begin{savaistu}[Did You Know?]
The modulus of elasticity \(\lambda\) (in newtons) is the force required to double the length of a string of unit natural length. In Cameroon's rubber industry, the elasticity of latex threads is tested by measuring \(\lambda\) to ensure quality for products like elastic bands and medical tubing.
\end{savaistu}

\begin{exercice}[Practice Problems]
\begin{itemize}
\item Two springs of stiffness \(\SI{30}{N/m}\) and \(\SI{50}{N/m}\) are joined in series and a \(\SI{40}{N}\) weight is attached. Find the total extension and the extension of each spring.
\item A light elastic string of natural length \(\SI{2}{m}\) and modulus \(\SI{100}{N}\) has its ends fixed to two points \(\SI{1.5}{m}\) apart on a horizontal line. A particle of weight \(\SI{40}{N}\) is attached to the midpoint and hangs in equilibrium. Find the extension of each half of the string.
\item A block of mass \(\SI{2}{kg}\) rests on a rough plane inclined at \(30^\circ\). It is held in equilibrium by a spring of stiffness \(\SI{100}{N/m}\) fixed up the plane. The spring's natural length is \(\SI{0.5}{m}\) and its length in equilibrium is \(\SI{0.7}{m}\). Find the coefficient of friction.
\item (Synthesis) A uniform rod \(AB\) of length \(\SI{2}{m}\) and weight \(\SI{20}{N}\) is hinged at \(A\). It is held horizontal by a light elastic string attached at \(B\) and to a point \(C\) vertically above \(A\) with \(AC = \SI{1.5}{m}\). The string has natural length \(\SI{1.5}{m}\) and modulus \(\SI{25}{N}\). Find the angle the string makes with the horizontal and the reaction at the hinge.
\end{itemize}
\end{exercice}

\begin{corrige}[Exercise 1]
Series: \(1/k_s = 1/30 + 1/50 = 8/150 = 4/75 \Rightarrow k_s = 18.75\ \mathrm{N/m}\). \(x_{\text{total}} = 40/18.75 = 2.133\ \mathrm{m}\). \(x_1 = 40/30 = 1.333\ \mathrm{m}\), \(x_2 = 40/50 = 0.8\ \mathrm{m}\).
\end{corrige}

\begin{corrige}[Exercise 2]
Each half has natural length \(l_0 = 1\ \mathrm{m}\). Horizontal distance from midpoint to each support \(= 0.75\ \mathrm{m}\). Let extension be \(x\). Length of each half \(l = 1+x\). Vertical drop \(h = \sqrt{(1+x)^2 - 0.75^2} = \sqrt{x^2+2x+0.4375}\). Tension \(T = \frac{\lambda x}{l_0} = 100x\). Vertical equilibrium: \(2T\frac{h}{l} = 40 \Rightarrow 200x\frac{h}{1+x}=40 \Rightarrow 5x h = 1+x\). Substitute \(h\): \(5x\sqrt{x^2+2x+0.4375} = 1+x\). Square: \(25x^2(x^2+2x+0.4375) = (1+x)^2\). This simplifies to \(25x^4+50x^3+10.9375x^2 = x^2+2x+1\). By inspection, \(x=0.25\) is a solution: LHS \(=25(0.00390625)+50(0.015625)+10.9375(0.0625)=0.09766+0.78125+0.68359=1.5625\); RHS \(=(1.25)^2=1.5625\). Thus \(x = \SI{0.25}{m}\).
\end{corrige}

\begin{corrige}[Exercise 3]
Spring extension \(x = 0.7-0.5 = 0.2\ \mathrm{m}\). Spring force \(F_s = kx = 100\times0.2 = 20\ \mathrm{N}\) up the plane. Weight component down plane \(= mg\sin30^\circ = 2\times9.8\times0.5 = 9.8\ \mathrm{N}\). Friction \(F_r\) acts down the plane (since spring pulls up). Equilibrium up plane: \(F_s = mg\sin30^\circ + F_r \Rightarrow 20 = 9.8 + F_r \Rightarrow F_r = 10.2\ \mathrm{N}\). Normal reaction \(R = mg\cos30^\circ = 19.6\times\sqrt{3}/2 \approx 16.97\ \mathrm{N}\). \(\mu = F_r/R = 10.2/16.97 \approx 0.60\).
\end{corrige}

\begin{corrige}[Exercise 4]
Geometry: \(AC=1.5\), \(AB=2\), angle at \(A\) is \(90^\circ\) (rod horizontal, \(AC\) vertical). So \(BC = \sqrt{1.5^2+2^2} = 2.5\ \mathrm{m}\). String natural length \(=1.5\ \mathrm{m}\), so extension \(x = 1.0\ \mathrm{m}\). Tension \(T = \lambda x/l_0 = 25\times1.0/1{1.5} = \SI{16.67}{N}\). Angle of string with horizontal: \(\tan\theta = AC/AB = 1.5/2 = 0.75 \Rightarrow \theta \approx 36.87^\circ\). Moments about \(A\): \(T\sin\theta \times 2 = W \times 1 \Rightarrow 16.67\times\sin36.87^\circ\times2 = 20\times1\). \(\sin36.87^\circ = 0.6\), LHS \(= 16.67\times0.6\times2 = 20.0\), RHS \(=20\). Equilibrium satisfied. Reaction at hinge: Horizontal \(R_H = T\cos\theta = 16.67\times0.8 = \SI{13.33}{N}\) (towards \(C\)). Vertical \(R_V = W - T\sin\theta = 20 - 10 = \SI{10}{N}\) (upwards). Magnitude \(R = \sqrt{13.33^2+10^2} \approx \SI{16.67}{N}\) at angle \(\tan^{-1}(10/13.33) \approx 36.87^\circ\) above horizontal.
\end{corrige}

\newpage

\coursec{Integration activity}

\courssub{Aims of the activity}

By the end of this integrative task, you should be able to:

\begin{itemize}
\item locate the \cle{centre of gravity} of an irregular lamina experimentally by the plumb-line method and use it to choose a suspension point;
\item apply the \cle{conditions of equilibrium} of a particle (closed triangle of forces, Lami's theorem, resolution into components) to a real two-string suspension;
\item use a \cle{graphical method} (scale triangle of forces) to check an analytical result and comment on the accuracy of each approach;
\item calibrate a spring, verify \cle{Hooke's law} graphically, extract the spring constant from a gradient and graduate a scale in kilograms;
\item communicate technical results in the form of a short user manual, as a technician or engineer would.
\end{itemize}

\courssub{Situation}

\begin{experience}[The weighing station of the Mokolo market]
A women's cooperative that sells rice, groundnuts and dried fish at Mokolo market (Yaound\'e) wants a cheap, reliable weighing station. The committee has asked your Upper Sixth class to design it. The station consists of:

\begin{itemize}
\item a flat \textbf{tray} cut from a sheet of cardboard (an irregular quadrilateral $ABCD$), of mass $2\ \mathrm{kg}$, which must hang \emph{horizontally} so that goods do not slide off;
\item two light inextensible strings attached to the tray at two points and joined at a ring $O$ fixed to a beam; when loaded, the strings make angles of $40^{\circ}$ and $55^{\circ}$ with the horizontal;
\item a \textbf{spring scale}: a steel spring suspended from the beam, to be calibrated with standard masses and graduated up to $20\ \mathrm{kg}$.
\end{itemize}

The heaviest single item sold is a $15\ \mathrm{kg}$ bag of rice. Take $g = 9.81\ \mathrm{m\,s^{-2}}$. The cooperative also wants to know the greatest mass the spring can measure without being permanently damaged.

For the calibration, your group uses standard masses and a millimetre ruler. The unloaded spring has length $L_0 = 12.0\ \mathrm{cm}$. The measurements are:

\begin{center}
\begin{tabular}{|l|c|c|c|c|c|c|}
\hline
\rowcolor{popblueL} Mass $m$ (kg) & 0.5 & 1.0 & 1.5 & 2.0 & 2.5 & 3.0 \\
\hline
Length $L$ (cm) & 13.0 & 14.0 & 15.1 & 16.0 & 17.1 & 18.0 \\
\hline
\end{tabular}
\end{center}

Beyond an extension of $50\ \mathrm{cm}$ the spring does not return to its original length: this is its \cle{elastic limit}.
\end{experience}

\courssub{Tasks}

\begin{enumerate}
\item \textbf{Centre of gravity.} Describe, with a labelled diagram, the plumb-line experiment you would perform on the cardboard tray to locate its centre of gravity $G$. Explain why two pin-holes are enough in principle, and why a third is a good idea. Where should the two strings be attached (relative to $G$) for the tray to hang level?
\item \textbf{Equilibrium of the loaded tray.} The tray ($2\ \mathrm{kg}$) carries the $15\ \mathrm{kg}$ bag of rice placed so that the whole system hangs symmetrically; the junction ring $O$ is then in equilibrium under three forces: the tensions $T_1$ (at $40^{\circ}$ to the horizontal) and $T_2$ (at $55^{\circ}$), and the total weight $W$.
\begin{enumerate}
\item Compute $W$.
\item Draw the free-body diagram of the ring $O$.
\item Use \textbf{Lami's theorem} to find $T_1$ and $T_2$.
\item Check your answers by \textbf{resolving} horizontally and vertically.
\end{enumerate}
\item \textbf{Graphical verification.} Using a scale of $1\ \mathrm{cm} : 20\ \mathrm{N}$, draw the triangle of forces for the three concurrent forces at $O$. Measure $T_1$ and $T_2$ and compare with Task 2. Comment on the agreement and on the sources of error of the graphical method.
\item \textbf{Calibrating the spring.} For each mass compute the load $F = mg$ and the extension $x = L - L_0$. Plot $F$ (vertical axis) against $x$ (horizontal axis). Is Hooke's law verified over this range? Determine the spring constant $k$ from the gradient, giving its unit.
\item \textbf{Graduating the scale.} Using your value of $k$, compute the extension for loads of $5$, $10$, $15$ and $20\ \mathrm{kg}$, and explain how you would mark the scale in kilograms. Why are the graduations equally spaced?
\item \textbf{Safety limit.} The elastic limit corresponds to $x = 50\ \mathrm{cm}$. Find the greatest mass that may be hung without permanently deforming the spring. Is the requested $20\ \mathrm{kg}$ range safe?
\item \textbf{User manual.} Write a one-page manual for the trader: how to hang the tray, where to place the goods, how to read the spring scale, and the two safety warnings (maximum load, elastic limit).
\end{enumerate}

\courssub{Model answer}

\textbf{Task 1. Locating the centre of gravity.}

Punch a small hole near a corner $A$ of the tray and suspend the tray freely from a horizontal nail through the hole, so that it can turn without sticking. Hang a plumb line (a string with a small mass) from the same nail. When the tray is completely at rest, it is in equilibrium under exactly two forces: its \cle{weight}, acting vertically downwards through $G$, and the \cle{reaction} of the nail at the hole. Two forces in equilibrium must be equal, opposite and \emph{collinear}; therefore $G$ lies on the vertical line through the nail, which is exactly the line of the plumb line. Draw this line on the cardboard with a pencil.

Repeat the experiment from a second hole near corner $B$: $G$ also lies on the new plumb line. The intersection of the two lines is $G$. Two holes are enough in principle because two distinct lines meet in exactly one point. A third hole is a good idea as a \emph{check}: the three lines should be concurrent. If they are not, the experiment was careless (friction at the nail, tray not perfectly still, plumb line swinging, pencil line drawn too thickly).

For the loaded tray to hang level, the strings must be attached at points such that the vertical through the junction ring $O$ passes through the centre of gravity $G$ of the \emph{loaded} system (tray plus goods); otherwise the tray would tilt until this condition is met. In practice the strings are fixed at points placed symmetrically about $G$, and the goods are placed at the centre of the tray.

\begin{popfigure}
\begin{center}
\begin{tikzpicture}[scale=0.85]
\coordinate (N) at (0,3.2);
\coordinate (A) at (-0.4,2.6);
\draw[thick, popdark, fill=popblueL] (-0.4,2.6) -- (2.2,2.2) -- (2.6,0.2) -- (0.6,-0.8) -- (-1.8,0.4) -- cycle;
\fill[popdark] (N) circle (0.06) node[above, popdark]{nail};
\draw[thick, popink] (N) -- (A);
\draw[popmuted, dashed] (A) -- (0.1,-1.4);
\fill[popmuted] (0.1,-1.4) circle (0.09);
\draw[thick, poporange] (A) -- (0.1,-1.1);
\fill[poporange] (0.55,0.75) circle (0.07) node[right=1pt, popdark]{$G$};
\node[popblue] at (1.7,1.4) {lamina};
\node[popmuted, align=center, text width=2.2cm] at (-1.9,-1.0) {plumb line};
\end{tikzpicture}
\end{center}
\legende{Plumb-line method: $G$ lies on the vertical line drawn from the suspension hole.}
\end{popfigure}

\textbf{Task 2. Equilibrium of the loaded tray.}

(a) Total mass $M = 2 + 15 = 17\ \mathrm{kg}$, so
\[ W = Mg = 17 \times 9.81 = 166.77\ \mathrm{N} \approx 166.8\ \mathrm{N}. \]

(b) The ring $O$ is modelled as a particle in equilibrium under three concurrent forces: $T_1$ along the left string, $T_2$ along the right string, and $W$ vertically downwards (the strings are light, so each tension is transmitted unchanged).

\begin{popfigure}
\begin{center}
\begin{tikzpicture}[scale=0.9]
\coordinate (O) at (0,2.6);
\coordinate (A) at (-2.6,0);
\coordinate (B) at (2.2,0);
\draw[thick, popink] (O) -- (A) node[midway, above left, popblue]{$T_1$};
\draw[thick, popink] (O) -- (B) node[midway, above right, popblue]{$T_2$};
\draw[thick, popdark, fill=popblueL] (-2.6,0) -- (2.2,0) -- (2.6,-0.5) -- (-2.2,-0.5) -- cycle;
\fill[poporange] (0,-0.25) circle (0.07) node[below=2pt, popdark]{$G$};
\draw[->, thick, poporange] (0,-0.25) -- (0,-1.6) node[below]{$W$};
\draw[popmuted] (A) ++(0.75,0) arc (0:40:0.75) node[midway, right, popmuted]{$40^{\circ}$};
\draw[popmuted] (B) ++(-0.85,0) arc (180:125:0.85) node[midway, left, popmuted]{$55^{\circ}$};
\fill[popdark] (O) circle (0.08) node[above, popdark]{$O$};
\end{tikzpicture}
\end{center}
\legende{Free-body diagram of the loaded tray suspended by two strings.}
\end{popfigure}

(c) \textbf{Lami's theorem.} The angle between the two strings is $180^{\circ} - 40^{\circ} - 55^{\circ} = 85^{\circ}$; the angle between $T_1$ and the downward vertical is $90^{\circ} + 40^{\circ} = 130^{\circ}$; the angle between $T_2$ and the downward vertical is $90^{\circ} + 55^{\circ} = 145^{\circ}$. Lami's theorem states that each force is proportional to the sine of the angle between the \emph{other two}:
\[ \frac{T_1}{\sin 145^{\circ}} = \frac{T_2}{\sin 130^{\circ}} = \frac{W}{\sin 85^{\circ}}. \]
Hence
\[ T_1 = \frac{166.77 \times \sin 145^{\circ}}{\sin 85^{\circ}} = \frac{166.77 \times 0.5736}{0.9962} \approx 96.0\ \mathrm{N}, \]
\[ T_2 = \frac{166.77 \times \sin 130^{\circ}}{\sin 85^{\circ}} = \frac{166.77 \times 0.7660}{0.9962} \approx 128.2\ \mathrm{N}. \]

(d) \textbf{Resolution check.} Horizontally, the components must cancel:
\[ T_2 \cos 55^{\circ} = 128.2 \times 0.5736 \approx 73.5\ \mathrm{N}, \qquad T_1 \cos 40^{\circ} = 96.0 \times 0.7660 \approx 73.5\ \mathrm{N}. \]
They balance. Vertically, the upward components must support the weight:
\[ T_1 \sin 40^{\circ} + T_2 \sin 55^{\circ} = 96.0 \times 0.6428 + 128.2 \times 0.8192 \approx 61.7 + 105.0 = 166.7\ \mathrm{N} \approx W. \]
Both conditions of equilibrium are satisfied, which confirms the Lami calculation.

\begin{attention}[Choosing the angles in Lami's theorem]
The angle used with each force is the angle \emph{between the other two forces}, measured at the point of concurrence. A very common error is to use $40^{\circ}$ and $55^{\circ}$ directly instead of $130^{\circ}$, $145^{\circ}$ and $85^{\circ}$. Always draw the three forces meeting at one point and read off the three angles between them.
\end{attention}

\textbf{Task 3. Graphical verification.}

Draw $W = 166.8\ \mathrm{N}$ vertically downward as a segment of length $\dfrac{166.8}{20} = 8.34\ \mathrm{cm}$. From its upper end draw a line at $40^{\circ}$ to the horizontal (direction of $T_1$) and from its lower end a line at $55^{\circ}$ to the horizontal (direction of $T_2$); their intersection closes the triangle of forces. Measuring the two unknown sides gives about $4.8\ \mathrm{cm}$ and $6.4\ \mathrm{cm}$, i.e.\ $T_1 \approx 96\ \mathrm{N}$ and $T_2 \approx 128\ \mathrm{N}$, in excellent agreement with Task 2.

\begin{popfigure}
\begin{center}
\begin{tikzpicture}[scale=0.045]
\draw[->, very thick, poporange] (0,0) -- (0,-166.8) node[midway, left, popdark]{$W$};
\draw[->, very thick, popblue] (0,0) -- (-73.5,61.7) node[midway, above left, popdark]{$T_1$};
\draw[->, very thick, popblue] (0,-166.8) -- (-73.5,-105.1) node[midway, below left, popdark]{$T_2$};
\draw[popmuted, dashed] (0,0) -- (-73.5,0);
\draw[popmuted] (-30,0) arc (180:140:30) node[midway, left, popmuted]{$40^{\circ}$};
\end{tikzpicture}
\end{center}
\legende{Closed triangle of forces at $O$ (not to scale): the three vectors add to zero.}
\end{popfigure}

Small discrepancies come from protractor and ruler reading errors (about $\pm 0.5^{\circ}$ and $\pm 0.5\ \mathrm{mm}$). The analytical method is more precise; the graphical method is faster and gives an immediate visual check that the triangle closes.

\textbf{Task 4. Calibrating the spring.}

With $L_0 = 12.0\ \mathrm{cm}$, $F = mg$ and $x = L - L_0$:

\begin{center}
\begin{tabular}{|l|c|c|c|c|c|c|}
\hline
\rowcolor{popblueL} $F = mg$ (N) & 4.91 & 9.81 & 14.72 & 19.62 & 24.53 & 29.43 \\
\hline
$x$ (cm) & 1.0 & 2.0 & 3.1 & 4.0 & 5.1 & 6.0 \\
\hline
\end{tabular}
\end{center}

\begin{popfigure}
\begin{center}
\begin{tikzpicture}
\begin{axis}[width=9cm, height=6.5cm, xlabel={$x$ (cm)}, ylabel={$F$ (N)}, xmin=0, xmax=7, ymin=0, ymax=32, grid=both, grid style={popline}, legend pos=north west]
\addplot[only marks, mark=*, poporange, mark size=2pt] coordinates {(1.0,4.91) (2.0,9.81) (3.1,14.72) (4.0,19.62) (5.1,24.53) (6.0,29.43)};
\addlegendentry{data}
\addplot[domain=0:6.5, thick, popblue] {4.9*x};
\addlegendentry{best-fit line}
\end{axis}
\end{tikzpicture}
\end{center}
\legende{Load against extension: the points lie on a straight line through the origin.}
\end{popfigure}

The points lie on a straight line through the origin, so $F$ is proportional to $x$: \cle{Hooke's law} $F = kx$ is verified over this range. Using two well-separated points on the best-fit line, the gradient is
\[ k = \frac{\Delta F}{\Delta x} = \frac{29.43 - 4.91}{(6.0 - 1.0)\times 10^{-2}} = \frac{24.52}{0.05} \approx 490\ \mathrm{N\,m^{-1}}. \]
The unit is the newton per metre, since $k = F/x$.

\begin{attention}[Units in the gradient]
Convert the extension to metres before dividing: using centimetres would give $k = 4.9\ \mathrm{N\,cm^{-1}}$, which is the same constant but not the SI value expected at the GCE. State $k \approx 490\ \mathrm{N\,m^{-1}}$.
\end{attention}

\textbf{Task 5. Graduating the scale.}

With $k = 490\ \mathrm{N\,m^{-1}}$, the extension for a mass $m$ is $x = \dfrac{mg}{k}$:

\begin{center}
\begin{tabular}{|l|c|c|c|c|}
\hline
\rowcolor{popblueL} Mass (kg) & 5 & 10 & 15 & 20 \\
\hline
$F$ (N) & 49.05 & 98.1 & 147.2 & 196.2 \\
\hline
$x$ (cm) & 10.0 & 20.0 & 30.0 & 40.0 \\
\hline
\end{tabular}
\end{center}

Fix a strip of card behind the spring, mark the pointer positions corresponding to these extensions and label them $5, 10, 15, 20\ \mathrm{kg}$. Because $x$ is proportional to $m$ (Hooke's law), equal steps of mass produce equal steps of extension: the graduations are \emph{equally spaced}, here $2\ \mathrm{cm}$ per kilogram, so intermediate marks (every kilogram, or every $500\ \mathrm{g}$) can be added by simple division of the intervals.

\textbf{Task 6. Safety limit.}

At the elastic limit, $x_{\max} = 50\ \mathrm{cm} = 0.50\ \mathrm{m}$, so
\[ F_{\max} = kx_{\max} = 490 \times 0.50 = 245\ \mathrm{N}, \qquad m_{\max} = \frac{F_{\max}}{g} = \frac{245}{9.81} \approx 25.0\ \mathrm{kg}. \]
The requested $20\ \mathrm{kg}$ range is safe, with a margin of about $5\ \mathrm{kg}$ before the elastic limit is reached. Beyond $25\ \mathrm{kg}$ the spring would be permanently stretched and every subsequent reading would be wrong.

\textbf{Task 7. User manual (model content).}

\begin{methode}[Weighing station: instructions for the trader]
\begin{enumerate}
\item \textbf{Hanging the tray.} Suspend the tray by its two strings so that the vertical line through the ring $O$ passes through the centre mark $G$ painted on the tray; the tray then hangs level.
\item \textbf{Placing the goods.} Always place the goods at the centre of the tray, directly above $G$, so that the tray stays horizontal and the strings share the load as designed.
\item \textbf{Reading the scale.} Hang the item on the spring hook, wait until the pointer is still, and read it at eye level to avoid parallax error.
\item \textbf{Safety warning 1.} Never put more than $20\ \mathrm{kg}$ on the scale: the graduations stop there.
\item \textbf{Safety warning 2.} Never load the spring beyond $25\ \mathrm{kg}$: the spring would pass its elastic limit, stay permanently stretched, and all readings would become false.
\end{enumerate}
\end{methode}

\begin{aretenir}[What this activity mobilises]
Locating a centre of gravity experimentally by the plumb-line method; the conditions of equilibrium of a particle (zero resultant, closed triangle of forces); Lami's theorem and resolution into components as two equivalent analytical tools; the scale triangle of forces as a graphical check; Hooke's law $F = kx$, its graphical verification and the reading of a gradient to obtain $k$; and the practical meaning of the elastic limit as the greatest load an instrument can bear without permanent damage.
\end{aretenir}

\newpage

\coursec{Methods and worked examples}

\coursec{Equilibrium, Forces, Moments and Couples}

\courssub{Introduction and Learning Outcomes}

\begin{aretenir}[What you will master in this chapter]
\begin{itemize}
\item The concept of \cle{equilibrium} for particles and rigid bodies: translational and rotational equilibrium.
\item The \cle{centre of gravity} (or centre of mass) of a body and how to locate it experimentally and by calculation.
\item Forces as vectors: addition by the parallelogram/triangle rule, resolution into components, and the use of free-body diagrams.
\item The \cle{resultant} and \cle{equilibrant} of a system of coplanar forces; graphical and analytical methods.
\item The \cle{moment of a force} and the \cle{couple}; the principle of moments and conditions for equilibrium of a rigid body.
\item \cle{Hooke's law} for elastic springs and strings: experimental verification, spring constant, energy stored, and combinations of springs (series and parallel).
\item Solving complete equilibrium problems involving suspended bodies, bodies on inclined planes, and systems of springs.
\end{itemize}
\end{aretenir}

\begin{savaistu}[Why this matters]
From the roof trusses of a house in Douala to the suspension cables of the Wouri Bridge, every structure that stands still obeys the same two conditions: the vector sum of the forces is zero, and the sum of the moments about any point is zero. Mastering these ideas lets you analyse and design safe, stable systems — a core skill for any engineer or physicist.
\end{savaistu}

\coursec{Equilibrium and Centre of Gravity}

\begin{definition}[Equilibrium of a particle and a rigid body]
A \textbf{particle} is in equilibrium if the vector sum of all forces acting on it is zero:
\[ \sum \vec{F} = \vec{0}. \]
A \textbf{rigid body} is in \textbf{complete equilibrium} if two conditions hold simultaneously:
\begin{enumerate}
\item \textbf{Translational equilibrium:} $\sum \vec{F} = \vec{0}$ (no net force, centre of mass does not accelerate).
\item \textbf{Rotational equilibrium:} $\sum \vec{\mathcal{M}} = \vec{0}$ about \emph{any} point (no net moment, body does not rotate).
\end{enumerate}
For coplanar forces these give three independent scalar equations:
\[ \sum F_x = 0,\qquad \sum F_y = 0,\qquad \sum \mathcal{M}_O = 0 \quad (\text{about any point } O). \]
\end{definition}

\begin{propriete}[Centre of gravity (centre of mass)]
The \cle{centre of gravity} $G$ of a body is the point through which the resultant of the gravitational forces on all its particles (the \textbf{weight} $W=mg$) acts, for any orientation of the body. For a uniform body in a uniform gravitational field, $G$ coincides with the \textbf{geometric centre} (centroid).
\begin{itemize}
\item Uniform rod: midpoint.
\item Uniform rectangular lamina: intersection of diagonals.
\item Uniform triangular lamina: centroid (intersection of medians), at $\frac{1}{3}$ of the median from the base.
\item Uniform solid hemisphere: $3R/8$ from the flat face on the axis of symmetry.
\end{itemize}
\end{propriete}

\begin{propriete}[Equilibrium of a suspended body]
A body suspended from a fixed point $O$ by a string or a smooth hinge comes to rest with its centre of gravity $G$ \textbf{vertically below} $O$. The line of action of the weight passes through $O$, so the moment of the weight about $O$ is zero. If the body is slightly displaced, the weight produces a \emph{restoring moment} (stable equilibrium).
\end{propriete}

\begin{methode}[Testing the equilibrium of a suspended body]
\begin{enumerate}
\item Identify \textbf{all} forces: weight $W=mg$ acting vertically downward through $G$, and the tension(s) or reaction(s) at the support(s).
\item Apply the first condition: $\sum \vec{F} = \vec{0}$.
\item For a body hanging from \textbf{one point}: at equilibrium $G$ is vertically below the suspension point.
\item For a body supported by \textbf{two strings}: the three forces are concurrent (lines of action meet at one point).
\item If the body can rotate, apply the second condition: $\sum \mathcal{M} = 0$ about any convenient point (often the suspension point).
\end{enumerate}
\textbf{Reflex:} always sketch the body slightly displaced; the weight must produce a restoring moment for stable equilibrium.
\end{methode}

\begin{exoresolu}[A hanging signboard in Buea]
A uniform rectangular signboard $ABCD$ of mass $12\ \mathrm{kg}$ hangs from a nail at the midpoint $O$ of its top edge $AB$. The board is $1.2\ \mathrm{m}$ wide and $0.8\ \mathrm{m}$ high. (a) Describe the position of the board at equilibrium. (b) A horizontal wind force of $20\ \mathrm{N}$ acts at the centre of the board. Through what angle does the board turn? Take $g=\SI{9.81}{m.s^{-2}}$.
\tcblower
\textbf{Solution.}
\begin{enumerate}
\item \textbf{Forces.} Weight $W=mg=12\times 9.81=117.72\ \mathrm{N}$ acts at $G$, the geometric centre of the uniform board; the reaction of the nail acts at $O$.
\item \textbf{(a)} With no wind, the only two forces are $W$ and the reaction at $O$. For equilibrium they must be collinear, so $G$ hangs vertically below $O$: the board hangs straight, with $AB$ horizontal.
\item \textbf{(b)} With the wind force $F=20\ \mathrm{N}$ horizontal through $G$, take moments about $O$ when the board has turned through an angle $\theta$. The perpendicular distance from $O$ to the line of action of $F$ is $OG\cos\theta$; the perpendicular distance to the line of action of $W$ is $OG\sin\theta$, where $OG=0.4\ \mathrm{m}$ (half the height).
\item Moment balance about $O$ (clockwise positive):
\[ F\cdot OG\cos\theta = W\cdot OG\sin\theta \;\Longrightarrow\; \tan\theta=\frac{F}{W}=\frac{20}{117.72}=0.1699. \]
\item Hence $\theta=\tan^{-1}(0.1699)\approx 9.65^{\circ}$.
\end{enumerate}
\textbf{Answer:} (a) $G$ vertically below $O$; (b) the board tilts through about $9.7^{\circ}$.
\end{exoresolu}

\begin{experience}[Locating the centre of gravity of an irregular lamina (plumb-line method)]
\begin{enumerate}
\item Cut the lamina (cardboard, thin plywood) to shape. Make three small holes $A$, $B$, $C$ near the edge.
\item Suspend the lamina freely from a pin through hole $A$. Hang a \cle{plumb line} (string with a small bob) from the same pin.
\item When oscillations cease, mark the vertical line of the plumb line on the lamina. The centre of gravity $G$ lies on this line.
\item Repeat from holes $B$ and $C$, drawing the plumb line each time.
\item The three lines should intersect at a single point — that point is $G$.
\item \textbf{Verification:} balance the lamina horizontally on a sharp point (e.g. a needle) placed under $G$; it should remain horizontal.
\end{enumerate}
\textbf{Precautions:} the lamina must swing freely (no friction at the pin), the plumb line must be still before marking, holes must be small so as not to shift $G$ appreciably, and the lamina must be uniform in thickness and density.
\end{experience}

\begin{exoresolu}[Locating $G$ of a triangular lamina]
A student cuts a triangular lamina with vertices $A(0,0)$, $B(12,0)$, $C(0,9)$ (coordinates in cm) from uniform cardboard. (a) State how she would locate $G$ experimentally. (b) Compute the exact position of $G$ and compare. (c) The lamina is suspended from $A$; find the angle that $AB$ makes with the vertical at equilibrium.
\tcblower
\textbf{Solution.}
\begin{enumerate}
\item \textbf{(a)} She suspends the lamina from a pin through a hole near $A$, marks the plumb line; repeats from a hole near $B$; the intersection of the two lines gives $G$. A third suspension from $C$ serves as a check.
\item \textbf{(b)} For a uniform triangular lamina, $G$ is at the centroid:
\[ x_G=\frac{0+12+0}{3}=4\ \mathrm{cm},\qquad y_G=\frac{0+0+9}{3}=3\ \mathrm{cm}. \]
So $G=(4,3)\ \mathrm{cm}$, which should coincide with the experimental intersection.
\item \textbf{(c)} Suspended from $A$, $G$ hangs vertically below $A$. The line $AG$ has slope $\dfrac{3}{4}$ to the horizontal, so it makes $\tan^{-1}(3/4)=36.87^{\circ}$ with the horizontal, hence $90^{\circ}-36.87^{\circ}=53.13^{\circ}$ with the vertical. Since $AB$ is along the $x$-axis (horizontal), at equilibrium $AG$ is vertical, so $AB$ makes $53.13^{\circ}$ with the vertical (or $36.87^{\circ}$ with the horizontal).
\end{enumerate}
\textbf{Answer:} $G=(4,3)\ \mathrm{cm}$; $AB$ makes $53.1^{\circ}$ with the vertical.
\end{exoresolu}

\coursec{Force as a Vector; Addition and Resolution}

\begin{definition}[Force as a vector]
A \cle{force} is a vector quantity: it has magnitude (in newtons, N), direction, and sense. It is represented by an arrow whose length is proportional to the magnitude (according to a chosen scale) and whose direction and sense match the force. Forces obey the \textbf{parallelogram law of addition}.
\end{definition}

\begin{propriete}[Addition of forces — resultant]
The \cle{resultant} $\vec{R}$ of two forces $\vec{F}_1$ and $\vec{F}_2$ is the single force that has the same effect as the two acting together. It is the diagonal of the parallelogram constructed on $\vec{F}_1$ and $\vec{F}_2$:
\[ \vec{R} = \vec{F}_1 + \vec{F}_2. \]
Magnitude: $R = \sqrt{F_1^2 + F_2^2 + 2F_1F_2\cos\theta}$, where $\theta$ is the angle between $\vec{F}_1$ and $\vec{F}_2$.
Direction: $\tan\alpha = \dfrac{F_2\sin\theta}{F_1+F_2\cos\theta}$ (angle with $\vec{F}_1$).
\end{propriete}

\begin{propriete}[Resolution of a force into components]
Any force $\vec{F}$ can be replaced by two (or more) component forces whose vector sum is $\vec{F}$. For perpendicular axes $Ox$ and $Oy$, if $\vec{F}$ makes an angle $\theta$ with $Ox$:
\[ F_x = F\cos\theta,\qquad F_y = F\sin\theta. \]
\textbf{Memory aid:} the component \emph{adjacent} to the angle uses $\cos$; the component \emph{opposite} uses $\sin$.
\end{propriete}

\begin{methode}[Resolving a force — inclined plane reflex]
\begin{enumerate}
\item Draw the force $\vec{F}$ and choose two perpendicular axes (often \textbf{along} and \textbf{perpendicular to} a plane).
\item If $F$ makes an angle $\theta$ with an axis, the component \textbf{along} that axis is $F\cos\theta$ and the component \textbf{perpendicular} to it is $F\sin\theta$.
\item Replace $\vec{F}$ by its components and work with each axis separately.
\item \textbf{Inclined plane:} for a body of weight $W$ on a plane inclined at $\alpha$ to the horizontal, resolve the weight:
\[ W_{\parallel}=W\sin\alpha \ (\text{down the slope}),\qquad W_{\perp}=W\cos\alpha \ (\text{into the plane}). \]
\item Equilibrium: sum of components along each axis equals zero. On a smooth plane, the normal reaction is $R=W\cos\alpha$.
\end{enumerate}
\end{methode}

\begin{exoresolu}[Crate on a smooth ramp]
A crate of mass $50\ \mathrm{kg}$ rests on a smooth plank inclined at $25^{\circ}$ to the horizontal, held at rest by a rope parallel to the plank. Find (a) the tension in the rope, (b) the normal reaction of the plank. Take $g=\SI{9.81}{m.s^{-2}}$.
\tcblower
\textbf{Solution.}
\begin{enumerate}
\item \textbf{Forces:} weight $W=mg=50\times9.81=490.5\ \mathrm{N}$ vertically down; normal reaction $R$ perpendicular to the plank; tension $T$ up the plank.
\item \textbf{Resolve $W$:} along the plane $W\sin25^{\circ}$ (down-slope); perpendicular $W\cos25^{\circ}$ (into the plane).
\item \textbf{(a) Equilibrium along the plane (up positive):}
\[ T - W\sin25^{\circ} = 0 \;\Longrightarrow\; T = 490.5\times\sin25^{\circ}=490.5\times0.4226=207.3\ \mathrm{N}. \]
\item \textbf{(b) Equilibrium perpendicular to the plane (outward positive):}
\[ R - W\cos25^{\circ} = 0 \;\Longrightarrow\; R = 490.5\times\cos25^{\circ}=490.5\times0.9063=444.5\ \mathrm{N}. \]
\end{enumerate}
\textbf{Answer:} $T\approx 207\ \mathrm{N}$, $R\approx 445\ \mathrm{N}$.
\begin{center}
\begin{tikzpicture}[scale=0.9]
\draw[thick] (0,0) -- (5,0);
\draw[thick] (0,0) -- (5,2.33);
\draw[fill=popblueL] (2.2,1.03) rectangle (3.2,1.83);
\draw[->,thick,popink] (2.7,1.43) -- (2.7,0.03) node[below left]{$W$};
\draw[->,thick,poporange] (3.2,1.43) -- (4.4,1.99) node[right]{$T$};
\draw[->,thick,poppurple] (2.5,1.83) -- (1.77,2.72) node[above]{$R$};
\draw (1.2,0) arc (0:25:1.2) node[right]{$\alpha$};
\draw[dashed] (2.7,1.43) -- (2.7,2.5);
\draw[dashed] (2.7,1.43) -- (4.0,1.43);
\end{tikzpicture}
\end{center}
\end{exoresolu}

\coursec{Types of Forces and Free-Body Diagrams}

\begin{definition}[Common forces in mechanics]
\begin{itemize}
\item \textbf{Weight} $W=mg$: acts vertically downward through the centre of gravity.
\item \textbf{Normal reaction} $R$ (or $N$): perpendicular to the contact surface, pushes away from the surface.
\item \textbf{Friction} $F_f$: parallel to the contact surface, opposes relative motion (or tendency). Maximum static friction $F_{\max}=\mu_s R$; kinetic friction $F_k=\mu_k R$.
\item \textbf{Tension} $T$: in a taut string/rope, pulls along the string away from the body.
\item \textbf{Thrust/Compression}: in a rod/strut, pushes along the rod toward the body.
\item \textbf{Applied forces}: pushes, pulls, wind, etc., specified in the problem.
\end{itemize}
\end{definition}

\begin{methode}[Drawing a free-body diagram (FBD)]
\begin{enumerate}
\item \textbf{Isolate} the body of interest (imagine it cut free from its surroundings).
\item \textbf{Draw} the body as a simple shape (dot, box, outline).
\item \textbf{Show all external forces} acting \emph{on} the body, with arrows starting \emph{on} the body. Label each force clearly ($W$, $R$, $T_1$, $F$, etc.).
\item \textbf{Do not} show forces exerted \emph{by} the body on other objects, nor internal forces.
\item \textbf{Indicate} known angles and dimensions; choose coordinate axes.
\end{enumerate}
\textbf{Golden rule:} if you miss a force, your equations will be wrong. Check: gravity? contact? strings? applied?
\end{methode}

\begin{exoresolu}[Free-body diagrams practice]
Draw the FBD for: (a) a block at rest on a rough inclined plane; (b) a picture frame hanging from two strings attached to its top corners; (c) a uniform ladder leaning against a smooth wall on rough ground.
\tcblower
\textbf{Solution.}
\begin{enumerate}
\item \textbf{(a) Block on rough incline:} weight $W$ vertically down; normal reaction $R$ perpendicular to plane; friction $F$ up the plane (opposing tendency to slide down).
\item \textbf{(b) Picture frame:} weight $W$ down at centre; tensions $T_1$, $T_2$ along strings, pulling upward and inward at the corners.
\item \textbf{(c) Ladder:} weight $W$ down at midpoint; at top (smooth wall) normal reaction $R_W$ horizontal; at bottom (rough ground) normal reaction $R_G$ vertical and friction $F_G$ horizontal (toward wall to prevent slipping out).
\end{enumerate}
\begin{center}
\begin{tikzpicture}[scale=0.7]
% (a)
\begin{scope}[xshift=0cm]
\draw[thick] (0,0) -- (3,0) -- (3,1.5) -- (0,1.5) -- cycle;
\draw[->,thick,popink] (1.5,0.75) -- (1.5,-0.5) node[below]{$W$};
\draw[->,thick,poppurple] (1.5,1.5) -- (1.5,2.3) node[above]{$R$};
\draw[->,thick,poporange] (0.5,0.75) -- (1.5,0.75) node[right]{$F$};
\node at (1.5,-1) {(a)};
\end{scope}
% (b)
\begin{scope}[xshift=5cm]
\draw[thick] (0,0) -- (3,0) -- (3,1.5) -- (0,1.5) -- cycle;
\draw[->,thick,popink] (1.5,0.75) -- (1.5,-0.5) node[below]{$W$};
\draw[->,thick,poporange] (0,1.5) -- (-1,2.5) node[left]{$T_1$};
\draw[->,thick,poporange] (3,1.5) -- (4,2.5) node[right]{$T_2$};
\node at (1.5,-1) {(b)};
\end{scope}
% (c)
\begin{scope}[xshift=10cm]
\draw[thick] (0,0) -- (0.3,3);
\draw[thick] (0,0) -- (3,0);
\draw[->,thick,popink] (0.15,1.5) -- (0.15,0.5) node[below]{$W$};
\draw[->,thick,poppurple] (0.3,3) -- (1.3,3) node[right]{$R_W$};
\draw[->,thick,poppurple] (0,0) -- (0,-0.8) node[below]{$R_G$};
\draw[->,thick,poporange] (0,0) -- (0.8,0) node[right]{$F_G$};
\node at (0.15,-1.3) {(c)};
\end{scope}
\end{tikzpicture}
\end{center}
\end{exoresolu}

\coursec{Resultant and Equilibrant; Graphical Methods}

\begin{definition}[Resultant and equilibrant]
The \cle{resultant} $\vec{R}$ of a system of forces is the single force that produces the same effect as the whole system. The \cle{equilibrant} $\vec{E}$ is the single force that, when added to the system, brings it into equilibrium:
\[ \vec{E} = -\vec{R}. \]
The equilibrant has the \textbf{same magnitude} as the resultant but the \textbf{opposite direction}.
\end{definition}

\begin{methode}[Graphical determination of resultant — parallelogram and polygon rules]
\textbf{Two forces (parallelogram rule):}
\begin{enumerate}
\item Choose a suitable \cle{scale} (e.g. $1\ \mathrm{cm}\equiv 10\ \mathrm{N}$).
\item From a point $O$, draw the two forces to scale with the correct angle between them.
\item Complete the parallelogram; the diagonal from $O$ is the \textbf{resultant}. Measure its length and convert with the scale; measure its direction with a protractor.
\end{enumerate}
\textbf{Several forces (polygon rule):}
\begin{enumerate}
\item Draw the forces \emph{head-to-tail} in any order, each to scale and in the correct direction.
\item The resultant is the vector from the \textbf{start of the first} to the \textbf{end of the last}.
\item If the polygon \textbf{closes}, the resultant is zero: the forces are in equilibrium (the closing side, drawn in the opposite sense, is the \emph{equilibrant}).
\end{enumerate}
\textbf{Reflex:} the equilibrant is equal in magnitude to the resultant but opposite in direction.
\end{methode}

\begin{exoresolu}[Two tug-of-war teams]
Two teams pull a tyre at a school sports day in Bamenda with forces $F_1=300\ \mathrm{N}$ due east and $F_2=400\ \mathrm{N}$ at $60^{\circ}$ north of east. (a) Find the resultant graphically. (b) Confirm the result by calculation. (c) State the equilibrant.
\tcblower
\textbf{Solution.}
\begin{enumerate}
\item \textbf{(a) Graphical:} scale $1\ \mathrm{cm}\equiv 50\ \mathrm{N}$. Draw $\vec{F_1}$ as $6\ \mathrm{cm}$ east; from its head draw $\vec{F_2}$ as $8\ \mathrm{cm}$ at $60^{\circ}$ to the east direction. The closing vector from the start measures about $12.2\ \mathrm{cm}$, i.e. $R\approx 610\ \mathrm{N}$, at about $35^{\circ}$ north of east.
\item \textbf{(b) Analytical check.} Components (east = $x$, north = $y$):
\[ R_x = 300+400\cos60^{\circ}=300+200=500\ \mathrm{N}, \]
\[ R_y = 400\sin60^{\circ}=346.4\ \mathrm{N}. \]
\[ R=\sqrt{500^2+346.4^2}=\sqrt{250000+120000}=\sqrt{370000}=608.3\ \mathrm{N}. \]
\[ \tan\theta=\frac{346.4}{500}=0.6928 \;\Longrightarrow\; \theta=34.7^{\circ}. \]
The graphical and analytical results agree (small differences come from drawing accuracy).
\item \textbf{(c) Equilibrant:} $E=608.3\ \mathrm{N}$ directed at $34.7^{\circ}$ \emph{south of west} (opposite to $\vec R$).
\end{enumerate}
\textbf{Answer:} $R\approx 608\ \mathrm{N}$ at $34.7^{\circ}$ north of east; equilibrant equal and opposite.
\end{exoresolu}

\coursec{Analytical Methods and Equilibrium Problems}

\begin{methode}[Component method for concurrent coplanar forces]
\begin{enumerate}
\item Draw a clear \cle{free-body diagram}: isolate the body and show \emph{only} the external forces on it.
\item Choose axes $Ox$, $Oy$ (horizontal/vertical, or along/perpendicular to a plane).
\item Resolve every force: $F_x=F\cos\theta$, $F_y=F\sin\theta$, with careful signs.
\item Sum components: $R_x=\sum F_x$, $R_y=\sum F_y$.
\item Resultant: $R=\sqrt{R_x^2+R_y^2}$, direction $\tan\theta=\left|\dfrac{R_y}{R_x}\right|$ (fix the quadrant from the signs of $R_x,R_y$).
\item \textbf{Equilibrium:} $R_x=0$ \emph{and} $R_y=0$ — two scalar equations, enough to find two unknowns.
\end{enumerate}
\end{methode}

\begin{exoresolu}[Picture frame on two strings]
A picture frame of weight $60\ \mathrm{N}$ hangs from a nail by two strings attached to its top corners. The left string makes $30^{\circ}$ with the vertical and the right string $45^{\circ}$ with the vertical. Find the tensions $T_1$ and $T_2$ in the strings.
\tcblower
\textbf{Solution.}
\begin{enumerate}
\item \textbf{Free-body diagram:} three concurrent forces on the frame: $T_1$ at $30^{\circ}$ left of vertical, $T_2$ at $45^{\circ}$ right of vertical, $W=60\ \mathrm{N}$ downward.
\item \textbf{Horizontal equilibrium} ($\sum F_x=0$), right positive:
\[ T_2\sin45^{\circ}-T_1\sin30^{\circ}=0 \;\Longrightarrow\; 0.7071\,T_2=0.5\,T_1 \;\Longrightarrow\; T_1=1.4142\,T_2. \]
\item \textbf{Vertical equilibrium} ($\sum F_y=0$), up positive:
\[ T_1\cos30^{\circ}+T_2\cos45^{\circ}-60=0 \;\Longrightarrow\; 0.8660\,T_1+0.7071\,T_2=60. \]
\item \textbf{Substitute} $T_1=1.4142\,T_2$:
\[ 0.8660(1.4142\,T_2)+0.7071\,T_2=60 \;\Longrightarrow\; 1.2247\,T_2+0.7071\,T_2=60, \]
\[ 1.9319\,T_2=60 \;\Longrightarrow\; T_2=31.06\ \mathrm{N}. \]
\item Then $T_1=1.4142\times31.06=43.92\ \mathrm{N}$.
\item \textbf{Check:} vertical: $43.92\cos30^{\circ}+31.06\cos45^{\circ}=38.03+21.97=60.00\ \mathrm{N}$. \checkmark
\end{enumerate}
\textbf{Answer:} $T_1\approx 43.9\ \mathrm{N}$, $T_2\approx 31.1\ \mathrm{N}$.
\end{exoresolu}

\begin{propriete}[Moment of a force and couple]
The \cle{moment} (or torque) of a force $\vec{F}$ about a point $O$ is $\mathcal{M}_O = F \times d$, where $d$ is the perpendicular distance from $O$ to the line of action of $\vec{F}$. Units: $\mathrm{N\,m}$. Sign convention: anticlockwise positive.
A \cle{couple} consists of two equal and opposite parallel forces $F$ and $-F$ separated by a perpendicular distance $d$. Its moment is $Fd$ (same about any point) and it produces pure rotation without translation.
\end{propriete}

\begin{propriete}[Principle of moments]
For a body in rotational equilibrium under coplanar forces, the sum of the clockwise moments about any point equals the sum of the anticlockwise moments about that point:
\[ \sum \mathcal{M}_{\text{clockwise}} = \sum \mathcal{M}_{\text{anticlockwise}}. \]
Equivalently, $\sum \mathcal{M}_O = 0$ for any point $O$.
\end{propriete}

\begin{methode}[Solving rigid-body equilibrium problems]
\begin{enumerate}
\item Draw a complete \textbf{free-body diagram} showing all forces (weights, reactions, tensions, applied forces) with their points of application.
\item Choose a point to take moments about — ideally a point where \emph{unknown} forces act (their moments are then zero).
\item Write the moment equation $\sum \mathcal{M}=0$.
\item Write the two force-component equations $\sum F_x=0$, $\sum F_y=0$.
\item Solve the system (up to three equations for three unknowns).
\item Check: do the forces balance? Do the moments balance about a \emph{different} point?
\end{enumerate}
\end{methode}

\begin{exoresolu}[Uniform beam on two supports]
A uniform beam $AB$ of length $4.0\ \mathrm{m}$ and mass $20\ \mathrm{kg}$ rests on two supports at $A$ and $C$, where $C$ is $1.0\ \mathrm{m}$ from $B$. A load of mass $30\ \mathrm{kg}$ is placed at $B$. Find the reactions at $A$ and $C$. Take $g=\SI{9.81}{m.s^{-2}}$.
\tcblower
\textbf{Solution.}
\begin{enumerate}
\item \textbf{Forces:} $W_{\text{beam}}=20g=196.2\ \mathrm{N}$ at midpoint ($2.0\ \mathrm{m}$ from $A$); $W_{\text{load}}=30g=294.3\ \mathrm{N}$ at $B$ ($4.0\ \mathrm{m}$ from $A$); reactions $R_A$ at $A$, $R_C$ at $C$ ($3.0\ \mathrm{m}$ from $A$). All vertical.
\item \textbf{Moments about $A$ (anticlockwise +):}
\[ R_C \times 3.0 - 196.2 \times 2.0 - 294.3 \times 4.0 = 0 \]
\[ 3R_C = 392.4 + 1177.2 = 1569.6 \;\Longrightarrow\; R_C = 523.2\ \mathrm{N}. \]
\item \textbf{Vertical equilibrium:}
\[ R_A + R_C = 196.2 + 294.3 = 490.5 \;\Longrightarrow\; R_A = 490.5 - 523.2 = -32.7\ \mathrm{N}. \]
\item \textbf{Interpretation:} $R_A$ is negative, meaning the beam would lift off support $A$; the beam is not in equilibrium on both supports with this load at $B$ — it would tip about $C$. (If the problem insists both supports are in contact, the load position must be changed.)
\end{enumerate}
\textbf{Answer:} With the given data, $R_A$ is negative, so the beam tips about $C$; $A$ loses contact.
\end{exoresolu}

\begin{exoresolu}[Ladder against a smooth wall]
A uniform ladder of length $5.0\ \mathrm{m}$ and mass $15\ \mathrm{kg}$ leans against a smooth vertical wall with its foot on rough horizontal ground. The ladder makes an angle of $60^{\circ}$ with the horizontal. Find (a) the normal reaction of the wall, (b) the frictional force at the ground, (c) the minimum coefficient of friction to prevent slipping. Take $g=\SI{9.81}{m.s^{-2}}$.
\tcblower
\textbf{Solution.}
\begin{enumerate}
\item \textbf{Forces:} $W=15g=147.15\ \mathrm{N}$ at midpoint ($2.5\ \mathrm{m}$ from foot); $R_W$ horizontal at top (smooth wall); $R_G$ vertical at foot; $F$ horizontal at foot (friction, toward wall).
\item \textbf{Moments about foot $O$ (anticlockwise +):}
\[ R_W \times (5\sin60^{\circ}) - W \times (2.5\cos60^{\circ}) = 0 \]
\[ R_W \times 4.330 - 147.15 \times 1.25 = 0 \;\Longrightarrow\; R_W = \frac{183.94}{4.330} = 42.5\ \mathrm{N}. \]
\item \textbf{Horizontal equilibrium:} $F = R_W = 42.5\ \mathrm{N}$.
\item \textbf{Vertical equilibrium:} $R_G = W = 147.15\ \mathrm{N}$.
\item \textbf{Minimum $\mu$:} $F \le \mu R_G \;\Longrightarrow\; \mu_{\min} = \frac{F}{R_G} = \frac{42.5}{147.15} = 0.289$.
\end{enumerate}
\textbf{Answer:} (a) $R_W=42.5\ \mathrm{N}$; (b) $F=42.5\ \mathrm{N}$; (c) $\mu_{\min}=0.289$.
\end{exoresolu}

\coursec{Hooke's Law and Its Experimental Verification}

\begin{propriete}[Hooke's law]
Within the \cle{elastic limit}, the extension $e$ of a spring (or elastic string) is directly proportional to the stretching force $F$:
\[ F = k e, \]
where $k$ is the \cle{spring constant} (stiffness) in $\mathrm{N\,m^{-1}}$. The force exerted \emph{by} the spring is $\vec{F}_{\text{spring}} = -k\vec{e}$ (restoring force).
\end{propriete}

\begin{methode}[Experimental verification of Hooke's law]
\begin{enumerate}
\item Hang a \cle{spring} vertically from a clamp with a pointer and a millimetre ruler beside it. Record the pointer reading $x_0$ with no load (unstretched position).
\item Add a known mass $m$, wait until the spring is still, and record the new reading $x$. The \textbf{extension} is $e=x-x_0$ and the stretching force is $F=mg$.
\item Increase the load in equal steps, recording $x$ each time; \textbf{unload} step by step and check the readings repeat (no permanent stretch).
\item \textbf{Never exceed the elastic limit:} stop if the spring does not return to $x_0$ when unloaded.
\item Plot $F$ (y-axis) against $e$ (x-axis). Hooke's law is verified if the points lie on a \textbf{straight line through the origin}.
\item The gradient of the line is the \textbf{spring constant}: $k=\dfrac{F}{e}$, in $\mathrm{N\,m^{-1}}$.
\end{enumerate}
\end{methode}

\begin{exoresolu}[Analysing a Hooke's law experiment]
A Form Upper Sixth student in Yaoundé obtains the following results for a spring:
\begin{center}
\begin{tabular}{|c|ccccc|}
\hline
\rowcolor{popblueL} Load $m$ (g) & 0 & 50 & 100 & 150 & 200 \\
\hline
Reading $x$ (cm) & 20.0 & 22.5 & 25.0 & 27.4 & 30.1 \\
\hline
\end{tabular}
\end{center}
(a) Compute the extensions and the forces. (b) Determine the spring constant $k$. (c) Estimate the extension for a load of $120\ \mathrm{g}$. (d) The reading for $200\ \mathrm{g}$ seems off the line: comment. Take $g=\SI{9.81}{m.s^{-2}}$.
\tcblower
\textbf{Solution.}
\begin{enumerate}
\item \textbf{(a)} Extensions $e=x-20.0\ \mathrm{cm}$ and forces $F=mg$:
\begin{center}
\begin{tabular}{|c|ccccc|}
\hline
\rowcolor{popblueL} $e$ (cm) & 0 & 2.5 & 5.0 & 7.4 & 10.1 \\
\hline
$F$ (N) & 0 & 0.4905 & 0.981 & 1.4715 & 1.962 \\
\hline
\end{tabular}
\end{center}
\item \textbf{(b)} Using the point $(5.0\ \mathrm{cm},\,0.981\ \mathrm{N})$ on the line:
\[ k=\frac{F}{e}=\frac{0.981}{0.050}=19.62\ \mathrm{N\,m^{-1}}. \]
(Each $50\ \mathrm{g}$ gives about $2.5\ \mathrm{cm}$, confirming proportionality.)
\item \textbf{(c)} For $m=120\ \mathrm{g}$: $F=0.120\times9.81=1.1772\ \mathrm{N}$, so
\[ e=\frac{F}{k}=\frac{1.1772}{19.62}=0.0600\ \mathrm{m}=6.00\ \mathrm{cm}. \]
\item \textbf{(d)} The last point gives $e=10.1\ \mathrm{cm}$ where $10.0\ \mathrm{cm}$ is expected; the small excess ($0.1\ \mathrm{cm}$) may be reading error, but if the spring no longer returns to $20.0\ \mathrm{cm}$ on unloading, the \textbf{elastic limit is being approached} and the load should not be increased further.
\end{enumerate}
\textbf{Answer:} $k\approx 19.6\ \mathrm{N\,m^{-1}}$; $e\approx 6.0\ \mathrm{cm}$ for $120\ \mathrm{g}$.
\end{exoresolu}

\coursec{Applications of Hooke's Law and Synthesis}

\begin{propriete}[Elastic potential energy]
The work done in stretching (or compressing) a spring from extension $0$ to $e$ is stored as \cle{elastic potential energy}:
\[ E_p = \frac{1}{2} k e^2 = \frac{1}{2} F e = \frac{F^2}{2k}. \]
This is the area under the $F$–$e$ graph (a triangle).
\end{propriete}

\begin{methode}[Springs in series and in parallel]
\begin{enumerate}
\item \textbf{Series} (springs end to end): the tension $F$ is the \emph{same} in each spring; extensions add: $e=e_1+e_2=\dfrac{F}{k_1}+\dfrac{F}{k_2}$, so
\[ \frac{1}{k_{\text{eq}}}=\frac{1}{k_1}+\frac{1}{k_2}. \]
\item \textbf{Parallel} (springs side by side, sharing the load): the extension $e$ is the \emph{same} for each; forces add: $F=k_1e+k_2e$, so
\[ k_{\text{eq}}=k_1+k_2. \]
\item \textbf{Strategy:} write $F=ke$ for each spring, use the geometry (which quantities are equal), then solve.
\end{enumerate}
\end{methode}

\begin{exoresolu}[Two springs supporting a platform]
A light platform is supported horizontally by two identical springs in parallel, each of stiffness $k=800\ \mathrm{N\,m^{-1}}$. (a) A bag of cement of mass $40\ \mathrm{kg}$ is placed at the centre; find the compression of each spring. (b) The same two springs are now connected in series and the bag is hung from the lower end; find the total extension. (c) Compute the elastic energy stored in case (a). Take $g=\SI{9.81}{m.s^{-2}}$.
\tcblower
\textbf{Solution.}
\begin{enumerate}
\item Load: $W=mg=40\times9.81=392.4\ \mathrm{N}$.
\item \textbf{(a) Parallel:} $k_{\text{eq}}=k_1+k_2=1600\ \mathrm{N\,m^{-1}}$.
\[ e=\frac{W}{k_{\text{eq}}}=\frac{392.4}{1600}=0.24525\ \mathrm{m}=24.5\ \mathrm{cm}. \]
Each spring compresses by $24.5\ \mathrm{cm}$ (and carries half the load, $196.2\ \mathrm{N}$: check $800\times0.24525=196.2\ \mathrm{N}$ \checkmark).
\item \textbf{(b) Series:} $\dfrac{1}{k_{\text{eq}}}=\dfrac{1}{800}+\dfrac{1}{800}=\dfrac{2}{800}$, so $k_{\text{eq}}=400\ \mathrm{N\,m^{-1}}$.
\[ e_{\text{total}}=\frac{W}{k_{\text{eq}}}=\frac{392.4}{400}=0.981\ \mathrm{m}=98.1\ \mathrm{cm}. \]
(Each spring stretches $49.05\ \mathrm{cm}$; the full tension passes through both.)
\item \textbf{(c) Energy in (a):} for the two springs together,
\[ E_p=2\times\tfrac12 k e^2=800\times(0.24525)^2=800\times0.06015=48.1\ \mathrm{J}. \]
\end{enumerate}
\textbf{Answer:} (a) $24.5\ \mathrm{cm}$ each; (b) $98.1\ \mathrm{cm}$ total; (c) $E_p\approx 48.1\ \mathrm{J}$.
\end{exoresolu}

\begin{exoresolu}[Synthesis: suspended body with spring]
A uniform rod $AB$ of length $1.2\ \mathrm{m}$ and mass $4\ \mathrm{kg}$ is hinged at $A$ to a vertical wall. A light spring of natural length $0.5\ \mathrm{m}$ and stiffness $k=200\ \mathrm{N\,m^{-1}}$ is attached to $B$ and to a point $C$ on the wall $1.3\ \mathrm{m}$ directly above $A$. The rod is held horizontal and released. Find the extension of the spring when the rod comes to rest in equilibrium. Take $g=\SI{9.81}{m.s^{-2}}$.
\tcblower
\textbf{Solution.}
\begin{enumerate}
\item \textbf{Geometry:} when the rod is horizontal, $B$ is $1.2\ \mathrm{m}$ horizontally from $A$. $C$ is $1.3\ \mathrm{m}$ vertically above $A$. So $BC = \sqrt{1.2^2+1.3^2} = \sqrt{1.44+1.69} = \sqrt{3.13} = 1.769\ \mathrm{m}$.
\item Extension $e = BC - l_0 = 1.769 - 0.5 = 1.269\ \mathrm{m}$.
\item Spring force $F_s = k e = 200 \times 1.269 = 253.8\ \mathrm{N}$ along $BC$ (toward $C$).
\item \textbf{Moments about $A$:} Weight $W=4g=39.24\ \mathrm{N}$ acts at midpoint ($0.6\ \mathrm{m}$ from $A$), vertically down. The spring force at $B$ has a vertical component $F_s \sin\theta$ and horizontal component $F_s \cos\theta$, where $\sin\theta = 1.3/1.769 = 0.735$, $\cos\theta = 1.2/1.769 = 0.678$.
\item Only vertical forces contribute moment about $A$ (horizontal force at $B$ passes through $A$? No, $A$ is at $(0,0)$, $B$ at $(1.2,0)$, horizontal force at $B$ has moment arm $0$ vertically? Wait: horizontal force at $B$ is perpendicular to vertical distance from $A$? Actually, moment of a force about $A$ is force $\times$ perpendicular distance. The horizontal component of spring force acts along the rod (horizontal), so its line of action passes through $A$? No, $A$ is at the hinge, the rod is horizontal. A horizontal force at $B$ is along the rod, so its line of action passes through $A$ — moment zero. The vertical component of spring force at $B$ is upward, perpendicular distance $1.2\ \mathrm{m}$.
\item Moment equilibrium about $A$ (anticlockwise +):
\[ F_s \sin\theta \times 1.2 - W \times 0.6 = 0 \]
\[ 253.8 \times 0.735 \times 1.2 - 39.24 \times 0.6 = 223.7 - 23.5 = 200.2 \neq 0. \]
\item \textbf{Contradiction:} the rod is not in equilibrium in the horizontal position with this spring extension. The problem asks for the extension \emph{when the rod comes to rest in equilibrium} — so the rod will settle at some angle $\alpha$ below horizontal. We must solve for the equilibrium angle.
\item \textbf{Let the rod make angle $\alpha$ below horizontal.} Then coordinates: $A=(0,0)$, $B=(1.2\cos\alpha,\,-1.2\sin\alpha)$, $C=(0,1.3)$.
\item Length $BC = \sqrt{(1.2\cos\alpha)^2 + (1.3+1.2\sin\alpha)^2}$.
\item Extension $e = BC - 0.5$.
\item Spring force magnitude $F_s = k e = 200(BC-0.5)$.
\item Direction of spring force: from $B$ to $C$, vector $\vec{BC} = (-1.2\cos\alpha,\, 1.3+1.2\sin\alpha)$.
\item Moment of spring force about $A$: $\vec{r}_B \times \vec{F}_s$. $\vec{r}_B = (1.2\cos\alpha,\,-1.2\sin\alpha)$. $\vec{F}_s = F_s \frac{\vec{BC}}{BC}$.
\item Moment (scalar, out of plane): $M_s = x_B F_{s,y} - y_B F_{s,x}$.
\item $F_{s,y} = F_s \frac{1.3+1.2\sin\alpha}{BC}$, $F_{s,x} = F_s \frac{-1.2\cos\alpha}{BC}$.
\item $M_s = (1.2\cos\alpha) F_s \frac{1.3+1.2\sin\alpha}{BC} - (-1.2\sin\alpha) F_s \frac{-1.2\cos\alpha}{BC} = \frac{1.2 F_s \cos\alpha}{BC} \big[1.3+1.2\sin\alpha - 1.2\sin\alpha\big] = \frac{1.2 F_s \cos\alpha \times 1.3}{BC}$.
\item Moment of weight about $A$: $W \times (0.6\cos\alpha)$ clockwise (negative).
\item Equilibrium: $M_s - 0.6 W \cos\alpha = 0 \;\Longrightarrow\; \frac{1.2 \times 1.3 F_s}{BC} = 0.6 W \;\Longrightarrow\; \frac{2.6 F_s}{BC} = W$.
\item But $F_s = 200(BC-0.5)$, so $\frac{2.6 \times 200 (BC-0.5)}{BC} = 39.24$.
\item $520 \frac{BC-0.5}{BC} = 39.24 \;\Longrightarrow\; \frac{BC-0.5}{BC} = 0.07546 \;\Longrightarrow\; 1 - \frac{0.5}{BC} = 0.07546 \;\Longrightarrow\; \frac{0.5}{BC} = 0.92454 \;\Longrightarrow\; BC = 0.5408\ \mathrm{m}$.
\item Then $e = BC - 0.5 = 0.0408\ \mathrm{m} = 4.08\ \mathrm{cm}$.
\item Check: $BC=0.5408$, but $BC$ must be at least the vertical distance $1.3\ \mathrm{m}$? Wait, $C$ is at $(0,1.3)$, $B$ is at $(1.2\cos\alpha,\,-1.2\sin\alpha)$. The vertical distance is $1.3+1.2\sin\alpha \ge 1.3$. So $BC \ge 1.3$. Our computed $BC=0.54$ is impossible.
\item \textbf{Error:} The spring is attached to $C$ which is $1.3\ \mathrm{m}$ \emph{above} $A$. If the rod is horizontal, $B$ is at $(1.2,0)$, so $BC = \sqrt{1.2^2+1.3^2}=1.769 > 1.3$. The spring is stretched. But the moment calculation gave a huge anticlockwise moment from the spring. The weight moment is small. The spring would pull the rod up, not down. So the equilibrium position is with the rod \emph{above} horizontal? Let's re-read: "hinged at $A$ to a vertical wall. A light spring ... attached to $B$ and to a point $C$ on the wall $1.3\ \mathrm{m}$ directly above $A$." If the rod is held horizontal and released, the spring pulls upward on $B$, lifting the rod. The rod will rise until the spring force moment balances the weight moment. The angle $\alpha$ should be measured \emph{above} horizontal.
\item Let $\alpha$ be angle \emph{above} horizontal. Then $B=(1.2\cos\alpha,\, 1.2\sin\alpha)$. $C=(0,1.3)$. Vertical distance $= 1.3 - 1.2\sin\alpha$. This must be positive, so $\sin\alpha < 1.3/1.2$ (always true). $BC = \sqrt{(1.2\cos\alpha)^2 + (1.3-1.2\sin\alpha)^2}$.
\item Moment of spring force about $A$: similar derivation gives $M_s = \frac{1.2 F_s \cos\alpha \times 1.3}{BC}$ (anticlockwise? Let's check sign). $\vec{r}_B = (1.2\cos\alpha,\, 1.2\sin\alpha)$. $\vec{BC} = (-1.2\cos\alpha,\, 1.3-1.2\sin\alpha)$. $M_s = x_B F_{s,y} - y_B F_{s,x} = (1.2\cos\alpha) F_s \frac{1.3-1.2\sin\alpha}{BC} - (1.2\sin\alpha) F_s \frac{-1.2\cos\alpha}{BC} = \frac{1.2 F_s \cos\alpha}{BC} [1.3-1.2\sin\alpha + 1.2\sin\alpha] = \frac{1.56 F_s \cos\alpha}{BC}$. This is positive (anticlockwise).
\item Weight moment: $W \times (0.6\cos\alpha)$ clockwise (negative).
\item Equilibrium: $\frac{1.56 F_s \cos\alpha}{BC} - 0.6 W \cos\alpha = 0$. If $\cos\alpha \neq 0$, $\frac{1.56 F_s}{BC} = 0.6 W = 23.544$.
\item $F_s = 200(BC-0.5)$. So $\frac{1.56 \times 200 (BC-0.5)}{BC} = 23.544 \;\Longrightarrow\; 312 \frac{BC-0.5}{BC} = 23.544 \;\Longrightarrow\; \frac{BC-0.5}{BC} = 0.07546$.
\item Same equation! $BC = 0.5408\ \mathrm{m}$. But $BC \ge |1.3 - 1.2| = 0.1$? Actually minimum $BC$ is when $B$ is closest to $C$. $B$ moves on circle radius $1.2$. $C$ is at distance $1.3$ from $A$. Minimum $BC = 1.3 - 1.2 = 0.1\ \mathrm{m}$. Maximum $BC = 2.5\ \mathrm{m}$. So $BC=0.5408$ is possible! It corresponds to the rod pointing somewhat upward, with $B$ close to $C$.
\item Then $e = 0.5408 - 0.5 = 0.0408\ \mathrm{m} = 4.08\ \mathrm{cm}$.
\item Find $\alpha$: $BC^2 = 1.2^2 + 1.3^2 - 2\times1.2\times1.3\sin\alpha = 3.13 - 3.12\sin\alpha$.
\item $0.5408^2 = 0.2925 = 3.13 - 3.12\sin\alpha \;\Longrightarrow\; 3.12\sin\alpha = 2.8375 \;\Longrightarrow\; \sin\alpha = 0.9095 \;\Longrightarrow\; \alpha = 65.4^{\circ}$ above horizontal.
\end{enumerate}
\textbf{Answer:} The spring extends by $4.1\ \mathrm{cm}$ when the rod rests at $65.4^{\circ}$ above the horizontal.
\end{exoresolu}

\begin{attention}[Frequent errors in this chapter]
\begin{itemize}
\item Confusing \emph{resultant} and \emph{equilibrant}: the equilibrant has the \textbf{same magnitude} but the \textbf{opposite direction}.
\item Resolving with $\sin$ and $\cos$ interchanged: the component \emph{adjacent} to the given angle uses $\cos$.
\item On an inclined plane, resolving the weight into $W\sin\alpha$ and $W\cos\alpha$ along the wrong directions — always draw the angle first.
\item Measuring the extension of a spring from the top instead of from the \textbf{unstretched length}: $e=x-x_0$.
\item Applying $F=ke$ beyond the \textbf{elastic limit}, where the graph is no longer linear.
\item Forgetting that in equilibrium problems you may write \emph{two} independent scalar equations ($\sum F_x=0$, $\sum F_y=0$) — and a moment equation if the body can rotate.
\item Taking moments about a point where an unknown force acts, but forgetting that force's moment is zero only if its line of action passes through the point.
\item Sign errors in moment equations: define a clear convention (e.g. anticlockwise +) and stick to it.
\item In series/parallel springs: confusing which quantity is the same (force in series, extension in parallel).
\end{itemize}
\end{attention}

\begin{exercice}[Practice problems]
\begin{enumerate}
\item A uniform rod $AB$ of length $2\ \mathrm{m}$ and weight $50\ \mathrm{N}$ is suspended by two vertical strings attached at $A$ and $B$. A weight of $30\ \mathrm{N}$ is hung from the rod at a point $0.5\ \mathrm{m}$ from $A$. Find the tension in each string.
\item A force of $100\ \mathrm{N}$ acts at $30^{\circ}$ to the horizontal. Find its horizontal and vertical components.
\item Three coplanar forces act at a point: $20\ \mathrm{N}$ at $0^{\circ}$, $30\ \mathrm{N}$ at $120^{\circ}$, $40\ \mathrm{N}$ at $240^{\circ}$. Find the resultant analytically.
\item A spring of natural length $0.3\ \mathrm{m}$ and stiffness $150\ \mathrm{N\,m^{-1}}$ is stretched to $0.5\ \mathrm{m}$. Find the force and the energy stored.
\item Two springs of stiffness $200\ \mathrm{N\,m^{-1}}$ and $300\ \mathrm{N\,m^{-1}}$ are connected (a) in series, (b) in parallel. Find the equivalent stiffness in each case.
\item A uniform ladder of length $6\ \mathrm{m}$ and mass $20\ \mathrm{kg}$ leans against a smooth wall at $60^{\circ}$ to the horizontal. A man of mass $70\ \mathrm{kg}$ climbs $4\ \mathrm{m}$ up the ladder. Find the friction force at the ground and the minimum coefficient of friction. ($g=9.81\ \mathrm{m\,s^{-2}}$)
\item A triangular lamina of mass $2\ \mathrm{kg}$ has vertices at $(0,0)$, $(0.3,0)$, $(0,0.4)$ (in metres). It is suspended from the vertex at $(0,0)$. Find the angle the side on the $x$-axis makes with the vertical at equilibrium.
\end{enumerate}
\end{exercice}

\begin{corrige}[Exercise 1]
Let tensions be $T_A$ at $A$ and $T_B$ at $B$. Moments about $A$: $T_B \times 2 - 50 \times 1 - 30 \times 0.5 = 0 \Rightarrow 2T_B = 50 + 15 = 65 \Rightarrow T_B = 32.5\ \mathrm{N}$. Vertical equilibrium: $T_A + T_B = 80 \Rightarrow T_A = 47.5\ \mathrm{N}$.
\end{corrige}

\begin{corrige}[Exercise 2]
$F_x = 100\cos30^{\circ} = 86.6\ \mathrm{N}$, $F_y = 100\sin30^{\circ} = 50\ \mathrm{N}$.
\end{corrige}

\begin{corrige}[Exercise 3]
$R_x = 20 + 30\cos120^{\circ} + 40\cos240^{\circ} = 20 - 15 - 20 = -15\ \mathrm{N}$.
$R_y = 0 + 30\sin120^{\circ} + 40\sin240^{\circ} = 25.98 - 34.64 = -8.66\ \mathrm{N}$.
$R = \sqrt{15^2 + 8.66^2} = 17.3\ \mathrm{N}$ at $\tan^{-1}(8.66/15)=30^{\circ}$ south of west (i.e. $210^{\circ}$).
\end{corrige}

\begin{corrige}[Exercise 4]
$e = 0.5 - 0.3 = 0.2\ \mathrm{m}$. $F = 150 \times 0.2 = 30\ \mathrm{N}$. $E_p = \frac12 \times 150 \times 0.2^2 = 3\ \mathrm{J}$.
\end{corrige}

\begin{corrige}[Exercise 5]
(a) Series: $1/k_{\text{eq}} = 1/200 + 1/300 = 5/600 \Rightarrow k_{\text{eq}} = 120\ \mathrm{N\,m^{-1}}$.
(b) Parallel: $k_{\text{eq}} = 200 + 300 = 500\ \mathrm{N\,m^{-1}}$.
\end{corrige}

\begin{corrige}[Exercise 6]
Forces: $W_L=196.2\ \mathrm{N}$ at $3\ \mathrm{m}$ from foot; $W_M=686.7\ \mathrm{N}$ at $4\ \mathrm{m}$ from foot; $R_W$ at top ($6\sin60^{\circ}=5.196\ \mathrm{m}$ up); $F$ and $R_G$ at foot.
Moments about foot: $R_W \times 5.196 - 196.2 \times 3\cos60^{\circ} - 686.7 \times 4\cos60^{\circ} = 0$.
$R_W \times 5.196 = 196.2 \times 1.5 + 686.7 \times 2 = 294.3 + 1373.4 = 1667.7 \Rightarrow R_W = 321.0\ \mathrm{N}$.
$F = R_W = 321.0\ \mathrm{N}$. $R_G = 196.2 + 686.7 = 882.9\ \mathrm{N}$.
$\mu_{\min} = 321.0 / 882.9 = 0.364$.
\end{corrige}

\begin{corrige}[Exercise 7]
Centroid $G = (0.1, 0.1333)\ \mathrm{m}$. Suspended from $(0,0)$, $G$ hangs vertically below. Line $OG$ has slope $0.1333/0.1 = 1.333$, angle with horizontal $= \tan^{-1}(1.333) = 53.13^{\circ}$. The side on $x$-axis is horizontal, so it makes $90^{\circ} - 53.13^{\circ} = 36.87^{\circ}$ with the vertical.
\end{corrige}

\newpage

\coursec{Exercises}

\courssub{I check my knowledge}

\begin{exercice}[MCQ: resultant of two forces]
Two forces of magnitudes $6\ \mathrm{N}$ and $8\ \mathrm{N}$ act at a point $O$ at right angles to each other. The magnitude of their resultant is:
\begin{enumerate}
\item[(a)] $2\ \mathrm{N}$ \quad (b) $10\ \mathrm{N}$ \quad (c) $14\ \mathrm{N}$ \quad (d) $48\ \mathrm{N}$
\end{enumerate}
Choose the correct answer and justify your choice with a calculation.
\end{exercice}

\begin{exercice}[True or false]
For each of the following statements, say whether it is \textbf{true} or \textbf{false} and justify your answer briefly.
\begin{enumerate}
\item A body is in equilibrium whenever the sum of the magnitudes of the forces acting on it is zero.
\item The centre of gravity of a uniform body always lies inside the material of the body.
\item When three non-parallel coplanar forces keep a body in equilibrium, their lines of action are concurrent.
\item Hooke's law states that the extension of a spring is proportional to the load, whatever the load may be.
\end{enumerate}
\end{exercice}

\begin{exercice}[Definitions]
Define each of the following terms, giving one example in each case:
\begin{enumerate}
\item the \cle{resultant} of a system of coplanar forces;
\item the \cle{equilibrant} of a system of coplanar forces;
\item a \cle{couple} of forces;
\item the \cle{centre of gravity} of a body.
\end{enumerate}
\end{exercice}

\begin{exercice}[Direct calculation]
A force of $20\ \mathrm{N}$ acts at an angle of $40^{\circ}$ above the horizontal.
\begin{enumerate}
\item Resolve the force into horizontal and vertical components.
\item A spring of stiffness $k = 150\ \mathrm{N\,m^{-1}}$ is stretched by a force of $12\ \mathrm{N}$. Assuming Hooke's law holds, calculate the extension in centimetres.
\end{enumerate}
\end{exercice}

\courssub{I practise}

\begin{exercice}[Resultant of two forces]
Two forces of $8\ \mathrm{N}$ and $12\ \mathrm{N}$ act at a point $O$ with an angle of $60^{\circ}$ between them.
\begin{enumerate}
\item Using the cosine rule (parallelogram of forces), calculate the magnitude of the resultant.
\item Calculate the angle the resultant makes with the $12\ \mathrm{N}$ force.
\item Verify your answers with an accurate scale drawing (state the scale used).
\end{enumerate}
\end{exercice}

\begin{exercice}[Block on an inclined plane]
A block of weight $100\ \mathrm{N}$ rests on a smooth plane inclined at $30^{\circ}$ to the horizontal. It is held in equilibrium by a force $P$ acting parallel to the plane.
\begin{enumerate}
\item Draw a clearly labelled free-body diagram of the block.
\item By resolving parallel and perpendicular to the plane, find $P$ and the normal reaction $R$.
\end{enumerate}
\end{exercice}

\begin{exercice}[Springs in series and in parallel]
Two identical springs, each of stiffness $k = 50\ \mathrm{N\,m^{-1}}$, support a load of $10\ \mathrm{N}$.
\begin{enumerate}
\item The springs are connected in series. Find the extension of each spring and the total extension.
\item The springs are connected in parallel, sharing the load equally. Find the total extension.
\item Explain, in terms of Hooke's law, why the two arrangements give different total extensions.
\end{enumerate}
\end{exercice}

\begin{exercice}[Hanging picture]
A picture of weight $40\ \mathrm{N}$ is hung by a string passing over a nail. The two parts of the string make angles of $50^{\circ}$ and $70^{\circ}$ with the vertical.
\begin{enumerate}
\item Draw the free-body diagram of the picture.
\item Using Lami's theorem, find the tension in each part of the string.
\item Check your answers by resolving horizontally and vertically.
\end{enumerate}
\end{exercice}

\courssub{I prepare for the GCE}

\begin{exercice}[GCE-style: string and spring in equilibrium]
A body of mass $5\ \mathrm{kg}$ hangs in equilibrium. It is supported by a light inextensible string inclined at $35^{\circ}$ to the vertical and by a horizontal spring of stiffness $200\ \mathrm{N\,m^{-1}}$ attached to a wall. Take $g = 9.81\ \mathrm{m\,s^{-2}}$.
\begin{enumerate}
\item Draw a free-body diagram of the body, showing all the forces acting on it. \hfill (3 marks)
\item Write down the equations obtained by resolving the forces horizontally and vertically. \hfill (4 marks)
\item Calculate the tension in the string. \hfill (3 marks)
\item Calculate the extension of the spring, giving your answer in centimetres. State clearly the law you use and the assumption made. \hfill (5 marks)
\end{enumerate}
\end{exercice}

\begin{exercice}[GCE-style: experimental verification of Hooke's law]
In an experiment to investigate Hooke's law, a student suspends a spring vertically and records the extension $e$ for increasing loads $W$:
\begin{center}
\begin{tabular}{|l|ccccc|}
\hline
\rowcolor{popblueL} Load $W$ (N) & 1 & 2 & 3 & 4 & 5 \\
\hline
Extension $e$ (cm) & 2.1 & 4.0 & 6.2 & 8.1 & 10.3 \\
\hline
\end{tabular}
\end{center}
\begin{enumerate}
\item Describe briefly how the student should carry out this experiment, including the precautions taken to obtain accurate readings. \hfill (4 marks)
\item Plot the graph of $e$ against $W$ on graph paper, using suitable scales. \hfill (4 marks)
\item Explain how the graph verifies Hooke's law, and use it to determine the stiffness $k$ of the spring in $\mathrm{N\,m^{-1}}$. \hfill (4 marks)
\item The reading at $W = 5\ \mathrm{N}$ lies slightly off the line. Suggest two possible reasons for this and state what the student should check before concluding that the elastic limit has been reached. \hfill (3 marks)
\end{enumerate}
\end{exercice}

\begin{exercice}[GCE-style: centre of gravity and stability]
\begin{enumerate}
\item A uniform lamina in the shape of a rectangle with a triangular notch cut from one side is suspended freely from a corner. Describe, with the aid of a diagram, the experiment you would perform to locate its centre of gravity, and explain why the lamina hangs so that the line from the point of suspension through the centre of gravity is vertical. \hfill (6 marks)
\item A loaded wheelbarrow is modelled as a rigid body resting on the ground. Using the position of the centre of gravity and the base of support, explain why the wheelbarrow topples forward when the load is placed too far in front of the wheel. Illustrate your answer with a diagram showing the line of action of the weight in the stable and in the toppling cases. \hfill (6 marks)
\item Three coplanar forces act at a point and keep a body in equilibrium. Two of the forces are $10\ \mathrm{N}$ due east and $15\ \mathrm{N}$ at $120^{\circ}$ measured anticlockwise from east. Find, by resolving, the magnitude and direction of the third force, and explain how the same result could be obtained graphically by the polygon of forces. \hfill (8 marks)
\end{enumerate}
\end{exercice}

\newpage

\coursec{Solutions to the exercises}

\begin{corrige}[Exercise 1 — MCQ: resultant of two forces]
The correct answer is \textbf{(b)} $10\ \mathrm{N}$.

\textbf{Justification.} The two forces are perpendicular, so by the parallelogram rule (here a rectangle), the magnitude of the resultant is given by Pythagoras' theorem:
\[
R=\sqrt{6^{2}+8^{2}}=\sqrt{36+64}=\sqrt{100}=10\ \mathrm{N}.
\]
The other options correspond to common errors: $14\ \mathrm{N}$ (adding the magnitudes arithmetically), $2\ \mathrm{N}$ (subtracting them), $48\ \mathrm{N}$ (multiplying them). Forces are vectors: only the \emph{vector} sum is physically meaningful.
\end{corrige}

\begin{corrige}[Exercise 2 — True or false]
\begin{enumerate}
\item \textbf{False.} Equilibrium requires the \emph{vector} sum of the forces to be zero, $\sum\vec{F}=\vec{0}$, not the sum of the magnitudes. For example, two forces of $5\ \mathrm{N}$ and $3\ \mathrm{N}$ acting in opposite directions have magnitudes summing to $8\ \mathrm{N}$, yet their vector sum has magnitude $2\ \mathrm{N}\neq 0$, so the body is not in equilibrium. The correct condition is that the resultant \emph{vector} vanishes (and, for an extended body, that the total moment about any point is also zero).
\item \textbf{False.} The centre of gravity may lie outside the material of the body: for a uniform ring it is at the centre of the hole; for a boomerang or an L-shaped lamina it can lie in empty space. What matters is that the line of action of the total weight always passes through it.
\item \textbf{True.} If three non-parallel coplanar forces are in equilibrium, the resultant of any two of them must balance the third. This resultant acts through the point of intersection of the first two lines of action, so the third force must pass through the same point: the three lines of action are \cle{concurrent}. (This is the \emph{three-force principle}, very useful for locating unknown forces in GCE problems.)
\item \textbf{False.} Hooke's law ($F=ke$) holds only \emph{provided the elastic limit is not exceeded}. Beyond the elastic limit the extension is no longer proportional to the load and the spring suffers permanent deformation: it does not return to its original length when unloaded.
\end{enumerate}
\end{corrige}

\begin{corrige}[Exercise 3 — Definitions]
\begin{enumerate}
\item The \cle{resultant} of a system of coplanar forces is the single force, acting at the same point, which has the same effect as all the forces together; it is their vector sum $\vec{R}=\sum\vec{F}_{i}$. \emph{Example:} forces of $3\ \mathrm{N}$ due east and $4\ \mathrm{N}$ due north have a resultant of $5\ \mathrm{N}$ in the direction $\arctan(4/3)\approx 53.1^{\circ}$ north of east.
\item The \cle{equilibrant} is the single force which, added to the system, produces equilibrium; it is equal in magnitude to the resultant, acts along the same line, but in the opposite direction: $\vec{E}=-\vec{R}$. \emph{Example:} for the system above, the equilibrant is $5\ \mathrm{N}$ directed $53.1^{\circ}$ south of west.
\item A \cle{couple} is a pair of equal, parallel, oppositely directed forces whose lines of action are distinct. Its resultant force is zero, but it produces a pure turning effect (moment) $\tau = Fd$, where $d$ is the perpendicular distance between the lines of action. \emph{Example:} the two forces applied by the hands on a steering wheel, or the fingers turning a tap.
\item The \cle{centre of gravity} of a body is the point through which the line of action of its total weight always passes, whatever the orientation of the body; equivalently, it is the point about which the weights of all the particles of the body have zero total moment. \emph{Example:} the centre of gravity of a uniform metre rule is at the $50\ \mathrm{cm}$ mark.
\end{enumerate}
\end{corrige}

\begin{corrige}[Exercise 4 — Direct calculation]
\begin{enumerate}
\item Let $\vec{F}$ have magnitude $F=20\ \mathrm{N}$ at $40^{\circ}$ above the horizontal. The component adjacent to the angle uses cosine, the component opposite uses sine:
\[
F_{x}=F\cos 40^{\circ}=20\times 0.7660\approx 15.3\ \mathrm{N},\qquad
F_{y}=F\sin 40^{\circ}=20\times 0.6428\approx 12.9\ \mathrm{N}.
\]
The horizontal component is about $15.3\ \mathrm{N}$ and the vertical component about $12.9\ \mathrm{N}$ (upwards). \emph{Check:} $\sqrt{15.3^{2}+12.9^{2}}\approx 20.0\ \mathrm{N}$. \checkmark
\item By Hooke's law, $F=ke$, so
\[
e=\frac{F}{k}=\frac{12}{150}=0.08\ \mathrm{m}=8\ \mathrm{cm}.
\]
The extension is $8\ \mathrm{cm}$ (assuming the elastic limit is not exceeded).
\end{enumerate}
\end{corrige}

\begin{corrige}[Exercise 5 — Resultant of two forces]
\begin{enumerate}
\item By the cosine rule applied to the parallelogram (triangle) of forces, with the angle between the forces $\theta=60^{\circ}$:
\[
R^{2}=8^{2}+12^{2}+2(8)(12)\cos 60^{\circ}=64+144+192\times 0.5=208+96=304.
\]
\[
R=\sqrt{304}\approx 17.4\ \mathrm{N}.
\]
\item Let $\alpha$ be the angle between the resultant and the $12\ \mathrm{N}$ force. Resolving the $8\ \mathrm{N}$ force parallel and perpendicular to the $12\ \mathrm{N}$ force:
\[
\tan\alpha=\frac{8\sin 60^{\circ}}{12+8\cos 60^{\circ}}=\frac{8\times 0.8660}{12+4}=\frac{6.928}{16}=0.4330,
\qquad \alpha\approx 23.4^{\circ}.
\]
The resultant makes an angle of about $23.4^{\circ}$ with the $12\ \mathrm{N}$ force.
\item \textbf{Scale drawing.} Choose the scale $1\ \mathrm{cm}:1\ \mathrm{N}$. From a point $O$ draw a segment of $12\ \mathrm{cm}$ (the $12\ \mathrm{N}$ force) along a reference line; using a protractor, draw from $O$ a segment of $8\ \mathrm{cm}$ at $60^{\circ}$ to it. Complete the parallelogram and draw the diagonal from $O$. Measuring the diagonal gives about $17.4\ \mathrm{cm}$, i.e. $R\approx 17.4\ \mathrm{N}$, and measuring the angle between the diagonal and the $12\ \mathrm{cm}$ side gives about $23^{\circ}$. This confirms the calculation.
\end{enumerate}
\end{corrige}

\begin{corrige}[Exercise 6 — Block on an inclined plane]
\begin{enumerate}
\item \textbf{Free-body diagram.} Three forces act on the block: the weight $W=100\ \mathrm{N}$ vertically downwards, the normal reaction $R$ perpendicular to the plane, and the applied force $P$ parallel to the plane, directed up the slope.
\begin{popfigure}
\begin{center}
\begin{tikzpicture}[scale=1.0,>=latex]
\draw[popink,thick] (0,0) -- (5,0);
\draw[popink,thick] (0,0) -- (5,2.89);
\draw[popmuted] (4.2,0) -- (5,0.578);
\draw[popmuted] (1.2,0) arc (0:30:1.2) node[midway,right=2pt] {\small $30^{\circ}$};
\draw[fill=popblueL,draw=popink] (2.2,1.27) rectangle (3.0,1.97);
\draw[->,popblue,very thick] (2.6,1.62) -- (2.6,0.22) node[below right] {\small $W=100$ N};
\draw[->,popgreen,very thick] (2.6,1.97) -- (2.15,3.05) node[above left] {\small $R$};
\draw[->,poporange,very thick] (3.0,1.62) -- (4.2,2.31) node[above right] {\small $P$};
\end{tikzpicture}
\end{center}
\legende{Free-body diagram of the block on the smooth inclined plane.}
\end{popfigure}
\item The plane is smooth, so there is no friction. Resolve the weight: the component parallel to the plane (down the slope) is $W\sin 30^{\circ}$; the component perpendicular to the plane (into the plane) is $W\cos 30^{\circ}$. (The angle between the weight and the normal to the plane equals the angle of inclination, $30^{\circ}$.)

\textbf{Equilibrium parallel to the plane:}
\[
P=W\sin 30^{\circ}=100\times 0.5=50\ \mathrm{N}.
\]
\textbf{Equilibrium perpendicular to the plane:}
\[
R=W\cos 30^{\circ}=100\times\frac{\sqrt{3}}{2}\approx 86.6\ \mathrm{N}.
\]
Hence $P=50\ \mathrm{N}$ up the plane and $R\approx 86.6\ \mathrm{N}$. Note that $R<W$: on an incline the plane supports only part of the weight.
\end{enumerate}
\end{corrige}

\begin{corrige}[Exercise 7 — Springs in series and in parallel]
\begin{enumerate}
\item \textbf{Series.} In series, the same tension $T=10\ \mathrm{N}$ acts in each spring. By Hooke's law, the extension of each spring is
\[
e=\frac{T}{k}=\frac{10}{50}=0.2\ \mathrm{m}=20\ \mathrm{cm}.
\]
The total extension is the sum:
\[
e_{\text{total}}=0.2+0.2=0.4\ \mathrm{m}=40\ \mathrm{cm}.
\]
\item \textbf{Parallel.} The load is shared equally between the two identical springs, so each carries $\dfrac{10}{2}=5\ \mathrm{N}$. Each spring (and hence the combination, since both ends move together) extends by
\[
e=\frac{5}{50}=0.1\ \mathrm{m}=10\ \mathrm{cm}.
\]
The total extension is $0.1\ \mathrm{m}=10\ \mathrm{cm}$.
\item \textbf{Explanation.} In series, the \emph{full} load passes through each spring, so each stretches by $F/k$ and the extensions add: the combination is \emph{softer} (effective stiffness $k_{\text{eff}}=k/2=25\ \mathrm{N\,m^{-1}}$, since $e_{\text{total}}=2F/k$). In parallel, each spring carries only \emph{half} the load, so each stretches less and the extensions do not add: the combination is \emph{stiffer} (effective stiffness $k_{\text{eff}}=2k=100\ \mathrm{N\,m^{-1}}$, since the total force $2ke$ is shared). In both cases the result follows directly from $e=F/k$ applied to the force actually carried by each spring.
\end{enumerate}
\end{corrige}

\begin{corrige}[Exercise 8 — Hanging picture]
\begin{enumerate}
\item \textbf{Free-body diagram.} Three forces act on the picture: its weight $W=40\ \mathrm{N}$ vertically downwards, and the tensions $T_{1}$ and $T_{2}$ along the two parts of the string, at $50^{\circ}$ and $70^{\circ}$ to the vertical respectively (on opposite sides of the vertical).
\begin{popfigure}
\begin{center}
\begin{tikzpicture}[scale=1.0,>=latex]
\draw[fill=popblueL,draw=popink] (-1.2,-1.4) rectangle (1.2,0);
\node at (0,-0.7) {\small picture};
\draw[->,popblue,very thick] (0,-1.4) -- (0,-2.6) node[below] {\small $W=40$ N};
\draw[->,poporange,very thick] (0,0) -- (-1.75,2.14) node[above left] {\small $T_{1}$};
\draw[->,popgreen,very thick] (0,0) -- (1.88,1.71) node[above right] {\small $T_{2}$};
\draw[popmuted,dashed] (0,0) -- (0,2.4);
\draw[popmuted] (0,1.2) arc (90:140:1.2) node[midway,above left=-1pt] {\small $50^{\circ}$};
\draw[popmuted] (0,0.9) arc (90:20:0.9) node[midway,above right=-1pt] {\small $70^{\circ}$};
\end{tikzpicture}
\end{center}
\legende{Forces on the hanging picture.}
\end{popfigure}
\item \textbf{Lami's theorem.} Since the two strings lie on opposite sides of the vertical, the angle between $T_{1}$ and $T_{2}$ is $50^{\circ}+70^{\circ}=120^{\circ}$; the angle between $T_{1}$ and $W$ is $180^{\circ}-50^{\circ}=130^{\circ}$; the angle between $T_{2}$ and $W$ is $180^{\circ}-70^{\circ}=110^{\circ}$. (Check: $120^{\circ}+130^{\circ}+110^{\circ}=360^{\circ}$. \checkmark) Lami's theorem states that each force is proportional to the sine of the angle between the \emph{other two}:
\[
\frac{T_{1}}{\sin 110^{\circ}}=\frac{T_{2}}{\sin 130^{\circ}}=\frac{W}{\sin 120^{\circ}}
=\frac{40}{0.8660}=46.19\ \mathrm{N}.
\]
\[
T_{1}=46.19\sin 110^{\circ}=46.19\times 0.9397\approx 43.4\ \mathrm{N},
\]
\[
T_{2}=46.19\sin 130^{\circ}=46.19\times 0.7660\approx 35.4\ \mathrm{N}.
\]
\item \textbf{Check by resolution.} Horizontally (taking components of the tensions measured from the vertical):
\[
T_{1}\sin 50^{\circ}=43.4\times 0.7660\approx 33.2\ \mathrm{N},\qquad
T_{2}\sin 70^{\circ}=35.4\times 0.9397\approx 33.3\ \mathrm{N},
\]
which balance (the small difference is due to rounding). Vertically:
\[
T_{1}\cos 50^{\circ}+T_{2}\cos 70^{\circ}
=43.4\times 0.6428+35.4\times 0.3420
\approx 27.9+12.1=40.0\ \mathrm{N}=W. \ \checkmark
\]
\end{enumerate}
\end{corrige}

\begin{corrige}[Exercise 9 — GCE-style: string and spring in equilibrium]
\begin{enumerate}
\item \textbf{Free-body diagram} \hfill (3 marks: one for each correctly drawn and labelled force)

Three forces act on the body: the weight $W=mg=5\times 9.81=49.05\ \mathrm{N}$ vertically downwards; the tension $T$ along the string, at $35^{\circ}$ to the vertical, pulling up and away from the wall; the spring force $F$ horizontal, pulling the body towards the wall.
\begin{popfigure}
\begin{center}
\begin{tikzpicture}[scale=1.0,>=latex]
\fill[popblueL] (0,0) circle (0.28);
\draw[popink] (0,0) circle (0.28);
\draw[->,popblue,very thick] (0,-0.28) -- (0,-2.2) node[below] {\small $W=mg$};
\draw[->,poporange,very thick] (0,0.28) -- (-1.15,2.15) node[above left] {\small $T$};
\draw[->,popgreen,very thick] (0.28,0) -- (2.2,0) node[right] {\small $F=ke$};
\draw[popmuted,dashed] (0,0.28) -- (0,2.3);
\draw[popmuted] (0,1.3) arc (90:125:1.3) node[midway,above left=-2pt] {\small $35^{\circ}$};
\end{tikzpicture}
\end{center}
\legende{Free-body diagram: weight, tension, spring force.}
\end{popfigure}
\item \textbf{Equations of equilibrium} \hfill (4 marks: 2 per equation)

Resolving horizontally ($F$ to the right balancing the horizontal component of $T$):
\[
F=T\sin 35^{\circ}.
\]
Resolving vertically:
\[
T\cos 35^{\circ}=W=mg=49.05\ \mathrm{N}.
\]
\item \textbf{Tension} \hfill (3 marks: method, substitution, answer)
\[
T=\frac{49.05}{\cos 35^{\circ}}=\frac{49.05}{0.8192}\approx 59.9\ \mathrm{N}.
\]
\item \textbf{Extension of the spring} \hfill (5 marks: value of $F$, statement of Hooke's law, assumption, calculation, answer in cm)

First, $F=T\sin 35^{\circ}=59.9\times 0.5736\approx 34.3\ \mathrm{N}$.

\textbf{Law used:} Hooke's law, $F=ke$; \textbf{assumption:} the elastic limit of the spring is not exceeded, so the extension is proportional to the applied force.
\[
e=\frac{F}{k}=\frac{34.3}{200}=0.172\ \mathrm{m}\approx 17.2\ \mathrm{cm}.
\]
\emph{Examiner's expectations:} a clear diagram with three forces only; correct resolution of $T$ (sine horizontal, cosine vertical, since the angle is measured from the vertical); explicit mention of Hooke's law and of the elastic-limit condition; final answer in centimetres with correct conversion.
\end{enumerate}
\end{corrige}

\begin{corrige}[Exercise 10 — GCE-style: experimental verification of Hooke's law]
\begin{enumerate}
\item \textbf{Procedure and precautions} \hfill (4 marks)

Suspend the spring vertically from a firm retort stand beside a vertical metre rule. Attach a pointer (e.g. a horizontal pin) to the lower end of the spring. Record the pointer reading with no load, then add slotted masses one at a time, recording the pointer reading for each load; the extension is the difference from the unloaded reading. \emph{Precautions:} read the scale at eye level to avoid parallax error; wait until the spring is at rest before reading; avoid overloading the spring beyond its elastic limit; check that the spring returns to its original length when the load is removed; use small, regular load increments.
\item \textbf{Graph} \hfill (4 marks: axes labelled with units, suitable scale, correct plotting, best straight line)

Plot $e$ (cm) on the vertical axis against $W$ (N) on the horizontal axis, e.g. $2\ \mathrm{cm}:1\ \mathrm{N}$ horizontally and $1\ \mathrm{cm}:1\ \mathrm{cm}$ of extension vertically. The points $(1,2.1)$, $(2,4.0)$, $(3,6.2)$, $(4,8.1)$, $(5,10.3)$ lie almost on a straight line through the origin; draw the best-fit straight line.
\begin{popfigure}
\begin{center}
\begin{tikzpicture}
\begin{axis}[width=0.7\linewidth,height=6cm,
xlabel={$W$ (N)},ylabel={$e$ (cm)},
xmin=0,xmax=6,ymin=0,ymax=12,
xtick={0,1,2,3,4,5,6},ytick={0,2,4,6,8,10,12},
grid=both,grid style={popline,very thin}]
\addplot[only marks,mark=*,popblue] coordinates {(1,2.1) (2,4.0) (3,6.2) (4,8.1) (5,10.3)};
\addplot[poporange,thick,domain=0:5.6] {2.05*x};
\end{axis}
\end{tikzpicture}
\end{center}
\legende{Extension against load: points and best-fit line through the origin.}
\end{popfigure}
\item \textbf{Verification and stiffness} \hfill (4 marks)

The graph is a \emph{straight line through the origin}, so $e\propto W$: the extension is proportional to the load, which is exactly Hooke's law. The stiffness is the inverse of the gradient (with $e$ converted to metres):
\[
\text{gradient}=\frac{\Delta e}{\Delta W}\approx\frac{10.0\ \mathrm{cm}}{5.0\ \mathrm{N}}=2.0\ \mathrm{cm\,N^{-1}}=0.020\ \mathrm{m\,N^{-1}},
\qquad k=\frac{1}{0.020}\approx 50\ \mathrm{N\,m^{-1}}.
\]
\item \textbf{Anomalous point at $W=5\ \mathrm{N}$} \hfill (3 marks)

Possible reasons: (i) a reading error, e.g. parallax or reading the pointer while the spring was still oscillating; (ii) the load was not exactly $5\ \mathrm{N}$ (worn or mislabelled mass), or the spring was beginning to be stretched beyond its elastic limit. Before concluding that the elastic limit has been reached, the student should repeat the reading, then \emph{remove the load and check whether the spring returns to its original unstretched length}: if it does, the elastic limit has not been exceeded and the point is simply an experimental error.
\end{enumerate}
\end{corrige}

\begin{corrige}[Exercise 11 — GCE-style: centre of gravity and stability]
\begin{enumerate}
\item \textbf{Locating the centre of gravity} \hfill (6 marks)

Make a small hole near the corner of the lamina and suspend it freely from a horizontal pin (or nail) so that it can swing without friction. Hang a plumb line from the same pin. When the lamina is at rest, draw the line of the plumb line on the lamina: the centre of gravity lies on this line. Repeat from a second hole at a different edge; the intersection of the two lines is the centre of gravity $G$ (a third suspension checks the result).

\emph{Explanation:} in equilibrium under gravity, the weight acts through $G$ and the reaction at the pin acts at the point of suspension. Two forces in equilibrium must be equal, opposite and collinear, so $G$ must lie vertically below the point of suspension; hence the line from the suspension point through $G$ is vertical. (Equivalently, the weight exerts a moment about the pin until $G$ comes to rest directly below it.)
\begin{popfigure}
\begin{center}
\begin{tikzpicture}[scale=0.9,>=latex]
\draw[fill=popblueL,draw=popink] (0,0) -- (3,0) -- (3,1.2) -- (2.2,1.2) -- (2.6,2.0) -- (3,2.0) -- (3,2.8) -- (0,2.8) -- cycle;
\fill[popink] (0,2.8) circle (1.5pt) node[above left] {\small pin};
\draw[popmuted,dashed] (0,2.8) -- (1.05,0.6);
\fill[poporange] (1.05,1.55) circle (2pt) node[right] {\small $G$};
\draw[->,popblue,thick] (1.05,1.55) -- (1.05,0.35) node[below] {\small $W$};
\draw[popmuted,dashed] (1.05,1.55) -- (1.05,2.9);
\end{tikzpicture}
\end{center}
\legende{Plumb-line method: $G$ lies vertically below the point of suspension.}
\end{popfigure}
\item \textbf{Stability of the wheelbarrow} \hfill (6 marks)

A rigid body resting on the ground is stable as long as the vertical line of action of its weight (the line through $G$) falls \emph{inside the base of support} (the region bounded by the contact points with the ground). If the load is placed too far in front of the wheel, $G$ moves forward; when the line of action of $W$ passes beyond the front edge of the base of support, the weight produces an unbalanced moment about that edge and the wheelbarrow topples forward. Moving the load back brings the line of action inside the base again and restores stability.
\begin{popfigure}
\begin{center}
\begin{tikzpicture}[scale=0.85,>=latex]
\draw[popink,thick] (-0.5,0) -- (4.5,0);
\draw[fill=popblueL,draw=popink] (0.4,0.25) -- (3.6,0.25) -- (3.2,1.1) -- (0.8,1.1) -- cycle;
\draw[popink,fill=white] (0.9,0.25) circle (0.25);
\fill[popgreen] (2.0,0.68) circle (2pt) node[above] {\small $G$};
\draw[->,popgreen,thick] (2.0,0.68) -- (2.0,-0.55);
\node[popgreen] at (2.0,-0.85) {\small stable};
\begin{scope}[xshift=6.2cm]
\draw[popink,thick] (-0.5,0) -- (4.5,0);
\draw[fill=popblueL,draw=popink] (0.4,0.25) -- (3.6,0.25) -- (3.2,1.1) -- (0.8,1.1) -- cycle;
\draw[popink,fill=white] (0.9,0.25) circle (0.25);
\fill[poporange] (3.9,0.68) circle (2pt) node[above] {\small $G$};
\draw[->,poporange,thick] (3.9,0.68) -- (3.9,-0.55);
\draw[->,poporange,thick] (3.6,0.25) arc (-70:-120:0.9);
\node[poporange] at (3.9,-0.85) {\small topples};
\end{scope}
\end{tikzpicture}
\end{center}
\legende{Left: line of action of $W$ inside the base (stable). Right: outside the base (toppling).}
\end{popfigure}
\item \textbf{Third force by resolution} \hfill (8 marks)

Take axes: $x$ due east, $y$ due north. Force 1: $\vec{F}_{1}=(10,0)$. Force 2: magnitude $15\ \mathrm{N}$ at $120^{\circ}$ anticlockwise from east:
\[
F_{2x}=15\cos 120^{\circ}=-7.5\ \mathrm{N},\qquad F_{2y}=15\sin 120^{\circ}=7.5\sqrt{3}\approx 12.99\ \mathrm{N}.
\]
Equilibrium requires $\vec{F}_{3}=-(\vec{F}_{1}+\vec{F}_{2})$:
\[
F_{3x}=-(10-7.5)=-2.5\ \mathrm{N},\qquad F_{3y}=-(0+12.99)=-12.99\ \mathrm{N}.
\]
Magnitude:
\[
F_{3}=\sqrt{2.5^{2}+12.99^{2}}=\sqrt{6.25+168.75}=\sqrt{175}\approx 13.2\ \mathrm{N}.
\]
Direction: $\tan\theta=\dfrac{12.99}{2.5}=5.196$, so $\theta\approx 79.1^{\circ}$; with both components negative, $\vec{F}_{3}$ points into the third quadrant (south-west), i.e. $79.1^{\circ}$ south of west. Equivalently, the direction is $259.1^{\circ}$ measured anticlockwise from east, or a bearing of about $190.9^{\circ}$ (almost due south, slightly west).

\emph{Graphical check (polygon of forces):} to a suitable scale (e.g. $1\ \mathrm{cm}:2\ \mathrm{N}$), draw the $10\ \mathrm{N}$ force eastwards, then from its head draw the $15\ \mathrm{N}$ force at $120^{\circ}$. For equilibrium the triangle must close: the closing side, drawn from the head of the second force back to the start, is the third force. Measuring gives about $6.6\ \mathrm{cm}$, i.e. $13.2\ \mathrm{N}$, at about $79^{\circ}$ south of west — confirming the calculation.

\emph{Examiner's expectations:} correct components with signs; use of $\sum F_{x}=0$, $\sum F_{y}=0$; magnitude from Pythagoras; direction stated unambiguously (angle with a named direction); clear description of the closing side of the force triangle.
\end{enumerate}
\end{corrige}

\newpage

\coursec{Summary sheet}

\begin{popfigure}
\begin{tikzpicture}[
  central/.style={circle, draw=popdark, fill=popblue, text=white, minimum size=2cm, font=\bfseries\small, align=center},
  branch/.style={rectangle, rounded corners, draw=popdark, fill=popblueL, text=popdark, minimum width=2.8cm, minimum height=1.4cm, font=\small, align=center},
  line/.style={draw=popdark, thick}
]
\node[central, align=center] (center){Equilibrium\\Forces\\Moments\\Couples};
\foreach \angle/\text/\col in {
  90/Equilibrium \& Centre of Gravity/popgreen,
  30/Force as Vector/poporange,
  -30/Types of Forces \& Free-Body Diagrams/poppurple,
  -90/Resultant \& Equilibrant/poppink,
  -150/Analytical Methods/popgold,
  150/Hooke's Law/popblueL,
  210/Applications \& Synthesis/popmuted%
}{
  \node[branch, fill=\col!20, draw=\col] (b\angle) at (\angle:4.5cm) {\text};
  \draw[line, \col] (center) -- (b\angle);
}
\end{tikzpicture}
\legende{Mind map of Chapter PMM14: Equilibrium Forces, Moments and Couples}
\end{popfigure}

\begin{bilan}
\textbf{Key Definitions}
\begin{itemize}
\item \cle{Equilibrium}: A body is in equilibrium if the resultant force and resultant moment acting on it are both zero.
\item \cle{Centre of Gravity (CG)}: The point where the whole weight of a body can be considered to act.
\item \cle{Force as a Vector}: A force has magnitude, direction, and point of application; represented by a vector.
\item \cle{Resultant Force}: The single force that has the same effect as a system of forces.
\item \cle{Equilibrant Force}: The force that brings a system into equilibrium; equal in magnitude, opposite in direction to the resultant.
\item \cle{Moment of a Force}: $M = F \times d$, where $d$ is the perpendicular distance from the pivot to the line of action.
\item \cle{Couple}: Two equal and opposite parallel forces separated by a distance; produces pure rotation.
\item \cle{Hooke's Law}: For an elastic material, extension $x$ is proportional to the applied force $F$: $F = kx$, where $k$ is the stiffness.
\end{itemize}

\textbf{Laws, Formulas \& Theorems}
\begin{itemize}
\item \textbf{Conditions for Equilibrium (Coplanar Forces)}:
  $\sum F_x = 0$, $\sum F_y = 0$, $\sum M_O = 0$ (about any point $O$).
\item \textbf{Triangle of Forces}: If three forces are in equilibrium, they can be represented in magnitude and direction by the sides of a triangle taken in order.
\item \textbf{Polygon of Forces}: For any number of coplanar forces in equilibrium, the force polygon closes.
\item \textbf{Resolution of Forces}: $F_x = F\cos\theta$, $F_y = F\sin\theta$.
\item \textbf{Resultant of Two Forces}: $R = \sqrt{F_1^2 + F_2^2 + 2F_1F_2\cos\theta}$.
\item \textbf{Moment of a Couple}: $M = F \times d$ (independent of the point about which moments are taken).
\item \textbf{Hooke's Law}: $F = kx$; Elastic potential energy $E = \frac{1}{2}kx^2$.
\item \textbf{Experimental Verification}: Plot $F$ vs $x$; gradient gives $k$; linear region confirms Hooke's law.
\end{itemize}

\textbf{Methods}
\begin{enumerate}
\item \textbf{Free-Body Diagram}: Isolate the body; show all external forces (weight, reactions, tensions, friction) with correct directions and points of application.
\item \textbf{Graphical Method}: Draw force vectors to scale; use triangle/polygon law; measure resultant/equilibrant.
\item \textbf{Analytical Method}: Resolve forces into perpendicular components; apply $\sum F_x=0$, $\sum F_y=0$; take moments about a convenient point to eliminate unknowns.
\item \textbf{Centre of Gravity Experiment}: Suspend a lamina from different points; draw vertical lines; intersection gives CG.
\item \textbf{Hooke's Law Experiment}: Hang known masses on a spring; measure extensions; plot $F$ vs $x$; determine $k$ from slope.
\end{enumerate}

\textbf{Common Mistakes}
\begin{itemize}
\item Forgetting to include the weight of the body or the reaction at a support in the free-body diagram.
\item Confusing resultant and equilibrant: equilibrant is opposite to resultant.
\item Taking moments about a point where an unknown force acts, but forgetting that the moment of that force is zero only if its line of action passes through the point.
\item Using the wrong perpendicular distance for moments (must be perpendicular to the line of action).
\item Assuming Hooke's law holds beyond the elastic limit; the graph becomes non-linear.
\item Mixing units (e.g., cm vs m) in moment calculations.
\item In graphical methods, not drawing to scale or not maintaining the correct order of vectors.
\end{itemize}

\textbf{GCE Tips}
\begin{itemize}
\item Always start with a clear free-body diagram; it earns marks and prevents missing forces.
\item For equilibrium problems, choose the point for moments that eliminates as many unknowns as possible (often a hinge or support).
\item Show all steps: resolve forces, write equilibrium equations, solve simultaneously.
\item In graphical questions, use a sharp pencil, ruler, and protractor; state the scale used.
\item For Hooke's law, remember to convert mass to force ($F=mg$) and extension to metres.
\item The centre of gravity of a uniform lamina is at its geometric centre; for composite bodies, use the principle of moments: $\sum m_i x_i = M \bar{x}$.
\item Practice past GCE questions on ladder problems, suspended bodies, and spring systems.
\end{itemize}
\end{bilan}

\findecours

\end{document}
